%\documentclass[dvipdfmx,a4paper,12pt]{amsart}
\documentclass[dvipdfmx,a4paper,11pt]{article} %titleとe-mailをコメントアウトする．


%%% Packages %%%
\setlength{\lineskip}{0pt}

% --- 和文フォント設定 (ゴシックを使うなら) ---amsartを使う時はコメントアウト
\renewcommand{\kanjifamilydefault}{\gtdefault}
\usepackage{otf}          % min10を避けるため
\usepackage{pxrubrica}    % 和文ルビ

% --- 基本パッケージ ---
\usepackage{graphicx}
\usepackage[all]{xy}
\usepackage{wrapfig}
\usepackage{pgfplots}
\usepackage{color}
\usepackage[dvipsnames]{xcolor}

% --- 数学関連 ---
\usepackage{amsmath,amssymb,amsthm,amsfonts,mathtools}
\usepackage{amscd,dsfont,bigdelim,braket,physics,mathrsfs,bm}

% --- 書式・リスト関連 ---
\usepackage{latexsym}
\usepackage{setspace}
\usepackage{multirow}
\usepackage{enumerate}
\usepackage{enumitem}

% --- コメント・取消線など ---
\usepackage{comment}
\usepackage[normalem]{ulem} % \emph の下線化を抑止（cancelと共存）

% --- URL・文字コード ---
\usepackage{url}
% \usepackage[utf8]{inputenc}  % ← 使用しているエンジンがuplatexなら不要、pdflatexなら有効に

% --- showkeys（常に表示） ---
%\usepackage{showkeys}
%\renewcommand*{\showkeyslabelformat}[1]{%
  %\fbox{\parbox{1.6cm}{\normalfont\tiny\sffamily#1\vspace{6mm}}}%
%}
% --- hyperrefは最後に読み込む ---
\usepackage[dvipdfmx,colorlinks,linkcolor=blue,anchorcolor=blue,citecolor=blue]{hyperref}

%%% レイアウト調整 %%%
%%% レイアウト調整（geometryに統一） %%%
\usepackage[
  top=30mm,        % 上余白
  bottom=30mm,     % 下余白
  left=25mm,       % 左余白
  right=25mm,      % 右余白
  headheight=12pt, % ヘッダー高さ（必要なら）
  headsep=10mm,    % ヘッダーと本文の間
  footskip=32pt,   % 本文とフッターの距離
  includehead,     % ヘッダー分も高さに含める
  includefoot      % フッター分も高さに含める
]{geometry}

%%% 行間調整（適宜 1.2 などに変更） %%%

\setstretch{1.1}

% --- 段落設定 --- 文字の行間ならここを変更する
\setlength{\parskip}{5pt}   % 段落間のスペース
\setlength{\parindent}{0pt}   % 段落先頭の字下げをなくす


%%% 追加（重複なし）パッケージ・設定 %%%

% --- 目次の体裁調整 ---
\usepackage{tocloft}
%\renewcommand{\contentsname}{目次} % 日本語化
\setlength{\cftbeforesecskip}{0pt}
\setlength{\cftbeforesubsecskip}{0pt}
\setlength{\cftbeforesubsubsecskip}{0pt}

% --- セクション見出しの体裁調整 ---
%\usepackage{titlesec}
%\titleformat*{\section}{\Large\bfseries}
%\titleformat*{\subsection}{\large\bfseries}
%\titlespacing*{\section}{0pt}{1.5ex plus .2ex minus .2ex}{0.8ex plus .1ex}
%\titlespacing*{\subsection}{0pt}{1.0ex plus .2ex minus .2ex}{0.5ex plus .1ex}

% --- ヘッダー／フッター設定 ---
%\usepackage{fancyhdr}
%\pagestyle{fancy}
%\fancyhf{}
%\rhead{岩井 雅崇}
%\lhead{大阪大学 数学専攻}
%\cfoot{\thepage}

% -- enumerate, itemize行間設定
% デフォルト設定
\setlist[itemize]{itemsep=3pt, parsep=0pt}
\setlist[enumerate]{itemsep=3pt, parsep=0pt}
% "変更する際は右を使う" \setlength{\parskip}{0cm} % 段落間\setlength{\itemsep}{5pt} % 項目間


\setlist[enumerate]{label=$(\arabic*)$}


% --- tcolorbox 設定 ---%\begin{tcolorbox}[mybox]と使う
\usepackage[most]{tcolorbox}
\tcbuselibrary{breakable, skins, theorems}
\tcbset{
  mybox/.style={
    colback = white,
    colframe = green!35!black,
    fonttitle = \bfseries,
    breakable = true
  }
}

%--緑枠自動化
\AtBeginEnvironment{prop}{\begin{tcolorbox}[mybox]}
\AtEndEnvironment{prop}{\end{tcolorbox}}
\AtBeginEnvironment{lem}{\begin{tcolorbox}[mybox]}
\AtEndEnvironment{lem}{\end{tcolorbox}}
\AtBeginEnvironment{thm}{\begin{tcolorbox}[mybox]}
\AtEndEnvironment{thm}{\end{tcolorbox}}
\AtBeginEnvironment{defn}{\begin{tcolorbox}[mybox]}
\AtEndEnvironment{defn}{\end{tcolorbox}}
\AtBeginEnvironment{cor}{\begin{tcolorbox}[mybox]}
\AtEndEnvironment{cor}{\end{tcolorbox}}
\AtBeginEnvironment{ques}{\begin{tcolorbox}[mybox]}
\AtEndEnvironment{ques}{\end{tcolorbox}}
\AtBeginEnvironment{conj}{\begin{tcolorbox}[mybox]}
\AtEndEnvironment{conj}{\end{tcolorbox}}
\AtBeginEnvironment{claim}{\begin{tcolorbox}[mybox]}
\AtEndEnvironment{claim}{\end{tcolorbox}}
\AtBeginEnvironment{thmdefn}{\begin{tcolorbox}[mybox]}
\AtEndEnvironment{thmdefn}{\end{tcolorbox}}


% --- TikZ 設定 ---
\usepackage{tikz}
\usetikzlibrary{positioning, arrows.meta, fit, calc, backgrounds}
\pgfdeclarelayer{background}
\pgfdeclarelayer{foreground}
\pgfsetlayers{background,main,foreground}
\usepackage{tikz-cd}

% --- footnote がページをまたがない設定 ---
\interfootnotelinepenalty=10000

% --- 目次に表示する階層の深さ ---
\setcounter{tocdepth}{2}

% --- 日本語目次---
\usepackage{pxjahyper}

%--newtheorem%--newcommand----


\newtheorem{thm}{Theorem}[section] 
\newtheorem{theo}[thm]{Theorem}
\newtheorem{theorem}[thm]{Theorem}
\newtheorem{cor}[thm]{Corollary}
\newtheorem{corollary}[thm]{Corollary}
\newtheorem{prop}[thm]{Proposition}
\newtheorem{proposition}[thm]{Proposition}
\newtheorem{conj}[thm]{Conjecture}
\newtheorem{conjecture}[thm]{Conjecture}
\newtheorem*{mainthm}{Theorem}
\newtheorem{deflem}[thm]{Definition-Lemma}
\newtheorem{lem}[thm]{Lemma}
\newtheorem{lemma}[thm]{Lemma}
\theoremstyle{definition} 
\newtheorem{defn}[thm]{Definition}
\newtheorem{definition}[thm]{Definition}
\newtheorem{propdefn}[thm]{Proposition-Definition} 
\newtheorem{lemdefn}[thm]{Lemma-Definition} 
\newtheorem{thmdefn}[thm]{Theorem-Definition} 
\newtheorem{eg}[thm]{Example} 
\newtheorem{ex}[thm]{Example} 
\newtheorem{example}[thm]{Example} 
\newtheorem{ques}[thm]{Question}
\newtheorem{question}[thm]{Question}
\newtheorem{remin}[thm]{Reminder}
\newtheorem{suma}[thm]{Summary}
\theoremstyle{remark}
\newtheorem{rem}[thm]{Remark}
\newtheorem{remark}[thm]{Remark}
\newtheorem{setup}[thm]{Setup}
\newtheorem{obs}[thm]{Observation}
\newtheorem{notation}[thm]{Notation}
\newtheorem{cl}{Claim}
\newtheorem{claim}[thm]{Claim}
\newtheorem{assup}[thm]{Assumption}
\newtheorem{step}{Step}
\newtheorem*{clproof}{Proof of Claim}
\newtheorem{cln}[thm]{Claim}
\newtheorem*{ack}{Acknowledgements} 
\numberwithin{equation}{section}
\newtheorem{case}{Case}


\newcommand{\rk}[0]{\operatorname{rk}}
\newcommand{\supp}[0]{\operatorname{Supp}}
\newcommand{\Rad}[0]{\operatorname{Rad}}
\newcommand{\Sha}[0]{\operatorname{Sha}}
\newcommand{\sha}[0]{\operatorname{sha}}
\newcommand{\eend}[0]{\operatorname{End}}
\newcommand{\codim}[0]{\operatorname{codim}}
\newcommand{\nd}[0]{\operatorname{nd}}
\renewcommand{\rank}[0]{\operatorname{rank}}
\newcommand{\degree}[0]{\operatorname{deg}}
\newcommand{\Exc}[0]{\operatorname{Exc}}
\newcommand{\pr}{\operatorname{pr}}
\newcommand{\id}{\operatorname{id}}
\newcommand{\Sym}{\operatorname{Sym}}
\newcommand{\End}[0]{\operatorname{End}}
\newcommand{\Coker}[0]{\operatorname{Coker}}


\newcommand{\Supp}{\operatorname{Supp}}
\newcommand{\Hom}[0]{\mathscr{H}\!\textit{om}}
\newcommand{\Ext}[0]{\mathscr{E}\!\textit{xt}}
\newcommand{\GL}[0]{\operatorname{GL}}
\newcommand{\SheafHom[1]}{\mathscr{H}\!\!\!\text{\calligra om}_{\,{#1}}}
\newcommand{\PGL}[0]{\mathbb{P}\GL(r,\C)}

\newcommand{\Exp}{{\rm Exp}}
\newcommand{\Burn}{{\rm Burn}}
\newcommand{\Am}{{\rm Am}}
\newcommand{\Cr}{{\rm Cr}}
\newcommand{\Kum}{{\rm Kum}}
\newcommand{\BKR}{{\rm BKR}}
\newcommand{\Ann}{{\rm Ann}}
\newcommand{\Char}{{\rm Char}}
\newcommand{\Null}{{\rm Null}}
\newcommand{\rd}{{\rm rd}}
\newcommand{\fint}{{\rm fint}}
\newcommand{\Nlc}{{\rm Nlc}}
\newcommand{\NKLT}{{\rm NKLT}}
\newcommand{\NE}{{\rm NE}}

\newcommand{\Pic}{\operatorname{Pic}}
\newcommand{\Gr}{{\rm Gr}}
\newcommand{\NS}{{\rm NS}}
\newcommand{\Nef}{{\rm Nef}}
\newcommand{\Amp}{{\rm Amp}}
\newcommand{\Pos}{{\rm Pos}}
\newcommand{\Id}{{\rm Id}}
\newcommand{\Hilb}{{\rm Hilb}}
\newcommand{\VSP}{{\rm VSP}}
\newcommand{\Mov}{{\rm Mov}}
\newcommand{\Proj}{{\rm Proj}}
\newcommand{\Int}{{\rm Int}}
\newcommand{\Mon}{{\rm Mon}}
\newcommand{\disc}{{\rm disc}}
\newcommand{\Bir}{{\rm Bir}}
\newcommand{\Chow}{{\rm Chow}}
\newcommand{\RatCurves}{{\rm RatCurves}}
\newcommand{\Sing}{{\rm Sing}}
\newcommand{\mult}{{\rm mult}}
\newcommand{\length}{{\rm length}}
\newcommand{\PP}{\mathbb{P}}
\newcommand{\loc}{{\rm loc}}
\newcommand{\Grass}{{\rm Grass}}
\newcommand{\Deck}{{\rm Deck}}
\newcommand{\Spec}{\operatorname{Spec}}
\newcommand{\Var}{{\rm Var}}



\newcommand{\Fix}{{\rm Fix}}
\newcommand{\rang}{{\rm rang}}
\newcommand{\cJ}{\mathcal{J}}
\newcommand{\J}{\mathcal{J}}
\newcommand{\rg}{{\rm rg}}
\newcommand{\Ens}{{\rm Ens}}
\newcommand{\LL}{{\rm LL}}
\newcommand{\pgcd}{{\rm pgcd}}
\newcommand{\utype}{{\rm utype}}
\newcommand{\Mat}{{\rm Mat}}
\newcommand{\Span}{{\rm Span}}
\newcommand{\Curv}{{\rm Curv}}
\newcommand{\Fl}{{\rm Fl}}
\newcommand{\ass}{{\rm ass}}
\newcommand{\dra}{{\rm dra}}
\newcommand{\fk}{{\rm fk}}
\newcommand{\dom}{{\rm dom}}
\newcommand{\tdim}{{\rm tdim}}
\newcommand{\Type}{{\rm Type}}
\newcommand{\yng}{{\rm yng}}
\newcommand{\Rm}{{\rm Rm}}
\newcommand{\Rc}{{\rm Rc}}
\newcommand{\discrep}{{\rm discrep}}
\newcommand{\kst}{{\rm kst}}
\newcommand{\cord}{{\rm cord}}
\newcommand{\ind}{{\rm ind}}
\newcommand{\Area}{{\rm Area}}
\newcommand{\Nullspace}{{\rm Nullspace}}
\newcommand{\Graph}{{\rm Graph}}
\newcommand{\Volume}{{\rm Volume}}
\newcommand{\Card}{{\rm Card}}
\newcommand{\Bimer}{{\rm Bimer}}
\newcommand{\genus}{{\rm genus}}
\newcommand{\Mass}{{\rm Mass}}
\newcommand{\Cal}{{\rm Cal}}
\newcommand{\Vect}{{\rm Vect}}
\newcommand{\e}{e}
\newcommand{\inj}{{\rm inj}}
\newcommand{\Irr}{{\rm Irr}}
\newcommand{\len}{{\rm len}}

\newcommand{\Dol}{{\rm Dol}}
\newcommand{\Hol}{{\rm Hol}}
\newcommand{\Isom}{{\rm Isom}}
\newcommand{\Spin}{{\rm Spin}}
\newcommand{\Teich}{{\rm Teich}}
\newcommand{\Def}{{\rm Def}}
\newcommand{\vol}{{\rm vol}}
\newcommand{\slope}{{\rm slope}}
\newcommand{\ch}{{\rm ch}}
\newcommand{\td}{{\rm td}}
\newcommand{\Flag}{{\rm Flag}}
\newcommand{\Quot}{{\rm Quot}}
\newcommand{\coker}{{\rm coker}}
\newcommand{\Td}{{\rm Td}}
\newcommand{\Num}{{\rm Num}}
\newcommand{\cF}{\mathcal{F}}
\newcommand{\cO}{\mathcal{O}}
\newcommand{\trdeg}{{\rm trdeg}}
\newcommand{\red}{{\rm red}}
\newcommand{\Max}{{\rm Max}}
\newcommand{\an}{{\rm an}}
\newcommand{\Sup}{{\rm Sup}}
\newcommand{\Herm}{{\rm Herm}}


\newcommand{\ord}{{\rm ord}}
\newcommand{\Fut}{{\rm Fut}}
\newcommand{\gr}{{\rm gr}}
\newcommand{\lct}{{\rm lct}}
\newcommand{\cX}{\mathcal{X}}
\newcommand{\Val}{{\rm Val}}
\newcommand{\fa}{\mathfrak{a}}
\newcommand{\Ex}{{\rm Ex}}
\newcommand{\A}{\mathbb{A}}
\newcommand{\cY}{\mathcal{Y}}
\newcommand{\cD}{\mathcal{D}}
\newcommand{\Cent}{{\rm Cent}}
\newcommand{\Gm}{{\rm Gm}}
\newcommand{\Rees}{{\rm Rees}}
\newcommand{\Ding}{{\rm Ding}}
\newcommand{\Sch}{{\rm Sch}}
\newcommand{\Univ}{{\rm Univ}}
\newcommand{\im}{{\rm im}}
\newcommand{\Diff}{{\rm Diff}}
\newcommand{\sm}{{\rm sm}}
\newcommand{\NA}{{\rm NA}}
\newcommand{\cH}{\mathcal{H}}
\newcommand{\Iso}{{\rm Iso}}
\newcommand{\Ad}{{\rm Ad}}
\newcommand{\RHom}{R\mathscr{H}\!\textit{om}}
\newcommand{\hol}{{\rm hol}}
\newcommand{\Sh}{{\rm Sh}}
\newcommand{\lt}{{\rm lt}}
\newcommand{\cA}{\mathcal{A}}
\newcommand{\cM}{\mathcal{M}}
\newcommand{\cZ}{\mathcal{Z}}
\newcommand{\cU}{\mathcal{U}}
\newcommand{\Ivol}{{\rm Ivol}}
\newcommand{\coeff}{{\rm coeff}}
\newcommand{\Sets}{{\rm Sets}}
\newcommand{\PSH}{{\rm PSH}}
\newcommand{\Eff}{{\rm Eff}}
\newcommand{\Mor}{{\rm Mor}}
\newcommand{\Ncf}{{\rm Ncf}}
\newcommand{\Sec}{{\rm Sec}}
\newcommand{\et}{{\rm et}}
\newcommand{\cont}{{\rm cont}}
\newcommand{\curv}{{\rm curv}}
\newcommand{\dist}{\operatorname{dist}}
\newcommand{\diam}{\operatorname{diam}}
\newcommand{\Reg}{{\rm Reg}}


\newcommand{\Mab}{{\rm Mab}}
\newcommand{\KE}{{\rm KE}}
\newcommand{\MA}{{\rm MA}}
\newcommand{\Ent}{{\rm Ent}}
\newcommand{\triv}{{\rm triv}}
\newcommand{\EPDiv}{{\rm EPDiv}}
\newcommand{\OO}{\mathcal{O}}
\newcommand{\depth}{\rm{depth}}
\newcommand{\gon}{\rm{gon}}
\newcommand{\ddc}{dd^c}
\newcommand{\VCar}{{\rm VCar}}
\newcommand{\CPSH}{{\rm CPSH}}
\newcommand{\na}{{\rm na}}
\newcommand{\rad}{{\rm rad}}
\newcommand{\PL}{{\rm PL}}
\newcommand{\FS}{{\rm FS}}
\newcommand{\CL}{{\rm CL}}
\newcommand{\CX}{{\rm CX}}
\newcommand{\CH}{{\rm CH}}
\newcommand{\ext}{{\rm ext}}
\newcommand{\cL}{\mathcal{L}}
\newcommand{\qm}{{\rm qm}}
\newcommand{\CW}{{\rm CW}}
\newcommand{\fm}{{\rm fm}}
\newcommand{\Ch}{{\rm Ch}}
\newcommand{\Sign}{{\rm Sign}}
\newcommand{\Ar}{{\rm Ar}}
\newcommand{\sh}{{\rm sh}}
\newcommand{\exph}{{\rm exph}}
\newcommand{\Ca}{{\rm Ca}}
\newcommand{\nef}{{\rm nef}}
\newcommand{\psef}{{\rm psef}}
\newcommand{\Kahl}{{\rm Kahl}}
\newcommand{\kod}{{\rm kod}}
\newcommand{\exc}{{\rm exc}}
\newcommand{\CP}{\mathbb{P}}
\newcommand{\HSC}{{\rm HSC}}
\newcommand{\Fr}{{\rm Fr}}
\newcommand{\Jac}{{\rm Jac}}
\newcommand{\Conn}{{\rm Conn}}
\newcommand{\Hod}{{\rm Hod}}
\newcommand{\Nm}{{\rm Nm}}
\newcommand{\Bun}{{\rm Bun}}
\newcommand{\DMod}{{\rm DMod}}
\newcommand{\Fuk}{{\rm Fuk}}
\newcommand{\alb}{{\rm alb}}
\newcommand{\NNef}{{\rm NNef}}
\newcommand{\Diam}{{\rm Diam}}
\newcommand{\RCD}{{\rm RCD}}
\newcommand{\Ram}{{\rm Ram}}
\newcommand{\osc}{{\rm osc}}
\newcommand{\Hess}{{\rm Hess}}
\newcommand{\T}{\mathbb{T}}
\newcommand{\F}{\mathcal{F}}
\newcommand{\Todd}{{\rm Todd}}
\newcommand{\YM}{{\rm YM}}
\newcommand{\BC}{{\rm BC}}
\newcommand{\cC}{\mathcal{C}}
\newcommand{\glct}{{\rm glct}}
\newcommand{\Cone}{{\rm Cone}}
\newcommand{\Tsch}{{\rm Tsch}}
\newcommand{\wt}{{\rm wt}}
\newcommand{\val}{{\rm val}}
\newcommand{\rdim}{{\rm rdim}}
\newcommand{\CDiv}{{\rm CDiv}}
\newcommand{\Br}{{\rm Br}}
\newcommand{\Manin}{{\rm Manin}}
\newcommand{\Map}{{\rm Map}}
\newcommand{\nKLT}{{\rm nKLT}}
\newcommand{\Cl}{{\rm Cl}}
\newcommand{\Ass}{{\rm Ass}}
\newcommand{\Dom}{{\rm Dom}}
\newcommand{\Bs}{{\rm Bs}}

\newcommand{\ank}{{\rm ank}}
\newcommand{\sgn}{{\rm sgn}}
\newcommand{\orb}{{\rm orb}}
\newcommand{\Diag}{{\rm Diag}}
\newcommand{\refl}{{\rm refl}}
\newcommand{\Ind}{{\rm Ind}}
\newcommand{\Sp}{{\rm Sp}}
\newcommand{\U}{{\rm U}}
\newcommand{\SL}{{\rm SL}}
\newcommand{\SU}{{\rm SU}}
\newcommand{\PSL}{{\rm PSL}}
\newcommand{\g}{\mathfrak{g}}
\newcommand{\E}{\mathcal{E}}
\newcommand{\cE}{\mathcal{E}}
\newcommand{\cQ}{\mathcal{Q}}
\newcommand{\Oo}{\mathcal{O}}
\newcommand{\ft}{{\rm ft}}
\newcommand{\fsu}{{\rm fsu}}
\newcommand{\fso}{{\rm fso}}
\newcommand{\fe}{{\rm fe}}
\newcommand{\Out}{{\rm Out}}
\newcommand{\Inn}{{\rm Inn}}
\newcommand{\Lie}{{\rm Lie}}
\newcommand{\Coh}{{\rm Coh}}
\newcommand{\Gal}{{\rm Gal}}
\newcommand{\Stab}{{\rm Stab}}
\newcommand{\diag}{{\rm diag}}
\newcommand{\aut}{{\rm aut}}
\newcommand{\Scal}{{\rm Scal}}
\newcommand{\DF}{{\rm DF}}
\newcommand{\V}{{\rm  V}}
\newcommand{\X}{{\rm  X}}
\newcommand{\G}{{\rm  G}}
\newcommand{\Maps}{{\rm  Maps}}
\newcommand{\Bl}{{\rm  Bl}}
\newcommand{\Higgs}{{\rm  Higgs}}
\newcommand{\Rep}{{\rm  Rep}}
\newcommand{\pHiggs}{{\rm pHiggs}}
\newcommand{\sRep}{{\rm sRep}}
\newcommand{\image}{{\rm image}}
\newcommand{\tors}{{\rm tors}}
\newcommand{\sat}{{\rm sat}}
\newcommand{\torsion}{{\rm torsion}}
\newcommand{\codh}{{\rm codh}}
\newcommand{\Kod}{{\rm Kod}}
\newcommand{\Ban}{{\rm Ban}}
\newcommand{\Mult}{{\rm Mult}}

\newcommand{\Alb}{\operatorname{Alb}}
\newcommand{\verti}{{\rm vert}}
\newcommand{\hor}{{\rm hor}}
\newcommand{\univ}{{\rm univ}}
\newcommand{\Tor}{{\rm tor}}
\newcommand{\shaf}{\mathrm{sha}}
\newcommand{\Shaf}{\mathrm{Sha}}
\newcommand{\reg}{{\rm{reg}}}
\newcommand{\sing}{{\rm{sing}}}
\newcommand{\qt}{{\rm{qt}}}
\newcommand\sO{{\mathcal O}}
\newcommand{\Div}[0]{\operatorname{div}}
\newcommand{\ddbar}{dd^c}
\newcommand{\cV}{\mathcal{V}}
\newcommand{\deldel}{\sqrt{-1}\partial \bar\partial}
\newcommand{\dbar}{\bar\partial}
\newcommand{\I}[1]{\mathcal{I}(#1)}
\newcommand{\Aut}[1]{\operatorname{Aut}(#1)}
\newcommand{\Ker}[1]{\ker(#1)}
\newcommand{\Image}[1]{\operatorname{Im}(#1)}
\DeclareMathOperator{\Ric}{Ric}
\DeclareMathOperator{\Vol}{Vol}
 \newcommand{\pdrv}[2]{\frac{\partial #1}{\partial #2}}
 \newcommand{\drv}[2]{\frac{d #1}{d#2}}
  \newcommand{\ppdrv}[3]{\frac{\partial #1}{\partial #2 \partial #3}}
%\newcommand{\underalign}[2]{\quad \underset{\mathclap{\strut #1}}{#2}\quad}
\newcommand{\underalign}[2]{%
  \quad
  \underset{\mathclap{\raisebox{-1.6ex}{\ensuremath{\scriptstyle #1}}}}{#2}
  \quad
}
\newcommand{\polar}{\beta}
  
\newcommand{\R}{\mathbb{R}}
\newcommand{\Z}{\mathbb{Z}}
\newcommand{\N}{\mathbb{N}}
\newcommand{\C}{\mathbb{C}}
\newcommand{\Q}{\mathbb{Q}}
\newcommand{\D}{\mathbb{D}}
\newcommand{\mP}{\mathbb{P}}
\newcommand{\mO}{\mathcal{O}}
\newcommand{\B}{\mathds{B}}
%\newcommand{\tl}{\hspace{-0.8ex}<\hspace{-0.8ex}}
%\renewcommand{\tr}{\hspace{-0.8ex}>}

\newcommand{\xb}[1]{\textcolor{blue}{#1}}
\newcommand{\xr}[1]{\textcolor{red}{#1}}
\newcommand{\xm}[1]{\textcolor{magenta}{#1}}


\title{の和訳}
%\author{Masataka IWAI}
%\address{Department of Mathematics, Graduate School of Science, Osaka City University 3-3-138, Sugimoto, Sumiyoshi-ku Osaka, 558-8585Japan} 
%\email{{\tt masataka.math@gmail.com}}
%\email{{\tt masataka.math@gmail.com, masataka@sci.osaka-cu.ac.jp}}



\date{\today, version 0.01}


\renewcommand{\thefootnote}{\arabic{footnote}}

\baselineskip = 15pt
\footskip = 32pt

\begin{document}

\title{Yauの定理とその応用}
\author{Masataka IWAI}

\maketitle
%
\begin{abstract}
Yauの定理とその応用を勉強した時のノートです. 
\end{abstract}
%
\setcounter{tocdepth}{3}
\tableofcontents

\section*{初めに}

2020年に基本群に関して勉強していた内容を今になって復習したくなったため作りました. 
なおBishop-Gromovの定理などは引用に留めております. 
せっかくなのでChern類や榎(Enoki)先生のgeneric nefness theorem\cite{Eno88}, Junyan Caoさんの論文\cite{Cao13}なども盛り込みました. 

\begin{itemize}
\item 2章と3章は\cite[Section 2, 3]{Sze14}の内容. \cite{Sze14}のこの部分を改めて見ると, Monge--Amp\`ere方程式の証明がかなり簡潔にまとまっていて良い本だと思った. 著作権的にまずい場合はこの部分は消します. (ただ加筆が多いので元々の文章とだいぶ異なってる.)
\item 4章と5章は半分はChern類の内容. \cite[Section 4--6]{Iwa-Chern}の内容とかぶっている. 研究で榎(Enoki)先生のgeneric nefness theorem \cite{Eno88}を使っていたが, この際だしこの証明をつけることにした. 
\item 6章は榎(Enoki)先生のstability theorem \cite{Eno88}の内容. \cite{Eno88}は改めて見ても時代を先取りしている論文である. 面白い. 
\item 7章の基本群の話は, \cite{DPS96, Paun98}に基づく. 一部Junyan Caoさんの論文\cite{Cao13}の内容も入れた. これはずっと気になっていたが証明を読んでいなかったので, この際だしこの証明をつけることにした.
\end{itemize}

ちなみにこのノートを書くきっかけは, 2026年1月の集会でJuanyong Wangさんが\cite{FGSW26}の講演をしていたから. 7章の内容は何かしら今後進展があると思う. 

誤字脱字や間違いはあると思うので注意して読んでください. またかなりの部分をchatGPTを用いて加筆しております. 

\newpage

\section{微分幾何学の復習}
アインシュタイン規約, 複素多様体, Chern曲率などに関しては\cite{Iwa-Chern}参照.
ただし\ref{sec-Yau-App1}章まではChern曲率はでない. 

以後アインシュタイン規約を用いる. 
$\omega$ K\"ahler form, $\alpha$ $(1,1)$-formについて$\sharp$-operatorを
\[
\sharp^{\omega} \alpha  \colon T_X \to T_X \quad
\alpha(u, \bar{v})=\omega(\sharp^{\omega} \alpha(u), \bar{v})
\]
として定める. 
$\omega=\sqrt{-1}g_{i \bar{j}} dz^i \wedge d \bar{z}^{\bar{j}}$かつ
$\alpha = \sqrt{-1}\alpha_{i \bar{j}} dz^i \wedge d \bar{z}^{\bar{j}}$ならば
\[
\sharp^{\omega} \alpha
\coloneqq  g^{i \bar{k}}\alpha_{j \bar{k}}
dz^{j} \otimes \frac{\partial}{\partial z^i}
% = \operatorname{tr}_{\omega}(\alpha)
\]
である. 
また$\operatorname{tr}_{\omega}(\alpha)$を
\[
\operatorname{tr}_{\omega}(\alpha)\coloneqq  g^{i\bar{j}}\alpha_{i\bar{j}}
\]
として定める. 次の公式がある. 
\begin{tcolorbox}[mybox]
\[
n \alpha \wedge \omega^{n-1}
= \operatorname{tr}_{\omega}(\alpha) \omega^{n}
\]
\end{tcolorbox}
$\operatorname{tr}_{\omega}$は調和積分論での$\Lambda_{\omega}$と同じである. 
また定義から以下が成り立つ. 
\[
\operatorname{tr}(\sharp^{\omega} \alpha)
=\sum_{i}g^{i \bar{k}}\alpha_{i\bar{k}}
=\operatorname{tr}_{\omega}(\alpha)
\]
K\"ahlerの公式をまとめると次のとおり. 以下$\partial_{i}\coloneqq \frac{\partial}{\partial z^{i}}, \partial_{\bar j }\coloneqq \frac{\partial}{\partial \bar{z}^{\bar{j}}}$と略記する. 
\begin{tcolorbox}[mybox]
\begin{suma}
\label{suma-brief-diff}
$\omega=\sqrt{-1}g_{i \bar{j}} dz^i \wedge d \bar{z}^{\bar{j}}$をK\"ahler計量とする時
\begin{itemize}
\item Christoffel記号. $\Gamma_{jk}^{i} = g^{i\bar{l}} \partial_j g_{k\bar{l}} $. $\Gamma_{jk}^{i} = \Gamma_{kj}^{i}$
\item $\omega$によって$T_X$にHermite計量$h$を入れる時, Chern接続 
\[
D_h \colon A^0(T_X) \to A^{1}(T_X)
\]
が定まる. これは$e_i\coloneqq \frac{\partial}{\partial z^i}$, $s^i$を$C^\infty$関数とすると, connection form $\theta$は$(1,0)$-formで$\theta=h^{-1} \partial h$で与えられる. これはつまり
\[
D_h (s^i e_i) = (\partial + \bar\partial)s^i \otimes e_i + s^j \theta^{i}_{j} \otimes e_i
\quad
\theta^{i}_{j} = h^{i \bar{l}}\partial_j h_{k \bar{l}} dz^k =\Gamma_{jk}^{i} dz^k
\]
となる. 
\item curvature tensorを
\[
R^{j}_{i k \bar{l}}\coloneqq  - \partial_{\bar{l}}\Gamma^{j}_{ik}
\]
と定める. この時Chern曲率$F_h$は$\End(T_X)$値$(1,1)$-formなので
\[
F_h \coloneqq R^{j}_{i k \bar{l}}
\underbrace{\frac{\partial}{\partial z^{j}}\otimes dz^i}_{\text{$\End(T_X)$値}} 
\otimes \underbrace{dz^k \wedge\,d\bar{z}^{\bar{l}}}_{\text{$(1,1)$-form}}
\]
%($\frac{1}{2 \pi}$をつけているのは, $dd^c = \frac{i}{2\pi}\partial \bar{\partial}$とDGPでは採用しているので, その都合上)
\item curvature tensorを
\[
R_{i \bar{j}k \bar{l}}\coloneqq  g_{p\bar{j}} R^{p}_{i k \bar{l}}
\]とすると次が成り立つ
\begin{equation*}
R_{i \bar{j}k \bar{l}} = -\partial_i \partial_{\bar{j}} g_{k \bar{l}}
+g^{p\bar{q}}(\partial_i g_{k \bar{q}})(\partial_{\bar{j}}g_{p\bar{l}})
\end{equation*}
また$i,k$や$\bar{j}, \bar{l}$は入れ替えても良い.またGriffith positivityやNakano positivityはこれでわかる.
\item Ricci curvature tensorを
\[
R_{i\bar{j}}\coloneqq g^{k \bar{l}}R_{i \bar{j}k \bar{l}} 
\]
とすると次が成り立つ
\[
R_{i \bar{j}} = -\partial_i \partial_{\bar{j}} \log \det(g)
\]
Ricci曲率は以下のように定める. 
\[
\Ric(\omega)\coloneqq  \frac{\sqrt{-1}}{2 \pi} R_{i\bar{j}}dz^i \wedge\,d\bar{z}^{\bar{j}}
\in c_1(X) \in H^{2}(X, \R)
\]
これによって実$(1,1)$-formになり, $H^{2}(X, \Z)$に入る. 
\item Scalar curvature $R(\omega)$を以下のように定める. 
\[
R(\omega)\coloneqq  \operatorname{tr}_{\omega}\Ric(\omega)=
\frac{1}{2 \pi} g^{i\bar{j}}R_{i\bar{j}}
\]
\end{itemize}
\end{suma}
\end{tcolorbox}
Ricci曲率に関しては本によって$2\pi$のズレが出る. 
\cite{Sze14}は$\Ric(\omega) \in 2 \pi c_1(X)$を採用している

%$\omega$をKahler formとするとき, $T_X$にHermite 計量$h$が入る. 
%そのChern curvature tensorを$\Theta_{T_X, h}$と表す. 
%するとそれはEnd $T_X$値 $(1,1)$-formで次のようにかける. 
%$$\Theta_{T_X, h}= \frac{1}{2 \pi} R^{j}_{i k \bar{l}} \partial_z^{i} \otimes \partial_{j} \otimes \partial_z^{k} \wedge \partial_{\bar{l}} $$
%ここで便宜上$\frac{1}{2\pi}$を入れる. これは$$dd^c = \frac{\sqrt{-1}}{2 \pi}\partial \overline{\partial}$$を以後使っていくためである. 

\begin{lem}[{\cite{Gue16}}]
\label{lem-sharp-formula}
$\End(T_X)$値$(n, n)$-formの式として次が成り立つ.
\[
n\cdot \frac{\sqrt{-1}}{2 \pi}F_{h}\wedge \omega^{n-1}
= \sharp^{\omega}\Ric(\omega) \cdot \omega^{n} 
\]
\end{lem}
\begin{proof}
\begin{align*}
n\cdot \frac{\sqrt{-1}}{2 \pi}F_{h}\wedge \omega^{n-1}
&\underalign{\text{(公式)}}{=}  
 \frac{1}{2 \pi}\operatorname{tr}_{\omega}(\sqrt{-1}F_{h})\omega^{n} 
 \underalign{\text{(curvature tensorの定義)}}{=}  
\frac{1}{2 \pi} g^{k \bar{l}}R^{j}_{i k \bar{l}} \frac{\partial}{\partial z^{j}}\otimes dz^i\cdot \omega^{n}\\
&\underalign{\text{(Riemann curvatureの定義)}}{=}  
\frac{1}{2 \pi} g^{j \bar{p}}g^{k \bar{l}}R_{i \bar{p} k \bar{l}} \frac{\partial}{\partial z^{j}}\otimes dz^i\cdot \omega^{n}
\underalign{\text{(Ricciの定義)}}{=} 
\frac{1}{2 \pi} g^{j \bar{p}}R_{i \bar{p}} \frac{\partial}{\partial z^{j}}\otimes dz^i \cdot \omega^{n}\\
&\underalign{\text{($\sharp$とRicciの定義)}}{=} 
\sharp^{\omega}\Ric(\omega) \cdot \omega^{n} 
\end{align*}
\end{proof}



\newpage

\section{Schauder estimateと楕円型作用素}
Szekelyhidiのlecture note \cite[Chapter 2]{Sze14}による

\subsection{Harmonic function}
\begin{defn}
$U\subset \mathbb{R}^n$を開集合とする.$C^2$関数$f\colon U\to \mathbb{R}$がharmonicであるとは
\[
\Delta f\coloneqq \frac{\partial^2 f}{\partial x^1\partial x^1}+\cdots+\frac{\partial^2 f}{\partial x^n\partial x^n}=0
\quad\text{on }U
\]
が成り立つこと.
\end{defn}



\begin{thm}[平均値等式 {\cite[Theorem 2.1]{Sze14}}]
\label{thm-Sze-2.1}
$f\colon U\to\mathbb{R}$がharmonicで,$x\in U$かつ$B_r(x)\subset U$なら
\[
f(x)=\frac{1}{\Vol(\partial B_r)}\int_{\partial B_r(x)} f(y)\,dy.
\]
ここで任意の$x\in\mathbb{R}^n$, $r>0$について
\[
B_r(x)\coloneqq \{ y \in \mathbb{R}^n  \mid \lVert y-x\rVert<r \}
\] 
とし, $B_r\coloneqq B_r(0)$とする. 
\end{thm}

\begin{proof}
Greenの定理を思い出す.
\begin{thm}[{Greenの定理}]
\[
\int_U \left( \psi \Delta \varphi + \nabla \psi \cdot \nabla \varphi \right) \,dV
=
\int_{\partial U} \psi \left( \nabla \varphi \cdot \mathbf{n} \right) \,dS
=
\int_{\partial U} \psi \nabla \varphi \cdot\,d\mathbf{S}
\]
\end{thm}
さて$\rho\leq r$に対して
\[
F(\rho)\coloneqq \int_{\partial B_1} f(x+\rho y)\,dy
\]
と定める.このとき
\begin{align*}
\frac{\partial }{\partial \rho} F(\rho)
&\underalign{}{=}
\int_{\partial B_1} 
\sum_{i=1}^{n}
\partial_i f(x+\rho y)y_i
\,dy
\underalign{\text{(chain rule)}}{=}
\frac{1}{\rho}
\int_{\partial B_1} 
\underbrace{
\sum_{i=1}^{n}
\frac{\partial}{\partial y^i}f(x+\rho y)y_i
}_{=\nabla_y f(x+\rho y)\cdot y}
\,dy
\\
&\underalign{\text{(Greenの定理)}}{=}
\frac{1}{\rho}
\int_{B_1}
\sum_{i=1}^{n}
\frac{\partial^2}{\partial y^i\partial y^i}
f(x+\rho y)
\,dy
\underalign{\text{(chain rule)}}{=}
\rho
\int_{B_1}
\underbrace{
\sum_{i=1}^{n}
\partial_i\partial_i f(x+\rho y)
}_{=\Delta f(x+\rho y)}
\,dy
\underalign{\text{(harmonicの仮定)}}{=}
0.
\end{align*}
%%%%%%%%%%%%
\begin{comment}
\begin{align*}
\frac{\partial }{\partial \rho} F(\rho)
&\underalign{}{=}
\int_{\partial B_1} 
\underbrace{\sum_{i=1}^{n}\frac{\partial }{\partial y^i}f(x+\rho y) y_i}
_{=\nabla f(x+\rho y)\cdot y}
\,dy
\\
&\underalign{\text{(Greenの定理)}}{=}
\qquad \int_{B_1}
\underbrace{\sum_{i=1}^{n}\frac{\partial^2}{\partial y^i \partial y^i}f(x+\rho y) }
_{=\Delta f(x+\rho y)}
dy
\underalign{\text{(Harmonicの仮定)}}{=}
0.
\end{align*}
\end{comment}
%%%%%%%%%%%

よって$F$は定数である.また変数変換により
\[
F(r)=
\int_{\partial B_1} f(x+r y)\,dy
=
\frac{\Vol(\partial B_1)}{\Vol(\partial B_r)}\int_{\partial B_r(x)} f(y)\,dy.
\]
以上より
\[
F(\rho)
=\lim_{r\to 0}F(r)
=
\Vol(\partial B_1) 
\underbrace{\lim_{r\to 0}\frac{1}{\Vol(\partial B_r)}\int_{\partial B_r(x)} f(y)\,dy,}
_{=f(x) \text{  by 微積分の基本定理}}
=\Vol(\partial B_1)f(x)
\]
\end{proof}



\begin{cor}[{\cite[Corollary 2.2]{Sze14}}]
\label{cor-Sze-2.2}
$C^\infty$写像$\eta\colon \mathbb{R}^n\to\mathbb{R}$がradially symmetric, つまり
 $\lVert z\rVert = \lVert w\rVert$ならば$\eta(z) = \eta(w)$を満たすとする. 
さらに$\Supp(\eta) \subset B_1$かつ$\int_{ \mathbb{R}^n} \eta\,d\mu=1$とする. 
 $f\colon B_2\to\mathbb{R}$がharmonicなら,任意の$x\in B_1$に対して
\[
f(x)=\int_{\mathbb{R}^n} f(x-y)\eta(y)\,dy
=\int_{\mathbb{R}^n} f(y)\eta(x-y)\,dy
\]
\end{cor}
\begin{proof}
$\eta(y) = \phi(\lvert y\rvert)$となる$\phi \colon \R^n \to \R$をとると, 
\begin{align*}
&\int_{\mathbb{R}^n} f(x-y)\eta(y)\,dy
\underalign{\text{(置換積分)}}{=}\int_{0}^1
\int_{\Vol(\partial B_1)} f(x-r \theta)\phi(r) r^{n-1} d \theta\,dr
\\
&\underalign{}{=}
\int_{0}^1
\underbrace{ \left(\int_{\Vol(\partial B_1)} f(x-r \theta) d \theta \right)}_{=\Vol(\partial B_1) f(x) \text{ by 平均値等式 (Thm. \ref{thm-Sze-2.1})}}
 \phi(r) r^{n-1}\,dr
\underalign{}{=}
\int_{0}^1 \Vol(\partial B_1) f(x) \phi(r) r^{n-1}\,dr
\\
&\underalign{}{=}
f(x)  \underbrace{\int_{0}^1\Vol(\partial B_1)  \phi(r) r^{n-1}\,dr}_{=\int_{\R^n}\eta(y)\,dy \text{ by 置換積分}}
\underalign{}{=}
f(x)
\end{align*}

\end{proof}


%重要な帰結として、$B_2$ 上の harmonic function の $L^1$-norm は、より小さい球 $B_1$ 上での関数のすべての導関数を支配する。

\begin{cor}[{\cite[Corollary 2.3]{Sze14}}]
\label{cor-Sze-2.3}
ある定数$C_k$が存在して,$f\colon B_2\to\mathbb{R}$がharmonicなら
\[
\sup_{B_1}\lvert\nabla^k f\rvert\leq C_k\int_{B_2}\lvert f(y)\rvert\,dy
\]
特に$C^2$関数がharmonicならば$C^\infty$である. 
ただしmulti-index $l=(l_1,\ldots,l_n)$について
\[
\partial^l\coloneqq \frac{\partial}{\partial x^{l_1}}\cdots\frac{\partial}{\partial x^{l_n}}
\quad \text{and} \quad
\nabla^k f(x)
\coloneqq \left(\partial^lf\right)_{l=(l_1,\ldots,l_n), \lvert l\rvert =k}
\]と定義する. 
\end{cor}


\begin{proof}
Corollary~\ref{cor-Sze-2.2}のような$\eta \colon \R^n \to \R^n$をとる. 
$C^2$関数$f$がharmonicならば, 
 $x\in B_1$に対し, Corollary~\ref{cor-Sze-2.2}から
 \begin{equation}
\label{eq-Sze-2.3}
\nabla^k f(x)
\underalign{\text{(Cor. \ref{cor-Sze-2.2})}}{=}
\nabla^k \int_{\mathbb{R}^n} f(y)\eta(x-y)\,dy
=
\int_{\mathbb{R}^n} f(y)\nabla^k\eta(x-y)\,dy
\end{equation}
よってmollifierによる畳み込みの議論と同じく, $\eta \in C^{\infty}$なので, $f \in C^\infty$である. 

\[
\lvert\nabla^k f(x)\rvert
\underalign{\eqref{eq-Sze-2.3}}{\leq}
\int_{\mathbb{R}^n} \lvert f(y)\rvert \cdot \lvert\nabla^k\eta(x-y)\rvert\,dy
\leq\bigl(\sup_{B_1}\lvert\nabla^k\eta\rvert\bigr)\int_{B_2}\lvert f(y)\rvert\,dy.
\]
よって$C_k=\sup\lvert\nabla^k\eta\rvert$と取れば結論が従う.
\end{proof}


\begin{cor}
[Liouville's theorem {\cite[Corollary 2.4]{Sze14}}]
\label{cor-Sze-2.4}
%関数 $f\colon \mathbb{R}^n\to\mathbb{R}$ が sub-linear growth をもつとは
%\[
%\lim_{R\to\infty} \frac{\sup_{B_R}\lvert f\rvert}{R}=0
%\]
%が成り立つことをいう。
 $f$が$\mathbb{R}^n$上でharmonicかつ
 \[
\lim_{R\to\infty} \frac{\sup_{B_R}\lvert f\rvert}{R}=0
\]
を満たすならば, $f$は定数. 
\end{cor}
条件$\lim_{R\to\infty} \frac{\sup_{B_R}\lvert f\rvert}{R}=0$を"sub-linear growthをもつ"ともいう. 

\begin{proof}

$r>0$に対し$f_r(x) \coloneqq f(rx)$とおくと,これもharmonicである. 
Corollary~\ref{cor-Sze-2.3}より$r$によらない定数$C'$があって, 
\begin{equation}
\label{eq-Sze-2.4}
\lvert\nabla f_r(0)\rvert
\underalign{\text{(Corollary~\ref{cor-Sze-2.3})}}{ \leq }
C_1\int_{B_2} \lvert f_r(x)\rvert\,dx \leq C'\sup_{B_2}\lvert f_r\rvert=C'\sup_{B_{2r}}\lvert f\rvert,
\end{equation}
一方$\nabla f_r(0)=r\nabla f(0)$なので,
\[
\lvert\nabla f(0)\rvert
=\frac{\lvert\nabla f_r(0)\rvert}{r}
\underalign{\eqref{eq-Sze-2.4}}{\leq} C' \frac{\sup_{B_{2r}}\lvert f\rvert}{r}  \underset{r \to \infty}{\longrightarrow} 0
\]
よって$\nabla f(0)=0$. $f$を平行移動することにより, 任意の$x$について$\nabla f(x)=0$が成り立ち, $f$は定数である. 
\end{proof}

\subsection{楕円型微分作用素}

\begin{defn}
\label{defn-Sze-elliptic}
$\R^n$上の領域$\Omega\subset\mathbb{R}^n$について, $\Omega$上の
2階線形微分作用素
\[
L(f)=\sum_{j,k=1}^n a_{jk}\frac{\partial^2 f}{\partial x^j\partial x^k}
+\sum_{l=1}^n b_l\frac{\partial f}{\partial x^l}+cf,
\]
で$f,a_{jk},b_l,c\colon \Omega\to\mathbb{R}$は$\Omega$上の関数で,
$a_{jk},b_l,c$はすべて滑らかなものを考える. 
\begin{itemize}
\item $L$が楕円型であるとは, 係数行列が対称(つまり$a_{jk}=a_{kj}$)かつ, 行列$(a_{jk})$が正定値であること. 
\item $L$が一様楕円型であるとは, ある定数$\lambda,\Lambda>0$があって, 任意のベクトル$v$について
\[
\lambda\lvert v\rvert^2\leq \sum_{j,k=1}^n a_{jk}(x)v^jv^k \leq \Lambda\lvert v\rvert^2,\qquad \text{for all }x\in\Omega
\]
であること. 
\end{itemize}
\end{defn}
$\sum_{j,k=1}^n a_{jk}(x)v^jv^k $のことを$L$のprincipal symbolとも言う. 
$L$が楕円型であり, コンパクト多様体$M$について, $\Omega$は座標近傍ならば, 一様楕円型である. 

\begin{defn}
Definition~\ref{defn-Sze-elliptic}の状況下で, 
 $\Omega$上局所可積分な関数$f \in L^{1}_{\loc}(\Omega)$が方程式$L(f)=g$の弱解であるとは,
\[
\int_\Omega f\,L^*(\varphi)\,dV=\int_\Omega g\varphi\,dV
\quad \forall \varphi \in C^{\infty}_{c}(\Omega)
\]
が成り立つこと. 
ここで$C^{\infty}_{c}(\Omega)$はsupportがコンパクトな滑らかな関数とし, 
$dV$は$\mathbb{R}^n$上の通常の体積測度とする. 

ただし形式随伴$L^*$を
\[
L^* \colon C^{\infty}_{c}(\Omega) \to C^{\infty}_{c}(\Omega)
\quad
L^*(\varphi )\coloneqq \sum_{j,k=1}^n \frac{\partial^2}{\partial x^j\partial x^k}(a_{jk}\varphi )
-\sum_{l=1}^n \frac{\partial}{\partial x^l}(b_l \varphi )+c\varphi 
\]
と定義する. 
\end{defn}

随伴は,もし$f$が$L(f)=g$の弱解でしかも実際に滑らかならば,部分積分により通常の意味でも$L(f)=g$が成り立つように定義されている.
楕円型作用素の基本的性質として,弱解は自動的に滑らかになる.

\begin{thm}[{\cite[Theorem 2.5]{Sze14}}]
\label{thm-Sze-2.5}
$L$は滑らかな係数をもつ線形楕円型作用素, $g$は滑らかな関数ならば, 
方程式$L(f)=g$の弱解$f$も滑らか. 
\end{thm}
この定理をもう少し深掘りしよう. 
\footnote{\cite[Theorem 2.5]{Sze14}には証明の概略として次が書かれていた. 
「証明の一つの方法は,まず畳み込みを用いて$f$の平滑化$f_\varepsilon$を構成し,ついで$\varepsilon$に依らない形で種々のSobolev空間における$f_\varepsilon$の評価を導くことである.その後Sobolev embedding theoremにより,$\varepsilon\to 0$のときの$f_\varepsilon$の極限である$f$が滑らかであることが従う.」 多分この方法でも示せると思う.}


\subsection{Sobolev空間とSobolev埋め込み定理}


Sobolev埋め込みについてまとめておく. 

\begin{defn}
\(\Omega\subset \mathbb{R}^N\)を領域, \(m\in\mathbb{Z}_{\geq 0}\), \(1\leq p\leq\infty\)とする. 多重指数\(\beta=(\beta_1,\ldots,\beta_N)\)に対し\(\lvert\beta\rvert=\beta_1+\cdots+\beta_N\), \(\partial^{\beta}=\partial_1^{\beta_1}\cdots\partial_N^{\beta_N}\)と書く. 

\(u\in L^1_{\mathrm{loc}}(\Omega)\)の弱微分\(v=\partial^{\beta} u\)とは,
\[
\int_\Omega u\partial^{\beta}\varphi\,dx=(-1)^{\lvert\beta\rvert}\int_\Omega v\varphi\,dx
\quad
\forall \varphi\in C_c^\infty(\Omega)
\]
が成り立つことをいう.
\end{defn}

\begin{defn} 
Sobolev空間を
\[
W^{m,p}(\Omega)\coloneqq \{u\in L^p(\Omega)\mid \partial^{\beta} u\in L^p(\Omega)\text{ for all }\lvert\beta\rvert\leq m\}
\]
\begin{itemize}
\item  \(1\leq p<\infty\)のときノルムは
\(
\lVert u\rVert_{W^{m,p}(\Omega)}\coloneqq \left(\sum_{\lvert\beta\rvert\leq m}\lVert \partial^{\beta} u\rVert_{L^p(\Omega)}^p\right)^{1/p}
\)
\item  \(p=\infty\)のときは
\(
\lVert u\rVert_{W^{m,\infty}(\Omega)}\coloneqq \sum_{\lvert\beta\rvert\leq m}\lVert \partial^{\beta} u\rVert_{L^\infty(\Omega)}
\)
\end{itemize}
を備えた, 完備ノルム空間(Banach空間)とする. 
覚書として, 
\[
W^{m,p} \text{について}
\quad \text{$m$が微分の回数}
\quad \text{$p$が$L^p$の部分.}
\]
\end{defn}
$W^{m, p}$を$L^{p}_{m}$と書くこともある. 多分幾何学の論文ではこっちが多い気がする.
$W^{m,2}$を$H^{m}$ともかく. これは多分Hilbert空間になるから重要視するのだと思う. 

ノルム空間の用語も復習しておく.
\begin{defn}
\label{defn-norm-cpt}
ノルム空間\(X,Y\)について, 
\begin{itemize}
\item \(X\hookrightarrow Y\)が連続埋め込みであるとは, ある\(C>0\)が存在して
\(\lVert u\rVert_Y\leq C\lVert u\rVert_X
\)
がすべての\(u\in X\)に対して成り立つことを 
\item  \(X\hookrightarrow Y\)がコンパクト埋め込みであるとは, \(X\)の任意の有界列から\(Y\)で強収束する部分列が取れることをいう.
\end{itemize}
\end{defn}
以下の定理はSobolev空間の基本的な定理である. 

\begin{thm}[Sobolev埋め込み定理]
\label{thm-Sobolev}
\(\Omega\subset\mathbb{R}^N\)を有界Lipschitz領域とする. \(m\geq 1\), \(1\leq p<\infty\)とする.

\begin{itemize}
\item  \(mp<N\)のとき
\(
W^{m,p}(\Omega)\hookrightarrow L^q(\Omega) \)\(\left(1\leq q\leq \frac{Np}{N-mp}\right)
\)
が連続埋め込み.
\item \(mp=N\)のとき
\(
W^{m,p}(\Omega)\hookrightarrow L^q(\Omega)\)\( (1\leq q<\infty)
\)
が連続埋め込み.
\item  \(mp>N\)のとき, \(m-\ell-\frac{N}{p}>\alpha\), \(\ell\in\mathbb{Z}_{\geq0}\), \(0<\alpha<1\)なら
\[
W^{m,p}(\Omega)\hookrightarrow C^{\ell,\alpha}(\overline{\Omega})
\]
が連続埋め込みである. 
特に
\(
m-\frac{N}{p}>0
\)
なら
\[
W^{m,p}(\Omega)\hookrightarrow C^0(\overline{\Omega})
\]
\end{itemize}

\end{thm}
感覚として微分の回数$m$がめちゃくちゃ大きければ, 滑らかになりやすい. 

\begin{thm}[Rellich--Kondrachovの定理]
\label{thm-Rellich}
\(\Omega\subset\mathbb{R}^N\)を有界Lipschitz領域とする. \(m\geq 1\), \(1\leq p<\infty\)とする. 

\begin{itemize}
\item \(mp<N\)のとき
\(
W^{m,p}(\Omega)\hookrightarrow L^q(\Omega)\)\( \left(1\leq q<\frac{Np}{N-mp}\right)
\)
がコンパクト埋め込みである.
\item  \(mp=N\)のとき
\(
W^{m,p}(\Omega)\hookrightarrow L^q(\Omega)\)\( (1\leq q<\infty)
\)
がコンパクト埋め込みである. 
\item \(mp>N\)のとき, \(m-\ell-\frac{N}{p}>\alpha\)なら
\(
W^{m,p}(\Omega)\hookrightarrow C^{\ell,\alpha}(\Omega)
\)
がコンパクト埋め込みである.
\end{itemize}
\end{thm}
$C^{\ell,\alpha}(\Omega)$はH\"older空間である. この後のDefinition~\ref{defn-Holder}を参照. 
感覚として微分の回数$m$がめちゃくちゃ大きければ, $W^{m,p}(\Omega)$の有界関数列について, 部分点列取り直して, 収束先を$C^{\ell,\alpha}(\Omega)$から取れる. (その収束先はある程度滑らかである. )

楕円型微分作用素に関して次が成り立つ. 
\begin{thm}
[局所楕円型正則性 {\cite[Chapter 4, Theorem 3.1]{Mel17}}]
\label{thm-local-elliptic-regularity}
$\Omega\subset\R^n$を領域とし, $L$を$\Omega$上の滑らかな係数をもつ
$2$階楕円型微分作用素とする. 
$U\Subset V\Subset\Omega$, $s,M\in\R$とする.

$u$を$\Omega$上のdistributionとし,
\(
L(u)\in W^{s,2}_{\loc}(\Omega)
\)
と仮定する. この時
\(
u\in W^{s+2,2}_{\loc}(\Omega)
\)
である.

さらに$u\in W^{M,2}(V)$ならば, ある定数$C>0$が存在して
\begin{equation}
\label{eq-local-elliptic-estimate}
\lVert u\rVert_{W^{s+2,2}(U)}
\leq
C\left(
\lVert L(u)\rVert_{W^{s,2}(V)}
+
\lVert u\rVert_{W^{M,2}(V)}
\right)
\end{equation}
が成り立つ.
ここで$C$は$L,U,V,s,M$に依存するが$u$には依存しない.
\end{thm}
これに関しては省略. 定数係数の場合はFourier変換で, そうでない場合はその場合に帰着させる. 
また$s, M \in \R$でもSobolev空間は定義できる. (同じくMelroseのレクチャーノートに定義がある)

\begin{proof}[Proof of Theorem~\ref{thm-Sze-2.5}]
$f\in L^1_{\loc}(\Omega)$なので$f$は$\Omega$上のdistributionとみなせる. 
いま$g\in C^\infty(\Omega)$なので, 任意の$s\in\Z_{\geq 0}$について
\(
g\in W^{s,2}_{\loc}(\Omega)
\)
である. 
$L$は滑らかな係数をもつ$2$階楕円型微分作用素なので, 
局所楕円型正則性(Theorem~\ref{thm-local-elliptic-regularity})から
\(
f\in W^{s+2,2}_{\loc}(\Omega)
\)
を得る. $s$は任意なので, Sobolev埋め込み(Theorem~\ref{thm-Sobolev})から
$f\in C^\infty(\Omega)$を得る.
\end{proof}

%%%%%%%%%%%
\begin{comment}
\begin{proof}[Proof of Theorem~\ref{thm-Sze-2.5}]
厳密に証明できるわけもないので, chatGPTに聞いて腑に落ちた内容を書く. 
まず次を認める. 
\begin{claim}[locally elliptic estimate]
\label{claim-local-ES}
$U \Subset V \Subset \Omega$について, ある定数$C>0$があって, 任意の$u \in L^{1}_{\loc}(\Omega)$と$s \geq 0$について
\begin{equation}
\label{eq-local-ES}
\lVert u\rVert_{W^{s+2, 2}(U)}
\leq C(\lVert Lu\rVert_{W^{s+2, 2}(V)} + \lVert u\rVert_{L^{2}(V)})
\end{equation}
\end{claim}
これは$\R^n$上のFourier変換を使って示す. 
これを認めると, 次のように議論できる. 


まず$f \in W^{2,2}_{\loc}$として良い\footnote{ここには$s \in \R$についてSobolev空間を定義して$\lVert u\rVert_{W^{s+2, 2}(U)}\leq C(\lVert Lu\rVert_{W^{s+1, 2}(V)} + \lVert u\rVert_{W^{s+1, 2}(V)})$を使う. $f \in L^{1}_{\loc}\Rightarrow f \in W^{-n, 2}(U) \Rightarrow f \in W^{2, 2}(U)$と議論できる. ただしここにもmollifierを使うので, $Lu$の滑らかさや$L$がellipticを使っている.}



$\rho$をmollifier, つまり$\rho \in C^{\infty}_{c}(\Omega)$で$\rho \geq 0$かつ
$\int_{\R^n}\rho\,dV=1$となるものをとる. 
$\rho_{\varepsilon}(x)\coloneqq\varepsilon^{-n}\rho(\frac{x}{\varepsilon})$とする. 
mollifierとの畳み込みを$f_{\varepsilon}$とする. 
もし任意の$l \in \N$について, 
\begin{equation*}
\sup_{0 < \varepsilon \ll 1} \lVert f_{\varepsilon}\rVert_{C^{l}} < + \infty
\end{equation*}
が言えたとすると, Ascoli-Arzeraの定理 (Theorem~\ref{thm-Arzela-Ascoli})と対角線論法から, 
$f_{\varepsilon_k} \to f$となる$C^\infty$収束部分列が取れ, $f$が$C^\infty$であることがわかる. 

さて$g$とmollifierとの畳み込みを$g_{\varepsilon}$とすると
\[
L(f_{\varepsilon}) = g_{\varepsilon} \in C^{\infty}_{c}(\Omega)
\]
となる. 
$f \in W^{2,2}_{\loc}$と$g$が滑らかであるので, 
\begin{equation}
\label{eq-gvarepsilon}
\sup_{0 < \varepsilon \ll 1}\lVert f_{\varepsilon}\rVert_{L^{2}(V)} < + \infty
\quad 
\sup_{0 < \varepsilon \ll 1}\lVert g_{\varepsilon}\rVert_{W^{s+2, 2}(V)} < + \infty
\end{equation}
である. 
よってある定数$D>0$があって任意の$\varepsilon>0$について, 
\[
\lVert f_{\varepsilon}\rVert_{W^{s+2, 2}(U)}
\underalign{\eqref{eq-local-ES}}{\leq}
 C(\lVert g_{\varepsilon}\rVert_{W^{s+2, 2}(V)} + 
\lVert f_{\varepsilon}\rVert_{L^{2}(V)})
\underalign{\eqref{eq-gvarepsilon} 
}{\leq} D
\]
Theorem~\ref{thm-Sobolev}から, $s$を$l$に応じて大きくとれば, $W^{m,p}(\Omega)\hookrightarrow C^{\ell,\alpha}(\overline{\Omega})$は連続埋め込みになる. 
よってDefinition~\ref{defn-norm-cpt}より言えた. 
\end{proof}
\end{comment}
%%%%%%%%%%%%%%%%%%%%%%%%%%%%%%%


\subsection{Schauder評価}


\begin{defn}
\label{defn-Holder}
 $\Omega\subset\mathbb{R}^n$を$C^1$級境界をもつ有界開集合とする. ($\mathbb{R}^n$の開球としてもよい)
 \begin{itemize}
 \item  $l=(l_1,\ldots,l_n)$はmulti-indexについて
\(
\partial^l\coloneqq \frac{\partial}{\partial x^{l_1}}\cdots\frac{\partial}{\partial x^{l_n}}
\)
 \item $\alpha\in(0,1)$に対し, $\Omega$上の関数$f$の$C^\alpha$ H\"older係数を
\[
\lvert f\rvert_{C^\alpha}\coloneqq \sup_{x,y\in\Omega,\ x\ne y}\frac{\lvert f(x)-f(y)\rvert}{\lvert x-y\rvert^\alpha}
\]
\item $k\in\mathbb{N}$および$\alpha\in(0,1)$に対し$C^{k,\alpha}$-norm(H\"older norm) を
\[
\lVert f\rVert_{C^{k,\alpha}}
\coloneqq \sup_{\lvert l\rvert\leq k,\ x\in\Omega}\lvert\partial^l f(x)\rvert
+\sup_{\lvert l\rvert=k}\lvert\partial^l f\rvert_{C^\alpha},
\]
\item H\"older空間を次で定める. 
\[
C^{k,\alpha}(\Omega)
\coloneqq \left\{
f  \colon \Omega \to \R \mid  \lVert f\rVert_{C^{k,\alpha}} < + \infty\right\}
\]
ここで$f \in C^{k,\alpha}(\Omega)$は$k$回連続微分可能な関数である. H\"older normにおいて, $C^{k,\alpha}(\Omega)$は完備, つまり$C^{k,\alpha}$ normに関する任意のCauchy列は$C^{k,\alpha}$で収束する. 
 \end{itemize}
\end{defn}

\begin{thm}[Arzel\`a--Ascoli theorem {\cite[Theorem 2.6]{Sze14}}]
\label{thm-Arzela-Ascoli}
$\Omega$を有界集合とし,$u_k\colon \Omega\to\mathbb{R}$を関数列で$\lVert u_k\rVert_{C^{k,\alpha}}<C$がある定数$C$に対して成り立つとする.このとき$l+\beta<k+\alpha$を満たす任意の$l,\beta$に対し,$u_k$の部分列で$C^{l,\beta}$において収束するものが存在する.
\end{thm}
かっこよくいうと$C^{k, \alpha} \hookrightarrow C^{l, \beta}$がコンパクト埋め込みである. 
これは厳密には普通の場合のArzel\`a--Ascoli theoremのCorollaryである. まず普通の場合のArzel\`a--Ascoli theoremを思い出す. 
\begin{defn}
$K$を集合とし, $\{f_j\}_{j=1}^{\infty}$を$K$上の実数値関数列とする.
\begin{enumerate}
\item $\{f_j\}_{j=1}^{\infty}$が$K$上で一様有界であるとは, ある定数$C>0$が存在して, 任意の$j\geq 1$と任意の$x\in K$に対して
\(
  \lvert f_j(x)\rvert\leq C
\)
が成り立つこと
\item $\{f_j\}_{j=1}^{\infty}$が$K$上で同程度連続であるとは, 任意の$\varepsilon>0$に対して, ある$\delta>0$が存在して, 任意の$j\geq 1$と任意の$x,y\in K$に対して
\[
  \lvert x-y\rvert<\delta
  \quad\Longrightarrow\quad
  \lvert f_j(x)-f_j(y)\rvert<\varepsilon
\]
が成り立つこと
\end{enumerate}
\end{defn}

\begin{thm}[Arzel\`a--Ascoliの定理]
$K\subset \mathbb{R}^n$をコンパクト集合とし,
\(
  f_j\colon K\to \mathbb{R}
\)
を連続関数列とする. 関数列$\{f_j\}_{j=1}^{\infty}$が$K$上で一様有界かつ同程度連続であると仮定する.

このとき, ある部分列$\{f_{j_m}\}_{m=1}^{\infty}$とある連続関数
\(
  f\colon K\to \mathbb{R}
\)
が存在して,
\[
  f_{j_m}\longrightarrow f
  \quad \text{uniformly on } K
\]
が成り立つ. すなわち,
\[
  \lVert f_{j_m}-f\rVert_{C^0(K)}
  =
  \sup_{x\in K}\lvert f_{j_m}(x)-f(x)\rvert
  \longrightarrow 0
\]
\end{thm}

\begin{proof}[Proof of Theorem~\ref{thm-Arzela-Ascoli}]
$\lVert u_j\rVert_{C^{k,\alpha}(\Omega)}\leq C$と仮定する.
まず, 各多重指数$\gamma$ with $\lvert\gamma\rvert\leq k$に対して,
\[
  \lVert D^\gamma u_j\rVert_{C^0(\Omega)}\leq C
\]
である. さらに$\lvert\gamma\rvert\leq k-1$なら, 平均値の定理より
\[
  \lvert D^\gamma u_j(x)-D^\gamma u_j(y)\rvert
  \leq C\lvert x-y\rvert
  \qquad (x,y\in \Omega)
\]
が成り立つ. また$\lvert\gamma\rvert=k$なら, 仮定より
\[
  \lvert D^\gamma u_j(x)-D^\gamma u_j(y)\rvert
  \leq C\lvert x-y\rvert^\alpha
  \qquad (x,y\in \Omega)
\]
である. したがって, 各$\lvert\gamma\rvert\leq k$について
$\{D^\gamma u_j\}_j$は一様有界かつ同程度連続である.

$\overline{\Omega}$はコンパクトであるから, Arzel\`a--Ascoliの定理を
有限個の列$\{D^\gamma u_j\}_{\lvert\gamma\rvert\leq k}$に対して適用し,
対角線論法を用いることで, 部分列を取り直して
\[
  D^\gamma u_j \longrightarrow v_\gamma
  \quad \text{uniformly on } \Omega
  \qquad (\lvert\gamma\rvert\leq k)
\]
としてよい. 
$u_j \to v_0$かつ$D^\gamma u_j \to v_\gamma$なので, 
任意の$\lvert\gamma\rvert\leq k$について
\(
  v_\gamma=D^\gamma v_0
\)
である.\footnote{ここは微積分の基本定理から出る.} 
よって
\[
  u_j\longrightarrow u
  \quad \text{in } C^k(\Omega)
\]
である. 特に$u_j\to u$は$C^l$で収束する.

残るのはH\"olderノルムの収束である. まず$l<k$とする.
このとき$\lvert\gamma\rvert\leq l$なら$\lvert\gamma\rvert+1\leq k$であるから,
平均値の定理より
\[
  [D^\gamma(u_j-u)]_{C^{0,1}}
  \leq
  C\sum_{\lvert\eta\rvert=\lvert\gamma\rvert+1}
  \lVert D^\eta(u_j-u)\rVert_{C^0}.
\]
右辺は$C^k$収束により$0$に収束する. したがって任意の
$0\leq \beta\leq 1$に対して
\[
  [D^\gamma(u_j-u)]_{C^{0,\beta}}
  \leq
  \diam(\Omega)^{1-\beta}
  [D^\gamma(u_j-u)]_{C^{0,1}}
  \longrightarrow 0.
\]
従って$u_j\to u$ in $C^{l,\beta}(\Omega)$である.

次に$l=k$とする. 条件$l+\beta<k+\alpha$はこの場合
$\beta<\alpha$を意味する. すでに
\[
  D^\gamma u_j\to D^\gamma u
  \quad \text{uniformly}
  \qquad (\lvert\gamma\rvert=k)
\]
である. また$C^{k,\alpha}$ノルムの一様有界性より
\[
  [D^\gamma(u_j-u)]_{C^{0,\alpha}}\leq C'
\]
が成り立つ. ここでH\"older補間不等式
\[
  [f]_{C^{0,\beta}}
  \leq
  2^{1-\beta/\alpha}
  \lVert f\rVert_{C^0}^{1-\beta/\alpha}
  [f]_{C^{0,\alpha}}^{\beta/\alpha}
  \qquad (0\leq \beta<\alpha)
\]
を$f=D^\gamma(u_j-u)$に適用すると,
\[
  [D^\gamma(u_j-u)]_{C^{0,\beta}}
  \longrightarrow 0.
\]
従って$u_j\to u$ in $C^{k,\beta}(\Omega)$である.
以上より, $l+\beta<k+\alpha$を満たすすべての$l,\beta$に対して,
$\{u_j\}$は$C^{l,\beta}(\Omega)$で収束する部分列を持つ.
\end{proof}


\begin{thm}[Local Schauder estimates {\cite[Theorem 2.7]{Sze14}}]
\label{thm-Sze-2.7}
$\Omega$を有界領域, 
$\Omega'\subset\Omega$を$d(\Omega',\partial\Omega)>0$を満たすより小さい領域とする. 

\[
L(f)=\sum_{j,k=1}^n a_{jk}\frac{\partial^2 f}{\partial x^j\partial x^k}
+\sum_{l=1}^n b_l\frac{\partial f}{\partial x^l}+cf
\]
が滑らかな係数をもつ一様楕円型な 2階微分作用素である, つまり不等式
\begin{equation}
\label{Sze-Schauder-18}
\lambda\lvert v\rvert^2\leq \sum_{j,k=1}^n a_{jk}(x)v^jv^k \leq \Lambda\lvert v\rvert^2
\end{equation}
がある一様な$\lambda,\Lambda>0$に対して成り立つと仮定する.
この時$\alpha\in(0,1)$, $k\in\mathbb{N}$について, 
ある定数$C$が存在し, 任意の$f \in C^2$について
\[
\lVert f\rVert_{C^{k+2,\alpha}(\Omega')}\leq C\bigl(\lVert L(f)\rVert_{C^{k,\alpha}(\Omega)}+\lVert f\rVert_{C^0(\Omega)}\bigr)
\]
この$C$は$k,\alpha$,領域$\Omega,\Omega'$,$L$の係数の$C^{k,\alpha}$-norm,  一様楕円性の定数$\lambda,\Lambda$にのみ依存する. 
\end{thm}
\cite[Theorem 2.7]{Sze14}では"任意の$f \in C^2$について"の仮定はなかったがおそらく必要だと思う.
ただ結局は$f$がdistributionでいけるらしい. ただかなり苦労するため省略.  
\begin{proof}
$L$の係数が定数であり, $k=0$かつ$b_l,c=0$の場合を示す.(一般の場合はこれに帰着する)
また補間不等式を用いることで, 次の弱い主張を示せば良い.
\begin{claim}
\label{claim-weak-schauder}
Theorem~\ref{thm-Sze-2.7}の仮定のもとで, $L(f)=g$とすると, ある定数$C$があって, 
\[
\lVert f\rVert_{C^{2,\alpha}(\Omega')}\leq C\bigl(\lVert g\rVert_{C^\alpha(\Omega)}+\lVert f\rVert_{C^2(\Omega)}\bigr).
%\tag{20}
\]
もっと強く任意の$x\in\Omega$に対して$d_x=\min\{1,d(x,\partial\Omega)\}$とおくと, ある定数$C$に対して
\[
\min\{d_x,d_y\}^\alpha\frac{\lvert\partial^l f(x)-\partial^l f(y)\rvert}{\lvert x-y\rvert^\alpha}
\leq C\bigl(\lVert g\rVert_{C^\alpha(\Omega)}+\lVert f\rVert_{C^2(\Omega)}\bigr),
%\tag{21}
\]
が任意の$x,y\in\Omega$と2階偏微分$\partial^l$に対して成り立つ. 
\end{claim}
 
%\begin{proof}[Proof of Claim \ref{claim-weak-schauder}]
以下, ある定数$K,\lambda,\Lambda$で,  $a_{jk}$は$\lVert a_{jk}\rVert_{C^\alpha}\leq K$を満たし, 一様楕円性条件\eqref{Sze-Schauder-18}を満たすとして良い. 

これはLiouvilleの定理 (Corollary~\ref{cor-Sze-2.4})を用いた背理法で示す. 
$+ \infty$に単調に発散する数列$C_1, C_2 \ldots, $と
ある$\Omega$上の関数$a_{jk},f,g$が存在して以下を満たすとする\footnote{本当は$a_{jk},f,g$らも$C_i$によるので$i$で添え字をつけないといけないが, 省略する. }
\begin{itemize}
\item $L(f)=g$, つまり
\(
\sum_{j,k} a_{jk}\frac{\partial^2 f}{\partial x^j\partial x^k}=g
\)を満たす.
\item  $\lVert g\rVert_{C^\alpha},\lVert f\rVert_{C^2}\leq 1$ 
\item 
ある$p,q\in\Omega$ 
とある 2階偏微分$\partial^l$が存在して
\begin{equation}
\label{eq-Sze-22}
\min\{d_p,d_q\}^\alpha\frac{\lvert\partial^l f(p)-\partial^l f(q)\rvert}{\lvert p-q\rvert^\alpha}
= \max_{p,q \in \Omega, l} 
\left\{\min\{d_p,d_q\}^\alpha\frac{\lvert\partial^l f(p)-\partial^l f(q)\rvert}{\lvert p-q\rvert^\alpha}\right\}
= C_i
%=C,
%\tag{22}
\end{equation}
\end{itemize}
また$l$は$C_i$によらないと仮定して良い. そこで
$$
M \coloneqq \frac{\lvert\partial^l f(p)-\partial^l f(q)\rvert}{\lvert p-q\rvert^\alpha}
\ge C_i
\quad \text{and} \quad
r \coloneqq \lvert p-q\rvert
$$
とし, スケール変換した関数を
\[
\tilde f(x) \coloneqq 
\frac{r^{-2-\alpha}f(p+rx)}{M}
\quad
F(x)\coloneqq \tilde f(x)-\tilde f(0)-\sum_k x^k\partial_k\tilde f(0)-\frac12\sum_{j,k}x^jx^k\partial_j\partial_k\tilde f(0)
\]
とする. するとこの関数$F$は次の性質を満たす：
\begin{enumerate}
\item $F$は原点のまわり半径$d_p/r$の球上で定義される. 
\item $F(0)=\partial F(0)=\partial^2 F(0)=0$.  $\partial^2$は任意の二階微分である.
\item  原点を中心とする半径$d_p/(2r)$の球上で,$\lvert\partial^2 F\rvert_{C^\alpha}\leq 2^\alpha$が成り立つ.
\item  $y=\frac{q-p}{r}$とおくと$\lvert y\rvert=1$であり,
\(
\bigl\lvert\partial^l F(0)-\partial^l F(y)\bigr\rvert=1
\)
\item $F$は次を満たす. 
\[
\sum_{j,k} a_{jk}\frac{\partial^2 F(x)}{\partial x_j\,\partial x_k}=
\frac{r^{-\alpha}}{M}\bigl(g(p+rx)-g(p)\bigr).
%\sum_{j,k} a_{jk}(p+rx)\,\frac{\partial^2 F(x)}{\partial x_j\,\partial x_k}=M^{-1}r^{-\alpha}\bigl(g(p+rx)-g(p)\bigr)+M^{-1}r^{-\alpha}\sum_{j,k}\bigl(a_{jk}(p)-a_{jk}(p+rx)\bigr)\,\frac{\partial^2 f(p)}{\partial x_j\,\partial x_k}.
\]
\end{enumerate}

各々$C_i$について, それに対応する関数列$F^{(i)}, g^{(i)}$をとる. 
$y^{(i)} = \frac{q_i-p_i}{r_i}$なので, $\lvert y_i\rvert=1$より点列コンパクト性を使って, 
$y^{(i)}$はあるベクトル$y$に収束すると仮定して良い. 



また$\lVert f\rVert_{C^2}\leq 1$を仮定しているので, $i \to \infty$で$r=\lVert p - q\rVert \longrightarrow 0$より
\[
C_i
\underalign{\eqref{eq-Sze-22}}{=}
\min\{d_p,d_q\}^\alpha\frac{\lvert\partial^l f(p)-\partial^l f(q)\rvert}{\lvert p-q\rvert^\alpha}
\underalign{\text{(by $\lVert f\rVert_{C^2}\leq 1$)}}{\leq}
\left(\frac{d_p}{r}\right)^{\alpha} 
\underalign{i \to \infty}{\longrightarrow}
 \infty
\]
(1)から, $F^{(i)}$は$\frac{d_p}{r}$上で定義され, $\frac{d_p}{r}\to \infty$なので, 
 (2) と (3) から, $F^{(i)}$は$\frac{d_p}{2r}$上で一様な$C^{2,\alpha}$評価を満たす, 
 つまり
 \[
 F^{(i)} \in C^{2,\alpha}(B_{\frac{d_p}{2r}})
 \quad 
 \&
 \quad
 F^{(i)} \leq C^{2,\alpha}(B_{\frac{d_p}{2r}})
 \]
 である. 
Arzel\`a--Ascoliの定理(Theorem~\ref{thm-Arzela-Ascoli})を用いて
\[
 F^{(i)} 
  \to \exists G \in C^{2}(B_{d_p/r})
\]
とできる. 
対角線論法より\footnote{$r=1$の時に並べて, その中から$r=2$の時にに並べて...をやって, 最後に対角線上になある関数列をとってくる}
収束部分列を取り替えることで, 
上の$r$は任意の$r$として良い. 
この$G$は条件 (1)--(4) を満たす.
これより
\[
\sum_{j,k} 
a_{jk}\frac{\partial^2  F^{(i)} (x)}{\partial x_j\,\partial x_k}
\underalign{(5)}{=}
\frac{1}{M}r^{-\alpha}\bigl(\underbrace{g(p+rx)-g(p)}_{\le r^{\alpha}|x|^{\alpha}}\bigr)
%\sum_{j,k} \underbrace{a_{jk}(p+rx)}_{\eqqcolon a_{jk} \text{定数(仮定)}}\,\frac{\partial^2  F^{(i)} (x)}{\partial x_j\,\partial x_k}=\frac{1}{M}r^{-\alpha}\bigl(\underbrace{g(p+rx)-g(p)}_{\text{絶対値2以下}}\bigr)+\frac{1}{M}r^{-\alpha}\sum_{j,k}\bigl(\underbrace{a_{jk}(p)-a_{jk}(p+rx)}_{=0 \text{定数の仮定}}\bigr)\,\underbrace{\frac{\partial^2 f^{(i)}(p)}{\partial x_j\,\partial x_k}.}_{\text{絶対値1以下}}\bigr)
\underalign{i \to \infty}{\longrightarrow}
\sum_{j,k} a_{jk}\,\frac{\partial^2 G}{\partial x_j\,\partial x_k}=0.
\]
よって$a_{jk}$を用いたしかるべき行列$T$を用いて, 
$H(x)=G(Tx)$とすれば, $H$は$\mathbb R^n$全体で定義されたharmonic関数になる. 
Corollary~\ref{cor-Sze-2.3}により$H \in C^\infty$は滑らかで, 二階微分$\partial^l H$もまたharmonicである.

さらに(2) と (3) から$H$の各二階微分はLiouvilleの定理 (Corollary~\ref{cor-Sze-2.4})より定数である. 
よって$\partial^l G \equiv 0$だがこれは (4) に矛盾する.
\end{proof}

%%%%%%%%%%%%%%
\begin{comment}


(21) からは，Exercise 2.2 のものより精密な補間不等式（(21) を参照）を用いることで，必要な評価（$f$ の $C^2$-ノルムを $C^0$-ノルムで置き換えるもの）を導くことができる．


別法として，以下のように別の背理法で進めてもよい．
\begin{proof}


定理の主張と同じ仮定のもとで，ある定数 $C$ が存在して
\[
d_x^2\,\lvert\partial^l f(x)\rvert \le C\bigl(\lVert g\rVert_{C^\alpha(\Omega)}+\lVert f\rVert_{C^0(\Omega)}\bigr)
\]
がすべての $x\in\Omega$ と任意の二階微分 $\partial^l$ に対して成り立つことを示す．

背理法. 先ほどと非常によく似た議論. 
$\lVert g\rVert_{C^\alpha},\,\lVert f\rVert_{C^0}\le 1$ と仮定する．$p\in\Omega$ と multi-index $l$ を選んで
\[
d_p^2\,\lvert\partial^l f(p)\rvert=C
\]
となるようにする．ここで $C$ はこの式の取りうる最大値であるとする．
\[
\lvert\partial^l f(p)\rvert=M\ge C
\]
とおく．

リスケールした関数
\[
F(x)=M^{-1}d_p^{-2}\,f(p+d_p x)
\]
を定める．すると $F$ は次を満たす：
\begin{enumerate}
\item[(i)] $F$ は少なくとも原点を中心とする半径 $1$ の球上で定義される．
\item[(ii)] 原点を中心とする半径 $1/2$ の球上で，ある固定定数 $K$ に対して $\lVert F\rVert_{C^2}\le K$ が成り立つ．
\item[(iii)] $\lvert\partial^l F(0)\rvert=1$．
\item[(iv)] $F$ は方程式
\[
\sum_{j,k} a_{jk}(p+d_p x)\,\frac{\partial^2 F(x)}{\partial x_j\,\partial x_k}=M^{-1}g(p+d_p x)
\]
を満たす．
\item[(v)] $\lvert F\rvert\le M^{-1}d_p^{-2}\,\lVert f\rVert_{C^0}$．
\end{enumerate}
もしこのような関数 $F^{(i)}$ の族が $C$ をより大きくしながら存在するなら，$d_p\le 1$ かつ $d_p^2 M=C$ なので，性質 (iv) における係数および右辺は一様な $C^\alpha$ 評価を満たす．

前の評価 (20) と性質 (ii) から，関数 $F^{(i)}$ は原点を中心とする半径 $1/4$ の球上で一様な $C^{2,\alpha}$ 評価を満たす．したがってある部分列は $C^2$ で $B_{1/4}$ 上の極限関数 $G$ に収束し，ある二階偏微分 $\partial^l$ に対して $\lvert\partial^l G(0)\rvert=1$ を満たす（性質 (iii) による）．
しかし性質 (v) と $d_p^2 M=C\to\infty$ であることから，$G$ は $B_{1/4}$ 上恒等的に $0$ になってしまい，これは矛盾である．
\end{proof}
\end{proof}
\end{comment}
%%%%%%%%%%%%%%%%%%%%%

もっと頑張れば次を示すことができる
\begin{prop}[Local Schauder estimates 2 {\cite[Proposition 2.8]{Sze14}}]
\label{prop-Sze-2.8}
$\Omega$を有界領域, 
$\Omega'\subset\Omega$を$d(\Omega',\partial\Omega)>0$を満たすより小さい領域とする. 

\[
L(f)=\sum_{j,k=1}^n a_{jk}\frac{\partial^2 f}{\partial x^j\partial x^k}
+\sum_{l=1}^n b_l\frac{\partial f}{\partial x^l}+cf
\]
が滑らかな係数をもつ一様楕円型な 2階微分作用素である, つまり不等式
\begin{equation}
\label{Sze-Schauder-18-2}
\lambda\lvert v\rvert^2\leq \sum_{j,k=1}^n a_{jk}(x)v^jv^k \leq \Lambda\lvert v\rvert^2
\end{equation}
がある一様な$\lambda,\Lambda>0$に対して成り立つと仮定する.
この時$\alpha\in(0,1)$, $k\in\mathbb{N}$について, 
ある定数$C$が存在し, 任意の$f \in C^2$について
\[
\lVert f\rVert_{C^{k+2,\alpha}(\Omega')}\leq C\bigl(\lVert L(f)\rVert_{C^{k,\alpha}(\Omega)}+\lVert f\rVert_{L^1(\Omega)}\bigr)
\]
が成り立つ.
\end{prop}
これも結局は$f$がdistributionでいけるらしい. ただかなり苦労するため省略.  
\cite[Proposition 2.8]{Sze14}に証明がなかったのでchatGPTにつけてもらうことにした. 

\begin{proof}
Theorem~\ref{thm-Sze-2.7}から, 
\[
\lVert f\rVert_{C^0(\Omega_1)}
\leq
C\left(
\lVert L(f)\rVert_{C^\alpha(\Omega)}
+
\lVert f\rVert_{L^1(\Omega)}
\right)
\]
が$\Omega'\Subset\Omega_1\Subset\Omega$について成り立つことを示せば良い.
$\overline{\Omega_1}\subset\Omega$はcompactなので, 
ある$r>0$と有限個の$x_1,\ldots,x_N\in\Omega_1$を
\[
\Omega'\subset\bigcup_{i=1}^N B_r(x_i),
\qquad
B_{2r}(x_i)\Subset\Omega
\]
となるように取ることができる. よって$\Omega_1=B_r(x)\Subset B_{2r}(x) = \Omega$の場合に示せば良い. 

$n=\dim \Omega$とする. 
Gilbarg--Trudingerの局所maximum principle
\cite[Corollary 9.21]{GT01}を$p=1$として
$B_r(x)\Subset B_{2r}(x)$に適用すると, 
ある定数$C>0$が存在して
\begin{equation}
\label{eq-GT-local-max}
\lVert f\rVert_{L^\infty(B_r(x))}
\leq
C\left(
\lVert f\rVert_{L^1(B_{2r}(x))}
+
\lVert L(f)\rVert_{L^n(B_{2r}(x))}
\right)
\end{equation}
が成り立つ.\footnote{ここに$f \in C^2$の仮定を使っている. この定理を使うなら$f \in W^{2, n}$が必要. 有界領域上では$L^n$ノルムは$L^\infty$ノルムで抑えられる.}
ここで$r$は固定されているので, scalingによって現れる$r$のべきは
定数$C$に含めている. また
\begin{equation}
\label{eq-nbound-C0}
\lVert L(f)\rVert_{L^n(B_{2r}(x))}
\leq
\operatorname{Vol}(B_{2r}(x))^{1/n}
\lVert L(f)\rVert_{C^0(\Omega)}
\leq
C\lVert L(f)\rVert_{C^\alpha(\Omega)}.
\end{equation}
であるので, $C^0$ノルムと$L^\infty$ノルムが同じことから
\eqref{eq-GT-local-max}と\eqref{eq-nbound-C0}より従う. 
\end{proof}

%%%%%%%%%%%%%%%%%%%%%%%
\begin{comment}

これはTheorem~\ref{thm-local-elliptic-regularity}をもう少し頑張れば, 
\begin{equation*}
\lVert u\rVert_{W^{s+2, 2}(U)}
\leq C(\lVert Lu\rVert_{W^{s+2, 2}(V)} + \lVert u\rVert_{W^{-N, 2}(V))}
\end{equation*}
とでき$N>\frac{n}{2}$なら$ \lVert u\rVert_{W^{-N, 2}(V)} \leq C\lVert u\rVert_{L^{1}(V)}$であるので, Sobolev埋め込み (Theorem~\ref{thm-Sobolev})を用いれば出る. 
もしくはboundedness estimate
\(
\lVert f\rVert_{L^\infty(\Omega_1)}
\leq
C\left(
\lVert Lf\rVert_{L^\infty(\Omega_2)}
+
\lVert f\rVert_{L^1(\Omega_2)}
\right)
\)を用いればTheorem~\ref{thm-Sze-2.7}から出る. 

\begin{proof}
$k=0$の場合のみ示す. (一般の場合はGilbarg--Trudinger Chapter 6の議論がいるらしい.)
これを示すには，Theorem 2.7 の条件のもとで 
$$
\lVert f\rVert_{C^0} \lesssim \lVert f\rVert_{L^1} +  \lVert Lf\rVert_{C^\alpha}
$$
を示せば良い.  Sobolev 空間での同様の評価と Sobolev embedding theorem(後述)を用いる. 
%(特別な場合 $Lf=0$ では，この評価は harmonic 関数に対する基本的な内点評価である Corollary 2.3 を一般化する)


$\Omega' \subset \Omega$を相対コンパクトとしてSobolev型の楕円型評価から
$$
\lVert f\rVert_{W^{2,p}(\Omega')} \lesssim \lVert f\rVert_{L^1(\Omega)} +  \lVert Lf\rVert_{L^p(\Omega)}
$$
が成り立つ, 

Sobolev埋め込み\ref{thm-Sobolev}から
なら
\[
W^{2,p}(\Omega')\hookrightarrow C^0(\overline{\Omega'})
\Rightarrow 
\lVert f\rVert_{C^0 (\Omega')} 
\lesssim 
\lVert f\rVert_{W^{2,p}(\Omega')} 
\]
である. 

$U \Subset V \Subset \Omega$について, ある定数$C>0$があって, 任意の$u \in L^{1}_{\loc}(\Omega)$と$s \ge 0$について
\begin{equation*}
\lVert u\rVert_{C^2(U)} 
\le C\lVert u\rVert_{W^{s, 2}(U)}
\le C'(||Lu||_{W^{s, 2}(V)} + ||u||_{L^{2}(V)})
\le C'(||Lu||_{C^{0}(V)} + ||u||_{L^{2}(V)})
\end{equation}



よって
$$
\lVert f\rVert_{C^0 (\Omega')} 
\lesssim 
\lVert f\rVert_{W^{2,p}(\Omega')} 
なお
$$
 \lVert Lf\rVert_{L^p(\Omega)}
\underalign{\text{積分の中身を押さえる}}{\lesssim}
 \lVert Lf\rVert_{C^0 (\Omega)}
 \underalign{\text{定義から}}{\lesssim}
  \lVert Lf\rVert_{C^\alpha(\Omega)}
$$

%我々の議論からは従わない重要な点として，$$f\in C^2\quad \& L(f)=g\quad \& \text{$L$ の係数, $g$  $ \in C^{k,\alpha}$}\quad \Rightarrowf \in C^{k+2,\alpha}$$この点についてはより努力が必要であり，Gilbarg--Trudinger Chapter 6の議論がいる. 
\end{proof}
\end{comment}
%%%%%%%%%%%%%%%%%%%%%


\subsection{Riemann多様体のSchauder評価とkernelの有限次元性}
以下Riemann多様体の場合を考える. ただし多様体は常に第2可算で向きづけ可能とする.
一般の場合の H\"older空間などの定義は, Levi-Civita接続などを用いて定義する. 
が, 少し厄介なので, 今回は$M$がコンパクトRiemann多様体の場合のみを考える. 
この場合, $\mathscr{U}$という有限開被覆をとって, H\"older normを任意のテンソル$T$について
\[
\lVert T\rVert_{C^{k,\alpha}}
\coloneqq  \sup_{U \in \mathscr{U}}
\lVert T\vert_{U}\rVert_{C^{k,\alpha}(U)}
\]
と定義できる. このH\"olderノルムは, チャートによって異なるが, 同値なノルムになるので有限性などには影響がない.  
%%%%%%%%%%%%%
\begin{comment}
Riemann多様体の場合のH\"older空間を計量を用いて以下のように定義する.\footnote{流石に完備性がないとまずいのではと思ってchatGPTに聞いたところ, 完備性があってもこの定義はまずいらしい. locally H\"olderを使うか, bounded geometry conditionをかした方がいい. ここでbounded geometryとは$\inj(M,g)< \rho$である. 同義的に$\mathrm{exp}_x \colon T_xM \to M$という指数写像に関し, $r < \rho$ならば$B_{T_{x}M}(0, r)\subset T_xM$上で微分同相を与えること.}
\begin{defn}
$(M,g)$をRiemann多様体とする. 
テンソル$T$に対し
\[
\lvert T\rvert_{C^\alpha} \coloneqq \sup_{x,y \in }\frac{\lvert T(x)-T(y)\rvert}{d(x,y)^\alpha}
\]
と定義できる. 
%ただし,$\lvert T(x)-T(y)\rvert$に関しては, Levi-Civita connection に関する測地線に沿う平行移動を用いて異なる点におけるテンソルを比較している. 
ここで
\begin{itemize}
\item 上限は, 点$x,y$が一意な最短測地線で結ばれるような組$(x,y)$全体にわたる. 
\item 差$T(x)-T(y)$は, $T(y)$をLevi--Civita接続に関する最短測地線に沿って$x$へ平行移動して計算する.
\end{itemize}
そしてH\"older normを
\[
\lVert f\rVert_{C^{k,\alpha}}=\sup_M\bigl(\lvert f\rvert+\lvert\nabla f\rvert+\cdots+\lvert\nabla^k f\rvert\bigr)+\lvert\nabla^k f\rvert_{C^\alpha}
\]
と定め, H\"older空間$C^{k,\alpha}(M)$と定義する. 
\end{defn}
チャートが有限個である場合は, 例えば$M$がコンパクト多様体の場合, 
$\mathscr{U}$という有限開被覆をとって, H\"older空間やH\"older normを任意のテンソル$T$について
\[
\lVert T\rVert_{C^{k,\alpha}}
\coloneqq  \sup_{U \in \mathscr{U}}
\lVert T\vert_{U}\rVert_{C^{k,\alpha}(U)}
\]
と定義できる. 上の二つのH\"olderノルムは, チャートが有限個である場合は同値である.
\end{comment}
%%%%%%%%%%%%%%%%%%%%

\begin{defn}
Riemann多様体上の線形2回楕円型偏微分作用素とは, 各局所チャートで
2階線形微分作用素
\[
L(f)=\sum_{j,k=1}^n a_{jk}\frac{\partial^2 f}{\partial x^j\partial x^k}
+\sum_{l=1}^n b_l\frac{\partial f}{\partial x^l}+cf,
\]
とかけて行列$(a_{jk})$が正定値かつ対称($a_{jk}=a_{kj}$)となること
\end{defn}
$M$がコンパクトならば, Definition~\ref{defn-Sze-elliptic}より, 一様楕円型となるような座標チャート$V_i$によって多様体を被覆することができる. 
Theorem~\ref{thm-Sze-2.7}のlocal Schauder estimateを$U_i \Subset V_i$で使えば次を得る. 
\begin{thm}
[Schauder estimates {\cite[Theorem 2.9]{Sze14}}]
\label{thm-Schauder}
$(M,g)$をコンパクトRiemann多様体とし, $L$を$M$上の線形2回楕円型偏微分作用素とする.
任意の$k \in \N$と$\alpha\in(0,1)$に対し,ある定数$C$が存在して
\[
\lVert f\rVert_{C^{k+2,\alpha}(M)}\leq C\bigl(\lVert L(f)\rVert_{C^{k,\alpha}(M)}+\lVert f\rVert_{L^1(M)}\bigr)
\]
が成り立つ.
ここで$C$は$(M,g)$,$k$,$\alpha$,$L$の係数の$C^{k,\alpha}$-ノルム,および各チャートの一様楕円性の定義での定数$\lambda,\Lambda$にのみ依存する.

特に$L(f)$と$L$の係数が$C^{k,\alpha}$である時, $f\in C^2$ならば$f\in C^{k+2,\alpha}$である. 
\end{thm}
これのCorollaryはよく使う. 

\begin{cor}
[{\cite[Corollary 2.10]{Sze14}}]
\label{cor-finitedim-kernel}
$L$をコンパクトRiemann多様体$M$上の線形2回楕円型偏微分作用素とする.このとき$L$のkernel
\[
\ker L=\{\, f\in L^2(M)\mid f \text{ は } Lf=0 \text{ の弱解}\,\}
\]
は滑らかな関数からなる有限次元空間である.
\end{cor}
\begin{proof}
Theorem~\ref{thm-Schauder}から
$Lf=0$の任意の弱解は, $L$の係数が滑らかなので, $f$も滑らかである. 
$\ker L$が有限次元であることを示すために,$L^2$計量に関して$\ker L$の閉単位球が(点列)コンパクトであることを示す.(下のRemark~\ref{rem-norm-cpt}参照)

以下$f_k\in\ker L$を$\lVert f_k\rVert_{L^2(M)}\leq 1$を満たす関数列とする.
ここでH\"olderの不等式を思い出す. 
\begin{thm}[H\"olderの不等式]
$(X, \mu)$測度空間, 
$1 \leq p,q \leq +\infty$とし, 
$f \in L^p(X), g \in L^q(X)$, $\frac{1}{p} + \frac{1}{q}=1$とすると
\[
\int_{X}\lvert fg\rvert\,d\mu\leq
\left( \int_{X}\lvert f\rvert^p\,d\mu \right)^{\frac{1}{p}}
\cdot
\left( \int_{X}\lvert g\rvert^q\,d\mu \right)^{\frac{1}{q}}
\]
\end{thm}
特に今$M$がコンパクトなので, ある定数$C_1$があって
\[
\lVert f_k\rVert_{L^1(M)}
\leq 
\underbrace{\lVert 1\rVert_{L^2(M)}}_{<+\infty \text{ by $M$コンパクト}}
\lVert f_k\rVert_{L^2(M)}
\leq
C_1
\] 
とできる. Theorem~\ref{thm-Schauder}の
Schauder評価を適用すると, ある定数$C_2$が存在して
\[
\lVert f_k\rVert_{C^{2,\alpha}(M)}
\underalign{\text{(Thm. \ref{thm-Schauder})}}{\leq}
C\bigl(\lVert\underbrace{L(f_k)}_{=0 \text{ by 仮定 }}\rVert_{C^{2,\alpha}(M)}+\lVert f_k\rVert_{L^1(M)}\bigr)
\leq C_2
\]
を得る. Arzera-Ascoliの定理
Theorem~\ref{thm-Arzela-Ascoli}から, $C^2$収束する部分列
\[
f_{k_l}
\to \exists f \in C^2(M)
\] 
が取れる. 
今$M$がコンパクトなので, $L$の係数の関数は上に有界で, 特に
\[
\lvert L(f)\rvert_{C^0}\leq
\sum_{j,k=1}^n \lvert a_{jk}\rvert_{C^0} \lvert\frac{\partial^2 f}{\partial x^j\partial x^k}\rvert_{C^0}
+\sum_{l=1}^n \lvert b_l\rvert_{C^0} \lvert\frac{\partial f}{\partial x^l}\rvert_{C^0}+\lvert c\rvert_{C^0}\lvert f\rvert_{C^0}
\lesssim 
\lVert f\rVert_{C^2}
\]
である. 特に$L \colon C^2(M) \to C^0 (M)$はH\"olderノルムで連続である.
よって$f_{k_l}\to f \in C^2(M)$なので, $L$の連続性から$f\in \ker L$.
また$\lVert f\rVert_{L^2(M)}\leq 1$でもある. 
よって, $\ker L$の単位球内の任意の列が収束部分列をもち点列コンパクトである. 
\end{proof}
\begin{rem}
\label{rem-norm-cpt}
ノルム空間は距離空間なので
コンパクトと点列コンパクトは同値である. 
また次が成り立つ. 
\begin{lem}
ノルム空間が無限次元ならば, 閉単位球が$($点列$)$コンパクトでない
\end{lem}
\begin{proof}
$B$を閉単位球とする
$v_i \in B$で, $\lVert v_i\rVert =1$かつ$\lVert v_i - v_j \rVert \geq \frac{1}{2}$となるものを構成すれば良い. 
$v_1, \ldots, v_m$まで作れたとする. 
$W_m \coloneqq  \mathrm{Span}(v_1, \ldots, v_m)$とすると$W_m \subsetneq V$である. 
よってRieszの補題より
\[
d(v_{m+1}, W_m) \geq \frac{1}{2} \quad \lVert v_{m+1}\rVert =1
\]
が取れる. これを繰り返せば良い. 
\end{proof}
\end{rem}


\subsection{Laplacianの定義}
以下K\"ahler多様体で議論をする. 今後の表示を簡単にするために共変微分を定義する\footnote{なぜこのような定義にしたかというと岩井がRiemannian幾何学をよくわかっていないからである. ただK\"ahlerならばChern接続とLevi--Civita接続は同じなので, Levi--Civita接続をよくわかってなくてもChern接続を用いて共変微分は定義できる. 多分これは少数派である. 普通はLevi--Civita接続の方がメジャーである.同様にRiemann曲率もよくわかってないが, Chern曲率と同じなので, こっちで理解している. Hermite多様体を扱いたい人は別の定義を見てほしい. }
\begin{defn}
$(X, g)$をK\"ahler多様体とする. 
$h$を$g$が誘導する$T_X$のHermite計量とし, 
$D_h \colon A^0(T_X) \to A^1(T_X)$をChern接続とする. 

$(z^1, \ldots, z^n)$を局所座標として, 
local Chartで書くと, Chern接続 
\[
D_h \colon A^0(T_X) \to A^{1}(T_X)
\]
は$e_i\coloneqq \frac{\partial}{\partial z^i}$, $s^i$を$C^\infty$関数とすると, connection form $\theta$は$(1,0)$-formで$\theta=h^{-1} \partial h$で与えられる. これはつまり
\[
D_h (s^i e_i) = (\partial + \bar\partial)s^i \otimes e_i + s^j \theta^{i}_{j} \otimes e_i
\quad
\theta^{i}_{j} = \underbrace{h^{i \bar{l}}\partial_j h_{k \bar{l}}}_{=:\Gamma_{jk}^{i} } dz^k =\Gamma_{jk}^{i} dz^k
\]
となる. 
そこで共変微分を
\begin{align*}
\nabla_{k} &\colon A^0(T_X) \to A^0(T_X) \\
s^i e_i &\mapsto D_h(s^i e_i)(\partial_k)
=\left((\partial_k s^i dz^k+s^j \Gamma_{jk}^{i} dz^k ) \otimes e_i\right)(\partial_k)
=(\partial_k s^i +s^j \Gamma_{jk}^{i} )\otimes e_i
%=\left(\partial s^i+s^j \underbrace{h^{i \bar{l}}\partial_j h_{k \bar{l}} dz^k}_{=\theta^{i}_{j}}) \otimes e_i\right)(\partial_k)
%=\left(\partial_k s^i + s^j h^{i \bar{l}}\partial_j h_{k \bar{l}}\right)\otimes e_i
\\
\nabla_{\bar k}& \colon A^0(T_X) \to A^0(T_X)
%\quad \alpha \mapsto D_h(\alpha)(\partial_{\bar k})
\\
s^i e_i &\mapsto D_h(s^i e_i)(\partial_{\bar k})
=\left((\partial_{\bar k} s^i\,d\overline{z^k} \otimes e_i\right)(\partial_{\bar k})
=(\partial_{\bar k} s^i )\otimes e_i
\end{align*}
と定義する. 
\end{defn}

\begin{rem}
\label{rem-BG13-formula}
\cite[(3.6)]{BG13}にあった公式
\begin{equation*}
\nabla_k X^i
=
\partial_k X^i+\Gamma^i_{jk}X^j,
\quad
\nabla_{\bar{k}}X^i
=
\partial_{\bar{k}}X^i,
\quad
\nabla_k \bar{Y}^{\bar{i}}
=
\partial_k \bar{Y}^{\bar{i}},
\quad
\nabla_{\bar{k}}\bar{Y}^{\bar{i}}
=
\partial_{\bar{k}}\bar{Y}^{\bar{i}}
+
\overline{\Gamma^i_{jk}}\bar{Y}^{\bar{j}}.
\end{equation*}
は正確にいうと
\[
\nabla_k X^i e_i
=
\left(\partial_k X^i+\Gamma^i_{jk}X^j\right)e_i
\quad
\nabla_{\bar{k}}X^i e_i
=
\left(\partial_{\bar{k}}X^i\right)e_i
\]
のことである(この$i$成分を見ている.) 上の3,4つ目の式は$A^0(\overline{T_X}) \to A^1(\overline{T_X})$のChern接続を見ている.

\cite[(3.7)]{BG13}にあった公式も見ておくと, 
\[
\nabla_k a_i=\partial_k a_i-\Gamma^j_{ik}a_j,\qquad
\nabla_{\bar k}a_i=\partial_{\bar k}a_i,\qquad
\nabla_k b_{\bar i}=\partial_k b_{\bar i},\qquad
\nabla_{\bar k}b_{\bar i}
=\partial_{\bar k}b_{\bar i}
-\overline{\Gamma^j_{ik}}\,b_{\bar j}.
\]
覚書として上つき文字には$+\Gamma$して下つき文字には$-\Gamma$する. 
\end{rem}

\begin{defn}
\label{defn-cpx-Laplace}
Riemannian 多様体について, Laplacianを次で定める
\[
\Delta_{\mathrm{Riem}} f
\coloneqq  \operatorname{div}(\operatorname{grad} f)
\underalign{\text{(locally)}}{=}
g^{ab}_{\R}(\partial_a \partial_b f -\Gamma_{ab}^c \partial_c f)
\]


$X$を複素多様体として, 複素Laplacianを次で定める
\[
\Delta(f)
=
\operatorname{tr}_\omega(\sqrt{-1} \partial \bar\partial f)
\underalign{\text{(locally)}}{=}
g^{l\bar{k}}\partial_l \partial_{\bar k} f
\]
\end{defn}

\begin{lem}
\label{lem-cpx-Laplace}
K\"ahler多様体上では, 複素Laplacianは通常のRiemannian Laplacianの半分になり, 
局所正則座標を用いて
\[
\Delta f = g^{k\bar l}\nabla_k\nabla_{\bar l}f = g^{k\bar l}\partial_k\partial_{\bar l}f
\]
と書ける.特にellipticである. 
\end{lem}

%$\nabla_k(\partial/\partial \overline{z^l})=0$ を思い出すと，Riemannian の場合と対照的に，normal coordinates を用いていない場合でも偏微分を用いた表式が成り立つことが分かる．作用素を局所実座標で書き直すと，Laplace operator が楕円型であることが分かる．
\begin{proof}
K\"ahler多様体上では, 複素Laplacianは通常のRiemannian Laplacianの半分になることは調和積分論の議論から. 

$\nabla_k\nabla_{\bar l}f$の定義は
\[
\bar\partial \colon A^0(X) \to A^{0,1}(X)=A^{0}(\overline{\Omega_X})
\quad
f \mapsto D(f)=Df
\]
\[
D'_{\overline{\Omega_X}} \colon A^{0}(\overline{\Omega_X}) \to A^{1,0}(\overline{\Omega_X})
\quad
\alpha \mapsto D_{\overline{\Omega_X}}(\alpha)
\]
とした場合の
\[
\nabla_k\nabla_{\bar l}\colon
 A^{0}(X) \to A^{0}(X)
 \quad
 f \mapsto (D'_{\overline{\Omega_X}}\circ \bar\partial)(f)(\partial_k \otimes \partial_{\bar l})
\]
である. 
よって
\[
(D'_{\overline{\Omega_X}}\circ \bar\partial)(f)
=(D'_{\overline{\Omega_X}} (\partial_{\bar l}f\,d\overline{z^l})
=\partial_k \partial_{\bar l}f dz^k \otimes\,d\overline{z^l}
\]
であるので, $g^{k\bar l}\nabla_k\nabla_{\bar l}f = g^{k\bar l}\partial_k\partial_{\bar l}f$が言える. 

楕円型であるのは
\[
\partial_k \partial_{\bar k}
=
\frac{1}{4}\left(\frac{\partial^2}{\partial (x^j)^2} +  \frac{\partial^2}{\partial (y^j)^2} \right)
\]
となることから, $g_{i \bar{j}}=\delta_{i\bar{j}}$とみなせばよい.
\end{proof}

\begin{rem}
微分幾何学的な証明なら, 
localには
$
\Delta_{\mathrm{Riem}} f
=
g^{ab}_{\R}(\partial_a \partial_b f -\Gamma_{ab}^c \partial_c f)
=2g^{k\bar l}\nabla_k\nabla_{\bar l}f
$である. そして一般には
\(
\nabla_k\nabla_{\bar l}f
=
\partial_k\partial_{\bar l}f
-
\Gamma^{a}_{k\bar l}\partial_af.
\)である. K\"ahler計量ではmixed Christoffel symbol $\Gamma^{a}_{k\bar l}$が消えるから, ほしい式が得られる. 
\end{rem}


\begin{defn}
$M$をコンパクト実多様体とする. 
$d \colon A^0(M) \to A^1(M)$の形式的随伴を
$
d^{*}\colon A^1(M) \to A^0(M)
$
を$dV$をvolume formとして
\[
\int_M \langle \alpha,d f\rangle\,dV=\int_M \langle d^{*} \alpha,f\rangle\,dV
\quad 
\forall f\in A^0(M), \alpha \in A^1(M), 
\]
が成り立つように定める. 

$X$を複素多様体とする. 
$\bar \partial \colon A^0(X) \to A^{0,1}(X)$の形式的随伴
\(
\bar\partial^*\colon \Omega^{0,1}X\to C^\infty(X)
\)
を
\[
\int_X \langle \alpha,\bar\partial f\rangle\,dV=\int_X \langle \bar\partial^*\alpha,f\rangle\,dV,
\quad 
\forall f\in A^0(X), \alpha \in A^{0,1}(X), 
\]
が成り立つように定める. 
ここで$\langle\cdot,\cdot\rangle$は計量によって誘導される自然なHermitian formであり,$dV$はRiemannian volume formなので$dV=\dfrac{\omega^n}{n!}$である. 
\end{defn}
局所座標では$\alpha=\alpha_{\bar l}d\bar z^{l}$と書ける. 
任意のcompact supportを持つsmooth function $f$に対して
\[
\begin{aligned}
\int_{X}\langle \bar\partial^*\alpha,f\rangle\frac{\omega^n}{n!}
&\underalign{\text{(随伴)}}{=}
\int_{X}\langle \alpha,\bar\partial f\rangle\frac{\omega^n}{n!}\underalign{\text{(定義)}}{=}
\int_{X}
g^{k\bar l}\alpha_{\bar l}\,
\overline{\partial_{\bar k}f}
\frac{\omega^n}{n!}
=
\int_{X}
g^{k\bar l}\alpha_{\bar l}\,
\partial_k\bar f
\frac{\omega^n}{n!}
\end{aligned}
\underalign{\text{(部分積分)}}{=}
-\int_{X}
g^{k\bar l}\nabla_k\alpha_{\bar l}\,
\bar f
\frac{\omega^n}{n!}.
\]
したがって
\(
\bar\partial^*\alpha
=
-g^{k\bar l}\nabla_k\alpha_{\bar l}.
\)
さらにK\"ahler metricについて
\(
\Gamma_{k\bar l}^{\bar q}=0
\)
なので
\[
\bar\partial^*\alpha
=
-g^{k\bar l}
\left(
\partial_k\alpha_{\bar l}
-
\Gamma_{k\bar l}^{\bar q}\alpha_{\bar q}
\right)
=
-g^{k\bar l}\partial_k\alpha_{\bar l}.
\]
以上をまとめると
\begin{lem}
$ f\in A^0(X), \alpha \in A^{0,1}(X)$について
\[
\bar\partial^*\alpha = -g^{k\bar l}\nabla_k\alpha_{\bar l}
=-g^{k\bar l}\partial_k\alpha_{\bar l}
\quad
-\bar\partial^*\bar\partial f=
g^{k\bar l}\partial_k\partial_{\bar l}f
=\Delta f
\]
特に複素Laplacianは$\Delta=- \bar\partial^*\bar\partial$となる. 
\end{lem}
\begin{rem}
同様に任意の$(p,q)$-formsにHodge Laplacian
\[
\Delta=-\bar\partial^*\bar\partial-\bar\partial\bar\partial^*
\]
を与える. これは調和積分論(Subsection~\ref{subsec-harmonic})で使うLaplacianと符号が逆である. 
理由はHodge theoryのLaplacianは
\[
\int_{X}\langle\bar\partial^*\bar\partial \alpha, \beta \rangle\,dV
=\int_{X}\langle\bar\partial \alpha, \bar \partial \beta \rangle\,dV \geq 0
\]
が成り立つようにしているから. 
これは部分積分において
\[
\int_{\R} (f' g)\,dx
= - \int_{\R}f g'\,dx
= \int_{\R}-f g'\,dx
\]
にしていることに似ている. 
一方今使っているLaplacian$\Delta=- \bar\partial^*\bar\partial$は楕円型作用素にするためにこの符号を用いている. 
\end{rem}

\subsection{Poisson方程式の解の存在}
\begin{thm}[Poisson方程式の解の存在]
$(X,\omega)$をコンパクトK\"ahler多様体とし, 
$\rho\colon X\to\mathbb R$を$C^\infty$級関数で
\(
\int_X \rho\,dV = 0.
\)
を満たすものとする. ただし$dV=\frac{\omega^n}{n!}$とする. 

この時ある$C^\infty$級関数$f\colon X\to\mathbb R$で, $\Delta f=\rho$となるものが存在する. 
\end{thm}
逆に$\Delta f=\rho$ならば, 部分積分によって, 
\[
\int_X \Delta f,\,dV 
\underalign{\text{(定義)}}{=} 
-\int_X \langle \bar\partial^*\bar\partial f,1\rangle\,dV
\underalign{\text{(随伴)}}{=} 
-\int_X \langle \bar\partial f,\bar\partial 1\rangle\,dV
=0.
\]
であるので, $\int_X \rho\,dV = 0$という条件は必要である. 
\begin{proof}
まず変分問題, つまり汎関数
\[
E \colon W^{1,2}(X) \to \R \quad
E(f)=\int_X\left(\frac12\underbrace{\lvert\nabla f\rvert^2}_{=g^{k\bar l}\nabla_k f\,\nabla_{\bar l}f=g^{k\bar l}\partial_k f\,\bar\partial_{\bar l}f}+\rho f\right)\,dV
\]
を最小化する関数$f$を,制約$\int_X f\,dV=0$のもとで探す.

\begin{claim}
\label{claim-E-variance}
ある定数$\varepsilon,\,C$が存在して
\[
E(f)\geq \varepsilon\lVert f\rVert^{2}_{W^{1,2}}-C
\]
が$\int_X f\,dV=0$である$f \in W^{1,2}(X)$に対して成り立つ.
\end{claim}
\begin{proof}[Proof of Claim~\ref{claim-E-variance}]
Poincar\'e inequalityを思い出す
\begin{thm}[Poincar\'e inequality]
\label{thm-poincare-inequality}
\((M,g)\)をコンパクトRiemann多様体とする. 
ある定数\(C_P=C_P(M,g)>0\)が存在して,任意の\(f\in W^{1,2}(M)\)に対して
\[
\int_M \lvert f-f_M\rvert^2\,dV_g
\leq
C_P
\int_M \lvert\nabla f\rvert^2\,dV_g
\]が成り立つ.ただし
\(
f_M
\coloneqq 
\frac{1}{\Vol(M)}
\int_M f\,dV_g
\)を\(f\)の平均値とする. 
\end{thm}
これにより, 特に今仮定から$ \int_X f\,dV=0$なので, 
\begin{equation}
\label{eq-Poincare}
\lVert f\rVert_{L^2(X)}\lesssim  \lVert \nabla f \rVert_{L^2(X)} 
\end{equation}
以上よりYoung不等式$aB \leq \frac{1}{4}a^2 + b^2$を用いると, 
\begin{equation}
\label{eq-Poincare-2}
\lvert\int_X\rho f\,dV\rvert
\underalign{\text{(Cauchy--Schwarz)}}{\leq}
 \lVert \rho \rVert_{L^2(X)}  \lVert f \rVert_{L^2(X)} 
 \underalign{\text{\eqref{eq-Poincare}}}{ \leq}
C  \lVert \rho \rVert_{L^2(X)} \lVert \nabla f \rVert_{L^2(X)} 
 \underalign{\text{(Young)}}{ \leq  }
 \frac{1}{4} \lVert \nabla f \rVert^2_{L^2(X)}  + C^2  \lVert \rho \rVert^2_{L^2(X)} 
\end{equation}
以上より定数$C^2$を$C$に取り替えて, 
\begin{equation}
\label{eq-Poincare-3}
\int_X\left(\frac12 \lvert\nabla f\rvert^2 +\rho f\right)\,dV
\underalign{\eqref{eq-Poincare-2}}{\geq}
\frac{1}{2} \lVert \nabla f \rVert_{L^2(X)}^2 
- \frac{1}{4} \lVert \nabla f \rVert_{L^2(X)}^2  - C  \lVert \rho \rVert_{L^2(X)}^2 
=\frac{1}{4} \lVert \nabla f \rVert_{L^2(X)}^2  - C 
\end{equation}
もう一回Poincar\'e inequalityを使って
\[
\lVert f\rVert^{2}_{W^{1,2}}
\underalign{\text{(定義)}}{=}
\lVert f\rVert^{2}_{L^2}
+
\lVert\nabla f \rVert^{2}_{L^2}
\underalign{\eqref{eq-Poincare}}{\lesssim}
\lVert\nabla f \rVert^{2}_{L^2}
\]
よって\eqref{eq-Poincare-3}と合わせて, 
\[
E(f)
=
\int_X\left(\frac12 \lvert\nabla f\rvert^2 +\rho f\right)\,dV
\underalign{\eqref{eq-Poincare-3}}{\geq} 
\varepsilon\lVert f\rVert^{2}_{W^{1,2}}-C
\]
\end{proof}

今$\inf E(f)$に近づくような列$f_j$を考える. 
 Claim~\ref{claim-E-variance}より$\lVert f_j\rVert_{W^{1,2}}$は有界なので, Hilbert空間論の一般論から
 \[
 f_{j_k} \underset{\text{弱収束}}{\longrightarrow} \exists F \in W^{1,2}
 \]
 を取れる. \footnote{弱収束とは任意の連続汎函数を当てて収束. 強収束はノルム収束. Hilbert空間なら弱収束は任意の内積をとって収束と同じ. }
今連続汎函数$f \mapsto\int_X{}f\,dV$と
$f \mapsto\int_X{}\rho f\,dV$を考えると, 弱収束の定義から
\begin{equation}
\label{eq-weak-converge-1}
\int_X F\,dV
\underalign{\text{(弱収束の定義)}}{=}
\lim_{j \to \infty}\int_X f_{j_k}\,dV
\underalign{\text{($f_{j_k}$の定義)}}{=}
0
\quad \text{and} \quad
\int_X \rho F\,dV
\underalign{\text{(弱収束の定義)}}{=}
\lim_{j \to \infty}\int_X \rho f_{j_k}\,dV
\end{equation}
一方
ノルムの下半連続性\footnote{一般に弱収束するならば, ノルムは下半連続}より
\begin{equation}
\label{eq-weak-converge-2}
 f_{j_k} \underset{\text{弱収束}}{\longrightarrow} F \in W^{1,2}
 \Rightarrow
 \nabla f_{j_k} \underset{\text{弱収束}}{\longrightarrow} \nabla F \in L^{2}
\Rightarrow
\lVert  \nabla F\rVert_{L^{2}}
\leq
\liminf_{k \to\infty} \lVert \nabla f_{j_k}\rVert_{L^{2}}
\end{equation}
これを用いると, $\int_X F\,dV=0$から
\begin{align*}
E(F)
&\underalign{\text{($F$の定義)}}{=}
\int_X\left(\frac12\lvert\nabla F\rvert^2+\rho F\right)\,dV
\\
&\underalign{\eqref{eq-weak-converge-1}\text{ \& } \eqref{eq-weak-converge-2}}{\leq}
\liminf_{k \to\infty}
\int_X\left(\frac12\lvert\nabla f_{j_k}\rvert^2+\rho f_{j_k} \right)\,dV
=
\liminf_{k \to\infty} E(f_{j_k})
\underalign{\text{($f_{j_k}$の定義)}}{\leq}E(F)
\end{align*}
よって$F$は$E$の最小化関数である. 

あとはこの$F$が滑らかで$\Delta F=\rho$であることを示す. 
任意の滑らかな平均0の関数$\varphi$について
\[
H\colon (-\varepsilon, \varepsilon) \to \R
\quad
H(t)\coloneqq E(F +t \varphi)
\]を考えると, これは$t=0$で最小値を取る. よって微分は0となる

実際$t=0$で微分すると
\begin{align*}
0=
\frac{d}{dt}\big\vert_{t=0}H
&=
\frac{d}{dt}\big\vert_{t=0}\int_{X}\left\{ 
\frac{1}{2}\lvert\nabla F + t \nabla \varphi\rvert^2 + \rho\cdot (F + t \varphi)\right\}\,dV
\\
&=
\int_{X}\left\{ \langle \nabla F,  \nabla \varphi \rangle + \rho \varphi\right\}\,dV=
\int_{X}(-\Delta F + \rho )\varphi\,dV
\end{align*}
これが任意の滑らかな平均0の関数$\varphi$について成り立つので, $F$は
\(
\Delta F=\rho
\)
の弱解である. 以上よりTheorem~\ref{thm-Sze-2.5}から$F$は滑らかである.
\end{proof}

\subsection{コンパクトRiemann多様体上の楕円型偏微分作用素の解の有限次元性}
この内容はchatGPTに証明と議論を埋めてもらった. 

%Fredholm operators などの functional analysis の道具を用いると，コンパクト多様体上の H\"older 空間の間での linear elliptic operators の写像性を記述する次のかなり一般的な定理が得られる．

\begin{thm}
\label{thm-fredholm}
$M$をコンパクトRiemann多様体, 
 $L$を係数が$C^\infty$である線形2回楕円型偏微分作用素とする.
 この時$L$は以下の同型を誘導する.
\[
L\colon (\ker L)^\perp\cap C^{k+2,\alpha}\ \to\ (\ker L^*)^\perp\cap C^{k,\alpha}.
\]
さらに$\ker L$は有限次元である. 

特に任意の$k\geq 0$と$\alpha\in(0,1)$について, 
$\rho\in C^{k,\alpha}(M)$が$L^2$内積で$\rho\perp \ker L^*$となるならば, 
$f\in C^{k+2,\alpha}$で$f\perp \ker L$かつ$Lf=\rho$となるものがただ一つ存在する. 
\end{thm}

\begin{proof}
局所楕円型評価(Theorem~\ref{thm-local-elliptic-regularity})において$s=0$とすることで, ある定数$C>0$が存在して任意の$f\in W^{2,2}(M)$について
\begin{equation}
\label{eq-global-Sobolev-elliptic}
\lVert f\rVert_{W^{2,2}(M)}
\leq
C\bigl(
\lVert Lf\rVert_{L^2(M)}
+
\lVert f\rVert_{L^2(M)}
\bigr)
\end{equation}
が成り立つ. ここに$M$がコンパクトを用いている. 

(1). まず$\ker L$が有限次元であることを示す. 
Theorem~\ref{thm-Sze-2.5}から$Lf=0$の弱解は滑らかなので, 
$\ker L\subset C^\infty(M)$である. 
よって$\ker L\subset L^2(M)$なので, Corollary~\ref{cor-finitedim-kernel}より有限次元である. 

ただ後の補足(Remark~\ref{rem-fredholm})のことも考え局所楕円型評価\eqref{eq-global-Sobolev-elliptic}からでも出ることを示しておく. 
$\ker L$が無限次元であると仮定する. 
すると$L^2$内積に関する正規直交列
\(
f_i\in\ker L\)かつ
\(\lVert f_i\rVert_{L^2(M)}=1
\)となるものを取ることができる. 
$Lf_i=0$なので, \eqref{eq-global-Sobolev-elliptic}から
\[
\lVert f_i\rVert_{W^{2,2}(M)}
\underalign{\eqref{eq-global-Sobolev-elliptic}}{\leq}
C\lVert f_i\rVert_{L^2(M)}
=C
\]
となる. 
一方Rellich--Kondrachovの定理(Theorem~\ref{thm-Rellich})より
\(
W^{2,2}(M)\hookrightarrow L^2(M)
\)
はコンパクト埋め込みである. 特にDefinition~\ref{defn-norm-cpt}
より$f_i$は$L^2(M)$で収束する部分列を持つ. 
しかし$i\neq j$について
\[
\lVert f_i-f_j\rVert_{L^2(M)}^2
=
\lVert f_i\rVert_{L^2(M)}^2
+
\lVert f_j\rVert_{L^2(M)}^2
=
2
\]
であり, これは矛盾である. 
よって$\ker L$は有限次元である. 
同じ議論を$L^*$に適用すれば$\ker L^*$も滑らかな関数からなる有限次元空間である. 


(2). 次に$(\ker L)^\perp$上では\eqref{eq-global-Sobolev-elliptic}の$\lVert f\rVert_{L^2}$の項を除くことができることを示す. 
つまり, ある定数$C>0$が存在して
\begin{equation}
\label{eq-global-Sobolev-orth}
\lVert f\rVert_{W^{2,2}(M)}
\leq
C\lVert Lf\rVert_{L^2(M)}
\qquad
\text{for all }f\in W^{2,2}(M)\cap(\ker L)^\perp
\end{equation}
となることを示す. 
背理法. これが成り立たないとすると, 
ある$f_i\in W^{2,2}(M)\cap(\ker L)^\perp$が存在して
\[
\lVert f_i\rVert_{W^{2,2}(M)}=1,
\qquad
\lVert Lf_i\rVert_{L^2(M)}
\underset{i\to\infty}{\longrightarrow}0
\]
となる. 
Rellich--Kondrachovの定理(Theorem~\ref{thm-Rellich})から, 部分列を取り直すことで
$L^2(M)$内で
\(
f_i\to f
\)
と強収束する$f\in L^2(M)$が存在する. 
一方\eqref{eq-global-Sobolev-elliptic}を$f_i-f_j$に適用すると
\[
\begin{aligned}
\lVert f_i-f_j\rVert_{W^{2,2}(M)}
&\underalign{\eqref{eq-global-Sobolev-elliptic}}{\leq}
C\bigl(
\lVert Lf_i-Lf_j\rVert_{L^2(M)}
+
\lVert f_i-f_j\rVert_{L^2(M)}
\bigr)
\underset{i,j\to\infty}{\longrightarrow}
0.
\end{aligned}
\]
よって$f_i$は$W^{2,2}(M)$でCauchy列なので$W^{2,2}(M)$内で
\(
f_i\to f
\)
となる. 
特に
\(
Lf=0.
\)
また$f_i\perp\ker L$かつ$f_i\to f$が$L^2$収束なので
\(
f\perp\ker L.
\)
したがって
\[
f\in\ker L\cap(\ker L)^\perp=\{0\}
\]
より$f=0$である. 
しかし
\[
\lVert f\rVert_{W^{2,2}(M)}
=
\lim_{i\to\infty}
\lVert f_i\rVert_{W^{2,2}(M)}
=1
\]
となり矛盾である. 
よって\eqref{eq-global-Sobolev-orth}が成り立つ. 


(3). 
\[
R\coloneqq
\operatorname{Im}
\bigl(
L\colon W^{2,2}(M)\to L^2(M)
\bigr)
\]
とおく. 
まず$R$が$L^2(M)$の閉部分空間であることを示す. 

$Lf_i\to\rho$が$L^2(M)$で成り立つとする. 
$\ker L$への$L^2$直交射影を$P : L^2(M) \to \ker L$とする.\footnote{今$\ker L$は有限次元なので, 閉部分空間. ただそんなことを知らなくても$P$は$\ker L$の有限個の基底を用いて具体的に作れる. }
\[
\widetilde f_i\coloneqq f_i-Pf_i
\]
とおく. 
$Pf_i\in\ker L$なので
\(
\widetilde f_i\perp\ker L\)かつ
\(
L\widetilde f_i=Lf_i
\)
である. 
よって\eqref{eq-global-Sobolev-orth}から
\[
\begin{aligned}
\lVert \widetilde f_i-\widetilde f_j\rVert_{W^{2,2}(M)}
&\underalign{\eqref{eq-global-Sobolev-orth}}{\leq}
C\lVert L\widetilde f_i-L\widetilde f_j\rVert_{L^2(M)}
=
C\lVert Lf_i-Lf_j\rVert_{L^2(M)}
\underset{i,j\to\infty}{\longrightarrow}0.
\end{aligned}
\]
したがって, $W^{2,2}(M)$は完備なので, ある$f\in W^{2,2}(M)$が存在して$W^{2,2}(M)$内で
\(
\widetilde f_i\to f.
\)
今
$L\colon W^{2,2}(M)\to L^2(M)$は連続なので
\[
Lf
=
\lim_{i\to\infty}L\widetilde f_i
=
\lim_{i\to\infty}Lf_i
=
\rho.
\]
よって$\rho\in R$となり, $R$は閉である. 


次に
\begin{equation}
\label{eq-image-orth-kernel}
R^\perp=\ker L^*
\end{equation}
を示す. 
まず$h\in\ker L^*$とすると, 任意の$f\in W^{2,2}(M)$について
\[
\langle Lf,h\rangle_{L^2}
=
\langle f,L^*h\rangle_{L^2}
=0
\]
なので
\(
h\in R^\perp.
\)
よって
\(
\ker L^*\subset R^\perp.
\)
逆に$h\in R^\perp$とする. 
任意の$\varphi\in C^\infty(M)$について$L\varphi\in R$なので
\(
\langle L\varphi,h\rangle_{L^2}=0.
\)
弱解の定義より, これは$h$が
\(
L^*h=0
\)
の弱解であることを意味する. 
Theorem~\ref{thm-Sze-2.5}を$L^*$に適用すると$h$は滑らかであり, 
$h\in\ker L^*$となる. 
したがって
\(
R^\perp\subset\ker L^*.
\)
以上より\eqref{eq-image-orth-kernel}を得る. 

$R$は閉部分空間なので, Hilbert空間の基本的な性質から
\(
R=(R^\perp)^\perp.
\)
よって\eqref{eq-image-orth-kernel}から
\begin{equation}
\label{eq-image-kernel-star}
\operatorname{Im}
\bigl(
L\colon W^{2,2}(M)\to L^2(M)
\bigr)
=
(\ker L^*)^\perp.
\end{equation}


(4). 以上を用いてH\"older空間の間の全射性を示す. 
$\rho\in C^{k,\alpha}(M)$かつ
\(
\rho\perp\ker L^*
\)
とする. 
\eqref{eq-image-kernel-star}から, ある$f\in W^{2,2}(M)$が存在して
\(
Lf=\rho
\)
となる. 
さらに$f$から$\ker L$への直交射影を引くことで
\(
f\perp\ker L
\)
として良い. 
あとは$f\in C^{k+2,\alpha}(M)$となることを示す. 

以下$C>0$は$i$によらない定数で, 記号の濫用だが, 適宜定数を取り替えていく(つまり$C$は式の中で異なる可能性がある.)
$\rho$を滑らかな関数で近似する. 
有限個の座標チャート上でmollifierを用いることで, 
$\widetilde\rho_i\in C^\infty(M)$で
\[
\widetilde\rho_i\to\rho
\quad\text{in }L^2(M)
\quad \text{and} \quad 
\lVert\widetilde\rho_i\rVert_{C^{k,\alpha}(M)}
\leq
C\lVert\rho\rVert_{C^{k,\alpha}(M)}
\]
となるものを取ることができる. 

$\ker L^*$への$L^2$直交射影を$P^*$として
\(
\rho_i
\coloneqq
\widetilde\rho_i-P^*\widetilde\rho_i
\)
とおく. 
$\ker L^*$は滑らかな関数からなる有限次元空間なので
\(
\rho_i\in C^\infty(M)\),
\(
\rho_i\perp\ker L^*\),
\(
\rho_i\to\rho
\text{ in }L^2(M)
\)
であり, ある定数$C>0$があって
\begin{equation}
\label{eq-rhoi-Holder-bound}
\lVert\rho_i\rVert_{C^{k,\alpha}(M)}
\leq
C\lVert\rho\rVert_{C^{k,\alpha}(M)}
\end{equation}
となる. 
\eqref{eq-image-kernel-star}から, 各$i$について
\(
Lf_i=\rho_i\),
\(
f_i\perp\ker L
\)
となる$f_i\in W^{2,2}(M)$が存在する. 
$\rho_i$は滑らかなのでTheorem~\ref{thm-Sze-2.5}から$f_i$も滑らかである. 
$M$はコンパクトなので
\begin{equation}
\label{eq-fi-Holder-bound}
\begin{aligned}
\lVert f_i\rVert_{L^1(M)}
\leq
C\lVert f_i\rVert_{L^2(M)}
\leq
C\lVert f_i\rVert_{W^{2,2}(M)}
\underalign{\eqref{eq-global-Sobolev-orth}}{\leq}
C\lVert\rho_i\rVert_{L^2(M)}
\leq
C\lVert\rho_i\rVert_{C^{k,\alpha}(M)}
\underalign{\eqref{eq-rhoi-Holder-bound}}{\leq}
C\lVert\rho\rVert_{C^{k,\alpha}(M)}.
\end{aligned}
\end{equation}
またSchauder評価 (Theorem~\ref{thm-Schauder})を$f_i$に適用できて
\begin{equation}
\label{eq-Schauder-molifier}
\lVert f_i\rVert_{C^{k+2,\alpha}(M)}
\underalign{\text{(Thm. \ref{thm-Schauder})}}{\leq}
C\bigl(
\lVert\rho_i\rVert_{C^{k,\alpha}(M)}
+
\lVert f_i\rVert_{L^1(M)}
\bigr)
\underalign{\text{\eqref{eq-rhoi-Holder-bound} \& \eqref{eq-fi-Holder-bound}}}{\leq}
C\lVert\rho\rVert_{C^{k,\alpha}(M)}
\end{equation}
よって\eqref{eq-Schauder-molifier}とTheorem~\ref{thm-Arzela-Ascoli}から, 
任意の$0<\beta<\alpha$について部分列を取り直すことで
\[
f_i\to \widetilde f
\quad\text{in }C^{k+2,\beta}(M)
\]
となる$\widetilde f$が存在する. 
この$\widetilde f$は$f$に等しい. なぜならば, また$f_i-f\perp\ker L$なので, \eqref{eq-global-Sobolev-orth}から
\[
\lVert f_i-f\rVert_{W^{2,2}(M)}
\underalign{\eqref{eq-global-Sobolev-orth}}{\leq}
C\lVert\rho_i-\rho\rVert_{L^2(M)}
\xrightarrow{i\to\infty}0.
\]
特に$L^2$でも収束するので, 上の$W^{2,2}$収束と合わせて一意性から
\(
\widetilde f=f
\)
となるためである. よって$f\in C^{k+2,\alpha}(M)$である.  

(5). 最後に単射性を示す. 
$f\in(\ker L)^\perp\cap C^{k+2,\alpha}(M)$かつ$Lf=0$とする. 
すると
\(
f\in\ker L
\)
である一方$f\perp\ker L$なので
\[
f\in\ker L\cap(\ker L)^\perp=\{0\}.
\]
よって$f=0$となり, $L$は単射である. 

以上より
\[
L\colon
(\ker L)^\perp\cap C^{k+2,\alpha}(M)
\to
(\ker L^*)^\perp\cap C^{k,\alpha}(M)
\]
は同型である. 
\end{proof}


\begin{rem}
\label{rem-fredholm}
$k=\infty$の場合はSchauder評価 (Theorem~\ref{thm-Schauder})は不要で局所楕円型評価(Theorem~\ref{thm-local-elliptic-regularity})だけからわかる. これはSobolev空間をメインに出した次の定理からでる. 

\begin{thm}
\label{thm-fredholm-Sobolev}
$M$をコンパクトRiemann多様体,
 $L$を係数が$C^\infty$である線形2回楕円型偏微分作用素とする.
また$L^*$を$L^2$内積に関する形式的随伴とする.

この時$\ker L$および$\ker L^*$は
滑らかな関数からなる有限次元空間である.

さらに任意の$m\in\mathbb Z_{\geq 0}$について,
$L$はHilbert空間の同型
\[
L\colon
W^{m+2,2}(M)\cap(\ker L)^\perp
\overset{\sim}{\longrightarrow}
W^{m,2}(M)\cap(\ker L^*)^\perp
\]
を誘導する.
ここで$(\ker L)^\perp$および$(\ker L^*)^\perp$は
$L^2$内積に関する直交補空間とする.

さらにある定数$C_m>0$が存在して,
任意の
$f\in W^{m+2,2}(M)\cap(\ker L)^\perp$
について
\begin{equation}
\label{eq-global-Sobolev-orth-m}
\lVert f\rVert_{W^{m+2,2}(M)}
\leq
C_m\lVert Lf\rVert_{W^{m,2}(M)}
\end{equation}
が成り立つ.
特に$\rho\in W^{m,2}(M)$が
$L^2$内積で$\rho\perp\ker L^*$となるならば,
$f\in W^{m+2,2}(M)$で
$f\perp\ker L$かつ$Lf=\rho$となるものがただ一つ存在する.
\end{thm}
\begin{proof}
Theorem~\ref{thm-fredholm}の証明で変わる部分はSchauder評価を使っている(4)の全射性の部分のみである. しかも(4)において,  
$\rho\in W^{m,2}(M)$かつ
\(
\rho\perp\ker L^*
\)
とすると, $\rho\in L^2(M)$なので, 
\eqref{eq-image-kernel-star}から, ある$f\in W^{2,2}(M)$が存在して
\[
Lf=\rho \quad \text{and} \quad f\perp\ker L
\]
となる. よって示すのは$f \in W^{m+2,2}(M)$のみである. 
これは
局所楕円型正則性\eqref{eq-global-Sobolev-elliptic}よりすぐにわかる. 
$L$が連続なのは, 係数$L$の係数が$L^\infty$より
$$
\lVert Lf\rVert_{W^{m,2}(M)}
\leq
C_m
\lVert f\rVert_{W^{m+2,2}(M)}
$$
となるため. 逆写像の連続性は
局所楕円型評価(Theorem~\ref{thm-local-elliptic-regularity})
と\eqref{eq-global-Sobolev-orth}を組み合わせて
\[
\|f\|_{W^{m+2,2}}
\underalign{\text{(Thm. \ref{thm-local-elliptic-regularity})}}{\leq}
C_m\bigl(\|Lf\|_{W^{m,2}}+\|f\|_{L^2}\bigr)
\underalign{\text{\eqref{eq-global-Sobolev-orth}}}{\leq}
C'\|Lf\|_{W^{m,2}}
\]が成り立つことからわかる. 
\end{proof}

\begin{cor}
\label{cor-fredholm-smooth}
Theorem~\ref{thm-fredholm-Sobolev}と同じ仮定のもとで,
$C^\infty(M)$にseminormの族
\[
p_m(f)
\coloneqq
\lVert f\rVert_{W^{m,2}(M)},
\qquad
m=0,1,2,\ldots
\]
が生成するFr\'echet位相を入れる.
また
\(
C^\infty(M)\cap(\ker L)^\perp,
C^\infty(M)\cap(\ker L^*)^\perp
\)
には$C^\infty(M)$から誘導される相対位相を入れる.

この時$L$はFr\'echet空間の同型
\[
L\colon
C^\infty(M)\cap(\ker L)^\perp
\overset{\sim}{\longrightarrow}
C^\infty(M)\cap(\ker L^*)^\perp
\]
を誘導する.

特に$\rho\in C^\infty(M)$が
$L^2$内積で$\rho\perp\ker L^*$となるならば,
$f\in C^\infty(M)$で
$f\perp\ker L$かつ$Lf=\rho$となるものがただ一つ存在する.
\end{cor}

\begin{proof}
まず上で定めたFr\'echet位相が
通常の$C^\infty$位相と一致することを確認する.

$M$はコンパクトなので,
任意の$m\geq0$についてある定数$C_m>0$が存在して
\[
\lVert f\rVert_{W^{m,2}(M)}
\leq
C_m\lVert f\rVert_{C^m(M)}
\]
が成り立つ.
逆に$n=\dim M$とし,
$k\geq0$を固定する.
$s>k+\frac{n}{2}$となる整数$s$を取れば,
Sobolev埋め込み(Theorem~\ref{thm-Sobolev})から
ある定数$C_{k,s}>0$が存在して
\[
\lVert f\rVert_{C^k(M)}
\leq
C_{k,s}\lVert f\rVert_{W^{s,2}(M)}
\]
が成り立つ
したがってseminormの族
\(
\left\{
\lVert\bullet\rVert_{W^{m,2}(M)}
\right\}_{m\geq0}
\)
が生成する位相は通常の$C^\infty$位相と一致し,
特に$C^\infty(M)$はFr\'echet空間である.

Theorem~\ref{thm-fredholm-Sobolev}より
$\ker L$および$\ker L^*$は
滑らかな関数からなる有限次元空間である.
全単射性もTheorem~\ref{thm-fredholm-Sobolev}から明らか. 

最後にFr\'echet位相に関する連続性を確認する.
$L$は2階微分作用素なので,
任意の$m\geq0$について
\[
p_m(Lf)
=
\lVert Lf\rVert_{W^{m,2}(M)}
\leq
C_m
\lVert f\rVert_{W^{m+2,2}(M)}
=
C_mp_{m+2}(f).
\]
したがって$L$はFr\'echet位相について連続である.

一方Theorem~\ref{thm-fredholm-Sobolev}の
\eqref{eq-global-Sobolev-orth-m}から
\[
p_{m+2}(f)
=
\lVert f\rVert_{W^{m+2,2}(M)}
\leq
C_m
\lVert Lf\rVert_{W^{m,2}(M)}
=
C_mp_m(Lf)
\]
が$f\perp\ker L$について成り立つ.
また
\(
p_0(f)\leq p_2(f)
\leq C_0p_0(Lf)
\)
でもあるので$L^{-1}$もFr\'echet位相について連続である.
\end{proof}



なので, コンパクト多様体で調和積分論で使うぶんなら局所楕円型評価(Theorem~\ref{thm-local-elliptic-regularity})だけで事足りる. ただH\"older空間を使うならSchauder評価 (Theorem~\ref{thm-Schauder})が必要になる. 
\end{rem}


%%%%%%%%%%%%%%%%%%
\begin{comment}

\begin{proof}
(1). $\ker L \subset C^{k,\alpha}(M)$が有限次元となること. 
Theorem~\ref{thm-Arzela-Ascoli}より
\[
C^{k,\alpha}(M) \hookrightarrow C^{0}(M) 
\]
はコンパクト埋め込み\footnote{$C^{k,\alpha}(M)$の有界関数列は$C^0$で収束する部分列を持つ}
である. よって$\ker L \subset C^{0}(M)$
とみて有限次元を示せば良い. 
これは$ C^{0}(M)$において$\ker L$の閉単位球が(点列)コンパクト, つまり任意の点列が収束する部分列を持つことを示せば良い. 

$f_i\in \ker L$かつ$\lVert f\rVert_{C^{0}(M)} \leq 1$とする. 
$Lf_i=0$ならば, 
Schauder評価 (Theorem~\ref{thm-Schauder})により, 
\[
\lVert f_i\rVert_{C^{k+2,\alpha}(M)}
\underalign{\text{(Theorem~\ref{thm-Schauder})}}{\leq} C\bigl(\lVert \underbrace{L(f_i)}_{=0}\rVert_{C^{k,\alpha}(M)}+\lVert f_i\rVert_{L^1(M)}\bigr)
= C' \lVert f_i\rVert_{L^1(M)} 
\underalign{\text{($M$コンパクト)}}{\leq} C''\lVert f_i\rVert_{C^0(M)} 
\]
が成り立つ. 
よって$\lVert f_i\rVert_{C^{k, \alpha}(M)} \leq C$なので, $C^0$で収束する部分列を持つ, つまり
\[
f_{i_k} \to \exists f \in C^0(M)
\]
となる. よって言えた. 

(2). 次の閉値域定理とFredholm作用素の定義を思い出す.
\begin{thm}[閉値域定理]
\label{thm-closed-thm}
$T \colon H_1 \to H_2$をHilbert空間の有界線形作用素とするとき以下は同値
\[
\operatorname{Im}T \text{closed} 
\Leftrightarrow  \operatorname{Im}T^{*} \text{closed} 
\]
また上が成り立つ時, 次が成り立つ
\[
\operatorname{Im}T  = (\ker T^{*})^{\perp}
\quad
 \operatorname{Im}T^*  = (\ker T)^{\perp}
\]
\end{thm}
\begin{defn}[Fredholm作用素]
Banach空間の有界線形作用素$T \colon H_1 \to H_2$がFredholm作用素とは次を満たすこと
\begin{itemize}
\item kernelが有限次元
\item $\operatorname{Im}T$ closed
\item cokernelが有限次元
\end{itemize}
\end{defn}

次を認める
\begin{thm}
$L\colon C^{k+2,\alpha} \to  C^{k,\alpha}$はFredholm作用素である. 
\end{thm}
この証明は難しいので省略. 雰囲気をchatGPTが以下に教えてくれた. (あっているかはしらん)\footnote{ちなみにFredholm作用素のノリとして"$Tu=y$の解$u$の不定性がkernelに出る"と言っていた. わかりやすい.}
\begin{tcolorbox}[mybox]
楕円型作用素\(L\)に対して, ある擬微分作用素\(P\)が存在して,
\[
PL=I-K_1,
\qquad
LP=I-K_2
\]
となります. ここで\(K_1,K_2\)はsmoothing operatorです.

コンパクト多様体上では, smoothing operatorはコンパクト作用素になります. つまり
\[
K_1\colon C^{k+2,\alpha}\longrightarrow C^{k+2,\alpha},
\qquad
K_2\colon C^{k,\alpha}\longrightarrow C^{k,\alpha}
\]
はコンパクトです.

したがって\(L\)は
\[
\text{逆作用素を持つ, ただしコンパクト誤差を除いて}
\]
という状態になります.

一般に, Banach空間の作用素
\(
T\colon X\longrightarrow Y
\)
について
\[
ST=I-K_X,
\qquad
TS=I-K_Y
\]
を満たし, \(K_X,K_Y\)がコンパクトなら, \(T\)はFredholmです.
\end{tcolorbox}


(3). $L$がFredholm作用素なので, 閉値域定理が自由に使える. あとは全単射を定義から示す. 

$L$の全射性
$L \colon L^2(M) \to L^2(M)$とみなす. 
$\rho \in \ker L^*$について
\[
[\rho] =0 \in \Coker(L)=L^2(M)/\operatorname{Im} L \underalign{\text{(Hilbert 空間の性質)}}{=}\ker L^{*}
\]
である. 
よって
\[
\exists f \in \operatorname{Im} L \underalign{\text{(Theorem~\ref{thm-closed-thm} 閉値域定理)}}{=}(\ker L^{*})^{\perp}
\text{ s.t. } L(f)=\rho 
\]
またSchauder評価 (Theorem~\ref{thm-Schauder})により, $\rho \in C^{k, \alpha}$ならば$f \in C^{k+2, \alpha}$である. 

(4). 単射性
$L(f)=0$とすると
\[
f \in\ker L \underalign{\text{(Hilbert 空間の性質)}}{=} \Coker(L^{*})=L^2(M)/\operatorname{Im} L^{*}
\]
であるので, $f = L^{*} \rho$とかける. 
よって
\[
f = L^{*} \rho \in  \operatorname{Im}L^* \underalign{\text{(Theorem~\ref{thm-closed-thm} 閉値域定理)}}{=}(\ker L)^{\perp}
\]
である. よって, $f \in  \ker L  \cap (\ker L)^{\perp} =\{0\}$なので$f=0$である. 
\end{proof}


\end{comment}
%%%%%%%%%%%%%%%%%%%%%%%%

\subsection{調和積分論との関係}
\label{subsec-harmonic}
この辺りは\cite[Section 3]{Dem12}や小林先生の複素幾何の教科書を参照した
\begin{defn}
$X$をコンパクト複素多様体, $E$を正則ベクトル束とする. 
$E$のHermite計量を$h$, $X$のHermite計量$g$とし, Hermitian form $\omega$とする. 
\begin{itemize}
\item $E$値$(p,q)$-formの空間$A^{p,q}(E) $に内積を
\[
\langle \alpha, \beta \rangle\coloneqq 
\int_{X} 
\langle \alpha, \beta \rangle_{\omega, h}\frac{\omega^{n}}{n!}
\]
をいれる. 
ここで
\[
\langle \alpha, \beta \rangle_{\omega, h}
\coloneqq  \alpha_{I\overline{J}}^{i}\overline{\beta_{K\overline{L}}^{j}}
h_{i \bar{j}} 
g^{I \overline{K}}
g^{L \overline{J}}
\quad \text{ where }
\alpha=\alpha_{I\overline{J}}^{i}dz^I \wedge\,d\overline{z^J} \otimes e_{i}, 
\beta=\beta_{K\overline{L}}^{j}dz^K \wedge\,d\overline{z^L} \otimes e_{j}
\]

\item Hodge star operator
\[
* \colon A^{p,q}(E)\to A^{n-q,n-p}(E) 
\]
を$(n,n)$-formの次の等式が成り立つように入れる:
\[
\alpha \wedge \overline{* \beta}
=\langle \alpha, \beta \rangle_{\omega, h}  \frac{\omega^{n}}{n!}
\]
ただし左辺は$h : E \otimes \bar{E} \to \C$のcontractionを用いて$(n,n)$-formとみなしている. 
\item $\bar\partial_E^{*}\coloneqq  - * D'_E* \colon A^{p,q}(E) \to A^{p,q-1}(E) $と定義し, 
Hodge Laplacianを
\[
\square_{\mathrm{hodge}, E} 
\coloneqq  \bar\partial_E^{*} \bar\partial_E +\bar\partial_E \bar\partial_E^{*} 
\colon A^{p,q}(E) \to A^{p,q}(E) 
\]
と定義する. ただし$D'_E$を
Chern接続の次数を$(1,0)$上げるものとする. 
\end{itemize}
\end{defn}

\begin{thm}
\label{thm-harmonic-integral}
\[
\mathbf{H}^{p,q}(E)\coloneqq \{ \varphi \in A^{p,q}(E) \mid \square_{\mathrm{hodge}, E} \varphi =0\}
\]
とすると次が成り立つ. 
\begin{enumerate}
\item $\mathbf{H}^{p,q}$は有限次元
\item $\beta \in A^{p,q}(E) \cap (\ker \square_{\mathrm{hodge}, E})^\perp$ならば$\beta =\square_{\mathrm{hodge}, E} \varphi$となる$\varphi \in A^{p,q}(E)$が存在する.
\item 直交分解$A^{p,q}(E) = \mathbf{H}^{p,q}(E)\oplus \bar\partial_EA^{p,q-1}(E) \oplus \bar\partial_E^{*}A^{p,q+1}(E)$が存在する. 
\end{enumerate}
\end{thm}
\begin{proof}
(1). 
$\square_{\mathrm{hodge}, E}  \colon L^{p,q}_{2}(E) \to L^{p,q}_{2}(E)$と$L^2$完備化の空間で考える.\footnote{この書き方は実際は良くなくて, $\square_{\mathrm{hodge}, E}$は$ L^{p,q}_{2}(E)$全体で定義されてはいない. が証明では定義されている部分だけ考えるので, こう書いておく. }
$\square_{\mathrm{hodge}, E} $は楕円型なので, 
Corollary~\ref{cor-fredholm-smooth}より
\[
\square_{\mathrm{hodge}, E} \colon (\ker \square_{\mathrm{hodge}, E} )^\perp\cap A^{p,q}(E)  \to\ (\ker \square_{\mathrm{hodge}, E}^{*})^\perp\cap A^{p,q}(E)
\]
は同型で, $\ker \square_{\mathrm{hodge}, E} \eqqcolon \mathbf{H}^{p,q}(E)$は有限である. 

(2). 上に関して$\square_{\mathrm{hodge}, E} =\square_{\mathrm{hodge}, E}^{*}$であるので, 
\[
\square_{\mathrm{hodge}, E}=\square_{\mathrm{hodge}, E}^{*}
 \colon (\ker \square_{\mathrm{hodge}, E}^{*} )^\perp\cap A^{p,q}(E)  \to\ (\ker \square_{\mathrm{hodge}, E})^\perp\cap A^{p,q}(E)
\]
も同型になる. よって(全射性から)言えた. 

(3). $\supset$は明らか, $\subset$を示す. 
\[
H \colon A^{p,q}(E) \to \mathbf{H}^{p,q}(E)
\]
を直交射影とする. $\varphi \in  A^{p,q}(E)$について
\[
\varphi - H \varphi \in A^{p,q}(E) \cap (\ker \square_{\mathrm{hodge}, E})^\perp
\]
であるので, $\psi \in A^{p,q}(E) $があって
\[
\varphi - H\varphi  
\underalign{\text{(2)}}{=}
\square_{\mathrm{hodge}, E} \psi  
= \bar\partial_E^{*} \bar\partial_E \psi +\bar\partial_E \bar\partial_E^{*} \psi
\in \bar\partial_EA^{p,q-1}(E) \oplus \bar\partial_E^{*}A^{p,q+1}(E)
\]
よって言えた. 
\end{proof}

\begin{cor}
コンパクトHermite多様体$X$と正則ベクトル束$E$について, 
\[
H^{p,q}(X,E)
\coloneqq 
\frac{
\ker\bigl(
\bar\partial_E\colon A^{p,q}(X,E)\to A^{p,q+1}(X,E)
\bigr)
}{
\operatorname{Im}\bigl(
\bar\partial_E\colon A^{p,q-1}(X,E)\to A^{p,q}(X,E)
\bigr)
}
\]
とおくと$H^{p,q}(X,E)
\simeq
\mathbf H^{p,q}(X,E)$であり, 特に有限次元である. 特に大域holomorphic sectionの空間$H^0(X,E)$も有限次元である.
\end{cor}
\begin{proof}
任意の\(\alpha\in A^{p,q}(X,E)\)に対して
\[
\begin{aligned}
\langle
\square_{\mathrm{hodge},E}\alpha,\alpha
\rangle
&=
\langle
\bar\partial_E^*\bar\partial_E\alpha,\alpha
\rangle
+
\langle
\bar\partial_E\bar\partial_E^*\alpha,\alpha
\rangle
=
\lVert \bar\partial_E\alpha\rVert^2
+
\lVert \bar\partial_E^*\alpha\rVert^2.
\end{aligned}
\]
である. 特に
\(
\square_{\mathrm{hodge},E}\alpha=0
\)は\(
\bar\partial_E\alpha=0
\)かつ\(
\bar\partial_E^*\alpha=0
\)と同値である. 
さてTheorem~\ref{thm-harmonic-integral} (3)から
\[
A^{p,q}(E) = \mathbf{H}^{p,q}(E)\oplus \bar\partial_EA^{p,q-1}(E) \oplus \bar\partial_E^{*}A^{p,q+1}(E)
\]
である.

これを用いて
\(
H^{p,q}(X,E)
\simeq
\mathbf H^{p,q}(X,E)
\)を示す. 
これは
\(
\alpha\in A^{p,q}(X,E),
\bar\partial_E\alpha=0
\)について, 
\[
\alpha
=
h+\bar\partial_Eu+\bar\partial_E^*v,
\]
となる分解が取れる. \(\alpha\)は\(\bar\partial_E\)-closedより, 
\[
\langle\alpha,\bar\partial_E^*v\rangle
=
\langle\bar\partial_E\alpha,v\rangle
=
0.
\Rightarrow 
\alpha\perp\operatorname{Im}\bar\partial_E^*
\underalign{\text{(Hodge decomposition は直交分解)}}{\Rightarrow }
\bar\partial_E^*v=0.
\]よって
\(
\alpha=h+\bar\partial_Eu.
\)とかけて, したがって
\([\alpha]=[h] \in H^{p,q}(X,E)\)となり, つまりすべてのDolbeault classはharmonic representativeを持つ. 

このharmonic representativeは一意である, 実際
\(
h=\bar\partial_Eu \in\mathbf H^{p,q}(X,E)
\)なら
\[
\lVert h\rVert^2
=
\langle h,\bar\partial_Eu\rangle
=
\langle\bar\partial_E^*h,u\rangle
=0
\Rightarrow h=0
\] 
がいえることから. 
よってこの対応から
\(
H^{p,q}(X,E)
\simeq
\mathbf H^{p,q}(X,E)
\)がいえる.  
\end{proof}

\begin{thm}[Serre duality]
$X$を$n$次元コンパクト複素多様体, $E$を正則ベクトル束とする. このときpairing
\[
H^{p,q}(X,E)\times H^{n-p,n-q}(X,E^{\vee})\longrightarrow \C,
\qquad
([\alpha],[\beta])\longmapsto \int_X\alpha\wedge\beta
\]
はnon-degenerateである. 特に
\[
H^{p,q}(X,E)\simeq H^{n-p,n-q}(X,E^{\vee})^{\vee}
\]
が成り立つ. 特にDolbeault isomorphismより以下が成り立つ. 
\[
H^q(X,\Omega_X^p\otimes E)
\simeq
H^{p,q}(X,E)
\simeq
H^{n-p,n-q}(X,E^{\vee})^{\vee}
\simeq
H^{n-q}(X,\Omega_X^{n-p}\otimes E^{\vee})^{\vee}
\]
\end{thm}

\begin{proof}
$E$のHermite計量$h$を固定する. 
$h$が誘導する複素線形同型
\[
h^\flat\colon \overline{E}\to E^\vee,
\qquad
h^\flat(\overline v)(u)=h(u,v)
\]
を用いて
\[
*_h^\vee\colon A^{p,q}(E)\to A^{n-p,n-q}(E^\vee),
\qquad
*_h^\vee\alpha\coloneqq h^\flat(\overline{*\alpha})
\]
と定める. 定義から
\begin{equation}
\label{eq-harmonic-flat}
\alpha\wedge *_h^\vee\beta
=
\langle\alpha,\beta\rangle_{\omega,h}
\frac{\omega^n}{n!}
\end{equation}
が成り立つ.

$0\neq[\alpha]\in H^{p,q}(X,E)$とする. 
$H^{p,q}(X,E)\simeq\mathbf H^{p,q}(E)$によって
$\alpha$をnonzero harmonic formとして良い. 
特に
\(
\bar\partial_E\alpha=0\), \(\bar\partial_E^*\alpha=0\)
である. 任意の$\eta\in A^{p,q-1}(E)$についてStokesの定理から
\[
\begin{aligned}
\int_X\eta\wedge\bar\partial_{E^\vee}(*_h^\vee\alpha)
\underalign{\text{(Stokes)}}{=}
(-1)^{p+q}
\int_X\bar\partial_E\eta\wedge *_h^\vee\alpha
\underalign{\eqref{eq-harmonic-flat}}{=}
(-1)^{p+q}
\langle\bar\partial_E\eta,\alpha\rangle
\underalign{\text{(随伴)}}{=}
(-1)^{p+q}
\langle\eta,\bar\partial_E^*\alpha\rangle
=0.
\end{aligned}
\]
よって
\(
\bar\partial_{E^\vee}(*_h^\vee\alpha)=0
\)
であり, 
\(
[*_h^\vee\alpha]
\in
H^{n-p,n-q}(X,E^\vee)
\)
を定める.
さらに
\[
\int_X\alpha\wedge *_h^\vee\alpha
=
\int_X\lvert\alpha\rvert_{g,h}^2\frac{\omega^n}{n!}
=
\lVert\alpha\rVert_{L^2}^2
>0.
\]
よってpairingは第1成分についてnon-degenerateである.
第2成分についても同様である. 
\end{proof}



\newpage
\section{Aubin-Yauの定理とMonge--Amp\`ere方程式}
\label{Sec-MA-eq}


Szekelyhidiのlecture note \cite[Chapter 3]{Sze14}の内容. 
本来であれば$\Ric(\omega)=-\omega$となるK\"ahler--Einstein計量をメインに出すべきだが, これを次の章にとどめ, Monge--Amp\`ere方程式にメインを当てる. 
%%%%%%%%%%%%%%%%%
\begin{comment}


\subsection{The strategy}

我々の目標は次の定理を証明することである．

\begin{thm}[Aubin--Yau]
$M$が$c_1(M)<0$ ($K_M$ ample) のコンパクト K\"ahler 多様体とする. 
$\omega\in -c_1(M)$ で
$$
\Ric(\omega)
=-\omega$$
となるものが一意に存在する.

ただし$\omega \coloneqq  \sqrt{-1}g_{i \bar{j}}dz^i \wedge d \overline{z^j}$として
$$
\Ric(\omega)
\coloneqq -\frac{\sqrt{-1}}{2\pi} \partial \bar{\partial} \log(\det g)
=-\frac{\sqrt{-1}}{2\pi}\frac{\partial ^2}{\partial z^i \overline{\partial z^j}}  \log(\det g)
dz^i \wedge d \overline{z^j} \in c_1(M)
$$
である. $2\pi$で割ることに注意する. 
\end{thm}

$\omega_0 \in -c_1(M)$ を勝手に K\"ahler 計量とする． 
$\partial\bar\partial$-lemma により，
\begin{equation}
\begin{aligned}
\label{eq-ricci-1}
[\Ric(\omega_0)]=-[\omega_0 ] = c_1(M)
&\underalign{}{\Rightarrow} [\Ric(\omega_0) +\omega_0 ] =0 \in H^2(X, \R)
\\
&\underalign{\text{$\partial\bar\partial$-lemma}}{\Rightarrow }\Ric(\omega_0)=-\omega_0+\frac{\sqrt{-1}}{2\pi}\partial\bar\partial F
\quad \exists F \in C^\infty (M)
\end{aligned}
\end{equation}

$\varphi \in C^\infty(M)$について, $\omega=\omega_0+\frac{\sqrt{-1}}{2\pi}\partial\bar\partial\varphi$とする. 
\begin{equation}
\begin{aligned}
\label{eq-ricci-2}
\Ric(\omega)-\Ric(\omega_0)
\underalign{\text{定義}}{=}
-\frac{\sqrt{-1}}{2\pi} \partial \bar{\partial} \log(\det g)+\frac{\sqrt{-1}}{2\pi}\partial \bar{\partial} \log(\det g_0)
\underalign{\text{logでまとめる}}{=}
-\frac{\sqrt{-1}}{2\pi}\partial\bar\partial\log\frac{\omega^n}{\omega_0^n},
\end{aligned}
\end{equation}
以上より
\begin{align*}
\Ric(\omega)=-\omega
&\underalign{\eqref{eq-ricci-2} }{\Leftrightarrow }
\Ric(\omega_0)
-\frac{\sqrt{-1}}{2\pi}\partial\bar\partial\log\frac{\omega^n}{\omega_0^n}
=
-\omega_0-\frac{\sqrt{-1}}{2\pi}\partial\bar\partial\varphi
\\&\underalign{\eqref{eq-ricci-1} }{\Leftrightarrow }
-\frac{\sqrt{-1}}{2\pi}\partial\bar\partial\varphi
=
\frac{\sqrt{-1}}{2\pi}\partial\bar\partial F-\frac{\sqrt{-1}}{2\pi}\partial\bar\partial\log\frac{\omega^n}{\omega_0^n}
\\&\underalign{\text{logをとる}}{\Leftarrow }
\left(\omega_0+\frac{\sqrt{-1}}{2\pi}\partial\bar\partial\varphi\right)^n=e^{F+\varphi}\omega_0^n.
\end{align*}
よって次を示せば良い
\end{comment}
%%%%%%%%%%%%%%%%%%
\subsection{$c_1(X)<0$の場合のMonge--Amp\`ere方程式}

\begin{thm}
\label{thm-KE-ample}
$X$が$c_1(X)<0$のコンパクトK\"ahler多様体とする. 
K\"ahler計量$\omega_0 \in -c_1(X)$を固定し, 
\[
\Ric(\omega_0)=-\omega_0+\frac{\sqrt{-1}}{2\pi}\partial\bar\partial F
\]
となる$F \in C^\infty (X)$を固定する. 
この時(複素)Monge--Amp\`ere方程式の解$\varphi \in C^\infty (X)$で
 \begin{equation}
 \label{eq-MA-equation}
 \left(\omega_0+\frac{\sqrt{-1}}{2\pi}\partial\bar\partial\varphi\right)^n=e^{F+\varphi}\omega_0^n.
 \end{equation}
 が成り立つものがただ一つ存在する. 
\end{thm}
以下このMonge--Amp\`ere方程式の解の存在を示していく. 唯一性はすぐわかる. 
\begin{thm}
\label{thm-KE-uniqueness}
Theorem~\ref{thm-KE-ample}において, 解の唯一性が成り立つ. 
\end{thm}
\begin{proof}
\eqref{eq-MA-equation}の解を$\varphi, \psi$とし, $u=\varphi-\psi$とする. 
$u$の最大値を取る点を$p$とすると
\[
\sqrt{-1}\partial \bar\partial u \leq 0
\quad \text{ at $p$}
\]
である. 
以上より
\[
\sqrt{-1}\partial \bar\partial \varphi
\underalign{\text{(定義)}}{=}
\sqrt{-1}\partial \bar\partial \psi + 
\sqrt{-1}\partial \bar\partial u 
\underalign{\text{(最大値)}}{\leq} 
\sqrt{-1}\partial \bar\partial \psi 
\quad \text{ at $p$}
\]
となる. 
よって
\[
 \left(\omega_0+\frac{\sqrt{-1}}{2\pi}\partial\bar\partial\varphi\right)^n 
 \leq 
\left(\omega_0+\frac{\sqrt{-1}}{2\pi}\partial\bar\partial\psi \right)^n 
  \quad \text{ at $p$}
\]
である\footnote{これは正定値Hermite行列$A, B$について, $A \leq B$ならば$\det A \leq \det B$から. $A$の固有値が０以上で$B$より小さいので.}. 
よって
\[
1 \underalign{\text{(上の式)}}{\geq} \frac{ \left(\omega_0+\frac{\sqrt{-1}}{2\pi}\partial\bar\partial\varphi\right)^n }{\left(\omega_0+\frac{\sqrt{-1}}{2\pi}\partial\bar\partial\psi \right)^n }(p)
\underalign{\eqref{eq-MA-equation}}{=}
\frac{e^{\varphi}}{e^\psi}(p)
=e^u (p)
\Rightarrow 0 \geq u(p)
\]
$p$は$u$の最大値なので, $u \leq 0$が$X$全体でいえる. 
この議論を$-u$にも適応すると, $-u \leq 0$が$X$全体でいえるので, $u\equiv 0$, つまり
$\varphi=\psi$である. 
\end{proof}

この方程式\eqref{eq-MA-equation}を連続法 (continuity method) を用いて解く.
$t\in[0,1]$について
\[
\left(\omega_0+\frac{\sqrt{-1}}{2\pi}\partial\bar\partial\varphi\right)^n=e^{tF+\varphi}\omega_0^n,
\qquad
\omega_0+\frac{\sqrt{-1}}{2\pi}\partial\bar\partial\varphi\ \text{is a K\"ahler form}
\cdots (*)_t
\]
てという方程式を考える. 
\begin{enumerate}
\item $(*)_0$を解ける. これは$\varphi=0$が解になる.  
\item ある$t<1$で$(*)_t$が解をもつなら,十分小さい$\varepsilon>0$に対して$(*)_{t+\varepsilon}$も解ける.(陰関数定理).
\item ある$s\in(0,1]$に対し,すべての$t<s$で$(*)_t$を解けるなら,$(*)_s$も解ける. 
これは$t\to s$とするとき部分列に沿って極限をとれるように,H\"older空間における解の評価が必要になる.
\end{enumerate}



\subsection{解のopenness}

\begin{thm}[{\cite[Lemma 3.3]{Sze14}}]
\label{thm-open-2}
$(*)_t$がある$t<1$に対して滑らかな解をもつならば, 
十分小さい$\varepsilon>0$に対して$(*)_{t+\varepsilon}$の滑らかな解を持つ. 
\end{thm}


次の陰関数定理を思い出す.
\begin{defn}[Banach空間の$C^1$写像]
$\mathcal{X}, \mathcal{Y}, \mathcal{Z}$がBanach空間, $\Omega \subset \mathcal{Y}$が開集合とする. 
\begin{itemize}
\item $\Phi \colon \Omega \to \mathcal{Z}$がFrechet微分可能とは, 任意の$p \in \Omega$について
ある有界線形作用素$A \colon \mathcal{Y} \to \mathcal{Z}$があって
\[
\lim_{\lVert h\rVert_{\mathcal{Y}} \to 0}
\frac{\lVert \Phi(p+h) - \Phi(p) - Ah\rVert_{\mathcal{Z}}}{\lVert h\rVert_{\mathcal{Y}}}
=0
\]
を満たすこと. 
この$A$を$D\Phi(p)$とかき, $p$でのFrechet微分という. 
\item $\Phi \colon \Omega \to \mathcal{Z}$が$C^1$級であるとは
 任意の$p \in \Omega$についてFrechet微分可能かつ
 \[
D\Phi(p) \colon \Omega \to \operatorname{Hom}(\mathcal{Y}, \mathcal{Z})
\quad
p \to D\Phi(p)
 \]
 が連続となること. ただし$\operatorname{Hom}(\mathcal{Y}, \mathcal{Z})$は$\mathcal{Y} \to \mathcal{Z}$となる有界線形作用素で作用素ノルム$\lVert A\rVert\coloneqq \sup_{\lVert h\rVert_{\mathcal{Y}} =1}\lVert Ah \rVert_{\mathcal{Z}}$を入れる. 
 \item $\Phi \colon \mathcal{X} \times  \mathcal{Y} \to \mathcal{Z}$について$(x_0, y_0) \in \mathcal{X} \times \mathcal{Y}$を固定する. 
 $D_{x}\Phi(x_0, y_0)$を
 $
\Phi(\bullet, y_0) \colon \mathcal{X}\to \mathcal{Z}
 $
 に関する$x=x_0$でのFrechet微分とする. 
\end{itemize}
\end{defn}

\begin{thm}[Banach空間の陰関数定理]
\label{thm-implicit-theorem}
$\mathcal{X}, \mathcal{Y}, \mathcal{Z}$がBanach空間, $\Omega \subset \mathcal{X} \times \mathcal{Y}$を開集合とする. 
$\Phi \colon \Omega \to \mathcal{Z}$を$C^1$写像で, ある$(x_0, y_0) \in \Omega$について
\[
\Phi(x_0, y_0)=0 \quad \text{かつ} \quad
D_{x}\Phi(x_0, y_0) \colon \mathcal{X} \to \mathcal{Z} \text{がBanach 空間の同型}
\]
であるならば,  $x_0 \in U \subset \mathcal{X}$, $y_0 \in V \subset \mathcal{Y}$となる開近傍と$C^1$写像$g \colon V \to U$が存在して, 
\[
\Phi(g(y), y ) = 0 \quad \forall y \in V
\]
が成り立つ. 特に$(x_0, y_0)$のある近傍$W$があって, $(x,y) \in W$かつ$\Phi(x,y)=0$ならば$x = g(y)$である. 
\end{thm}


\begin{proof}[Proof of Theorem~\ref{thm-open-2}]
\[
\mathcal{H}^{3, \alpha}\coloneqq 
\{\varphi \in C^{3,\alpha}(X) \mid\omega_0+\frac{\sqrt{-1}}{2\pi}\partial\bar\partial\varphi >0
\}
\]
とする. "$>0$"はK\"ahler formという意味である. 
作用素を
\begin{equation}
\label{eq-mathcalF-defn}
\mathcal{F}\colon \mathcal{H}^{3, \alpha}\times[0,1]\to C^{1,\alpha}(X),
\qquad
(\varphi,t)\mapsto \log\frac{\left(\omega_0+\frac{\sqrt{-1}}{2\pi}\partial\bar\partial\varphi\right)^n}{\omega_0^n}-\varphi-tF
\end{equation}
と定める.
仮定により$\varphi_t \in \mathcal{H}^{3, \alpha}$で
\[
\log\frac{(\omega_0+\frac{\sqrt{-1}}{2\pi}\partial\bar\partial \varphi_t)^n}{\omega_0^n}-\varphi_t-tF
=
\mathcal{F}(\varphi_t,t)=0
\]
となるものがある. 
$\omega_t$を使って$X$上のH\"older normsを一つ固定する. 

\begin{claim}
\label{claim-inverse-function}
$\mathcal{X}=\mathcal{H}^{3, \alpha}, \mathcal{Y}=(0,1), \mathcal{Z}=C^{1,\alpha}(X)$, $\mathcal{F}\colon \mathcal{H}^{3, \alpha}\times[0,1]\to C^{1,\alpha}(X)$とすれば, 
点
\[(x_0, y_0)=(\varphi_t,t) \in \mathcal{H}^{3, \alpha}\times (0,1)=\mathcal{X} \times \mathcal{Y}
\]
において陰関数定理 (Theorem~\ref{thm-implicit-theorem})の仮定を満たす. 
\end{claim}
\begin{proof}[Proof of Claim~\ref{claim-inverse-function}]
点$(\varphi_t,t)$における$u$方向の$\mathcal{F}$のGateaux微分$D_u\mathcal{F}_{(\varphi_t,t)}$を計算する.
\begin{align*}
&D_u\mathcal{F}_{(\varphi_t,t)}
\\
&\underalign{\text{(定義)}}{=}
\lim_{s \to 0}\frac{1}{s} \left( \mathcal{F}_{(\varphi_t + s u ,t)} - \mathcal{F}_{(\varphi_t ,t)} \right)
\\
&\underalign{\text{\eqref{eq-mathcalF-defn}}}{=}
\lim_{s \to 0}\frac{1}{s} 
\left( 
\log\frac{(\omega_0+\frac{\sqrt{-1}}{2\pi}\partial\bar\partial \varphi_t + s\frac{\sqrt{-1}}{2\pi}\partial\bar\partial u)^n}{\omega_0^n}-(\varphi_t+su)-tF
- 
\log\frac{(\omega_0+\frac{\sqrt{-1}}{2\pi}\partial\bar\partial \varphi_t )^n}{\omega_0^n}
+\varphi_t+tF
\right)
\\
&\underalign{\text{(まとめる)}}{=}
\lim_{s \to 0}\frac{1}{s} 
\left( 
\log\frac{(\omega_0+\frac{\sqrt{-1}}{2\pi}\partial\bar\partial \varphi_t + s\frac{\sqrt{-1}}{2\pi}\partial\bar\partial u)^n}{(\omega_0+\frac{\sqrt{-1}}{2\pi}\partial\bar\partial \varphi_t )^n} 
\right)
-
\lim_{s \to 0}\frac{1}{s}(su)
\\
&\underalign{\text{(微分の計算)}}{=}
\underbrace{\frac{n}{2\pi}\frac{\sqrt{-1}\partial\bar\partial u \wedge \omega_t^{n-1}}{\omega_t^n}}_{=\frac{1}{2\pi}g_{t}^{i\bar j }\partial_i\bar\partial_{j}u = \frac{1}{2\pi}\Delta_t u}
-u
\underalign{\text{(Definition~\ref{defn-cpx-Laplace})}}{=}
\frac{1}{2\pi}\Delta_t u-u.
\end{align*}
ここで$\Delta_t$は$\omega_t$に関する複素Laplacianである. 
%ここで $\Delta_t$ は $\omega_t$ に関する Laplacian で$$\Delta_t = - dd^* -d^*d =\underalign{\text{関数に作用する場合}}{=}- dd^*$$である. 

\[
L\coloneqq D_{\bullet}\mathcal{F}_{(\varphi_t,t)}\colon C^{3,\alpha}(X)\to C^{1,\alpha}(X)
\quad
u \mapsto L(u)\coloneqq \frac{1}{2\pi}\Delta_tu-u
\]
とする. 
$\ker L =0$を示す. 
これは$\frac{1}{2\pi}\Delta_tu-u=0$ならば
\[
0 \leq 
\int_X\lvert u\rvert^2\,dV_t
\underalign{(u=\frac{1}{2\pi}\Delta_tu)}{=}
\frac{1}{2\pi}\int_Mu\Delta_tu\,dV_t
\underalign{\text{(随伴)}}{=}
=-\frac{1}{2\pi}\int_X \lvert \bar\partial u\rvert_t^2\,dV_t\leq 0
\]
より必ず$u=0$である. %\footnote{$\lvert du\rvert_t^2$とは$\omega_t$に関する$du$の内積}
$L=L^*$ (self-adjoint)なので, $L^*$も自明なkernelをもつ.
以上より$L$は楕円型なので, Theorem~\ref{thm-fredholm}から, 
\[
L\colon C^{3,\alpha}(X)=(\ker L)^\perp\cap C^{3,\alpha}
\underalign{\text{(Theorem~\ref{thm-fredholm})}}{\overset{\sim}{\to}}\ (\ker L^*)^\perp\cap C^{1,\alpha}= C^{1,\alpha}(X)
\]
は同型写像である.
よって
$
L=D_{\bullet}\mathcal{F}_{(\varphi_t,t)}
\colon C^{3,\alpha}(X)\to C^{1,\alpha}(X)
$
は同型である. 

そして$(\varphi_t,t)$の周りの点$(x, t)$について 
\[
\mathcal{H}^{3,\alpha}(X) \to \operatorname{Hom}(C^{3,\alpha}(X), C^{1,\alpha}(X))
\quad
x \mapsto (u \mapsto D_u\mathcal{F}_{(x, t)}=\frac{1}{2\pi}\Delta_x u-u)
\]
は連続である.
よって$\mathcal{F}\colon \mathcal{H}^{3, \alpha}\times[0,1]\to C^{1,\alpha}(X)$は$(\varphi_t,t)$の周りで$C^1$である\footnote{2変数$x,y$に関する関数について, $x, y$の偏微分が連続ならば全微分が存在すると同様. つまりGateaux微分が偏微分, Frechet微分が全微分に対応する. }
\end{proof}

以上より陰関数定理 (Theorem~\ref{thm-implicit-theorem})から, ある開集合$\varphi_t \in U \subset \mathcal{H}^{3, \alpha}$, $ t \in V \subset (0,1)$と$C^1$写像$g \colon V \to U$が存在して, $s \in V$ならば
\begin{equation}
\label{eq-exist-s}
0=\mathcal{F}(\underbrace{g(s)}_{\eqqcolon \varphi_s \in \mathcal{H}^{3, \alpha}}, s)
\underalign{\eqref{eq-mathcalF-defn}}{=}
\log\frac{(\omega_0+\frac{\sqrt{-1}}{2\pi}\partial\bar\partial\varphi_s)^n}{\omega_0^n}-\varphi_s-sF
\end{equation}
となる. 

$\varphi_s$が$C^\infty$であることを示す. 
\begin{itemize}
\item $\omega_0 = \sqrt{-1}g_{j\bar k}dz^j \wedge\,d\overline{z^k}$, 
\item  $(g_s)^{j\bar k}$は計量$(g_s)_{j\bar k}=g_{j\bar k}+\frac{1}{2\pi}\partial_j\partial_{\bar k}\varphi_s$の逆行列. 係数は$C^{1, \alpha}$である. 
\end{itemize}
として計算すると
\begin{align*}
\eqref{eq-exist-s}
&\underalign{\text{(localに展開)}}{\Rightarrow}
\log\det(g_{j\bar k}+\frac{1}{2\pi}\partial_j\partial_{\bar k}\varphi_s)-\log\det(g_{j\bar k})-\varphi_s-sF=0
\\
&\underalign{\text{($z^l$に関して微分)}}{\Rightarrow}
(g_s)^{j\bar k}\left(\partial_l g_{j\bar k}+\frac{1}{2\pi}\partial_l\partial_j\partial_{\bar k}\varphi_s\right)
-\partial_l\log\det(g_{j\bar k})-\partial_l\varphi_s-s\partial_l F=0
\\
&\underalign{\text{(移項)}}{\Rightarrow}
\left( \underbrace{\frac{1}{2\pi}(g_s)^{j\bar k}\partial_j\partial_{\bar k}
-1 }_{\eqqcolon E \text{係数$C^{1,\alpha}$の楕円型}}\right)
(\partial_l\varphi_s)
=
\underbrace{s\partial_l F+\partial_l\log\det(g_{j\bar k})-(g_s)^{j\bar k}\partial_l g_{j\bar k}}_{\eqqcolon h \in  C^{1,\alpha}}
\end{align*}
よって, $\partial_l\varphi_s$に対する線形楕円型方程式が$E(\partial_l\varphi_s)=h$得られる. 
以上よりSchauder評価 (Theorem~\ref{thm-Schauder})から
\[
\text{$E$と$h$の係数が$C^{1, \alpha}$}
\underalign{\text{(Thm. \ref{thm-Schauder})}}{\Rightarrow}
\partial_l\varphi_s \in C^{3, \alpha}
\underalign{\text{(定義)}}{\Rightarrow}
\varphi_s \in C^{4, \alpha}
\underalign{\text{(定義)}}{\Rightarrow}
\text{$E$と$h$の係数が$C^{2, \alpha}$}
\Rightarrow \cdots
\]
となっていき,  $\varphi_s \in C^{\infty}$が言える. \footnote{このように方程式を線形化してより良い正則性を次々に得る手法はbootstrappingと呼ばれる.}
\end{proof}


\subsection{解のclosedness}

次のPropositionを考える. 
\begin{prop}
[{\cite[Proposition 3.4]{Sze14}}]
\label{prop-uniform-estimate}
$X,\omega_0,F$のみに依存する定数$C>0$が存在して,もし$\varphi_t$がある$t\in[0,1]$について$(*)_t$を満たすなら
\[
(g_{j\bar k}+\frac{1}{2\pi}\partial_j\partial_{\bar k}\varphi_t)>C^{-1}(g_{j\bar k}),
\]
が成り立つ.ここで$g_{j\bar k}$は局所座標における$\omega_0$の成分であり,行列の不等式は差が正定値であることを意味する.さらに
\[
\lVert\varphi_t\rVert_{C^{3,\alpha}(X)}\leq C,
\]
も成り立つ.ここでH\"older normは計量$\omega_0$に関して測っている.
\end{prop}

これさえ示れば, 解のclosednessが出る. 

\begin{lem}
[{\cite[Lemma 3.5]{Sze14}}]
Proposition~\ref{prop-uniform-estimate}を仮定する.
$s\in(0,1]$とし,すべての$t<s$について方程式$(*)_t$が解けると仮定する.このとき$(*)_s$も解ける.
\end{lem}
\begin{proof}
$\lim t_i=s$を満たす$t_i<s$の数列をとる.すると
\[
\left(\omega_0+\frac{\sqrt{-1}}{2\pi}\partial\bar\partial\varphi_i\right)^n=e^{t_iF+\varphi_i}\omega_0^n
\]
を満たす関数列$\varphi_i$が得られる.
Proposition~\ref{prop-uniform-estimate}より$i$によらない$C>0$があって
\[
\lVert \varphi_i\rVert_{C^{3,\alpha}}
\leq C
\] 
となる. よってArzela-Ascoliの定理(Theorem~\ref{thm-Arzela-Ascoli})により,
$\alpha'<\alpha$となる$\alpha'$に対して, 部分列を取ることで
\[
\varphi_i
\to \exists \varphi \in C^{3,\alpha'}
\]
と収束するものが取れる. 
この収束は$C^{3,\alpha'}$収束なので,  
\[
\left(\omega_0+\frac{\sqrt{-1}}{2\pi}\partial\bar\partial\varphi\right)^n=e^{sF+\varphi}\omega_0^n
\]
を得る. 
さらにProposition~\ref{prop-uniform-estimate}より 
\[
\left(g_{j\bar k}+\frac{1}{2\pi}\partial_j\partial_{\bar k}\varphi \right)>C^{-1}(g_{j\bar k})>0
\]
であるので, 特に$\omega_0+\frac{\sqrt{-1}}{2\pi}\partial\bar\partial\varphi_i$はK\"ahler計量を定める. 
Theorem~\ref{thm-open-2}の証明と同じくして,  $\varphi$は滑らかである. 
\end{proof}

\subsection{$C^0$評価と$C^2$評価}
以下$F$を$tF$に置き換え, $\varphi_t$の$t$の添字をなくすことで次を示せば良い. 
\begin{tcolorbox}[mybox]
Monge--Amp\`ere方程式
\begin{equation}
\label{eq-MA-eq}
\left(\omega+\frac{\sqrt{-1}}{2\pi}\partial\bar\partial\varphi\right)^n=e^{F+\varphi}\omega^n,
\end{equation}
とし, $\omega = \sqrt{-1}g_{j\bar k}dz^j \wedge d \overline{z^k}$とする. H\"older normを$\omega$で測る.  
この時$X,\omega,F$のみに依存する定数$C>0$が存在して, 次の$C^2$評価が成り立つ. 
\[C^{-1}(g_{j\bar k})<(g_{j\bar k}+\frac{1}{2\pi}\partial_j\partial_{\bar k}\varphi)<C (g_{j\bar k})\]
%Proposition~\ref{prop-uniform-estimate} を示すためには$X,\omega_0,F$ のみに依存する定数 $C>0$ が存在して次が成り立つことを示せば良い. 
%\begin{itemize}
%\item (Lemma~\ref{lem-C2-estimate}) $C^2$評価が成り立つ. $C^{-1}(g_{j\bar k})<(g_{j\bar k}+\frac{1}{2\pi}\partial_j\partial_{\bar k}\varphi)<C (g_{j\bar k})$
%\item (Lemma~\ref{lem-C2-estimate}) $C^3$評価が成り立つ. $\lVert\varphi\rVert_{C^{3,\alpha}(X)}\le C$\end{itemize}
\end{tcolorbox}
これは最終的にLemma~\ref{lem-C2-estimate}で示される. 証明はかなり長い. 


\begin{lem}[{$C^0$ estimate \cite[Lemma 3.6]{Sze14}}]
\label{lem-sze-3.6}
$\varphi$が方程式\eqref{eq-MA-eq}を満たすなら, 
\[\sup_X\lvert\varphi\rvert\leq \sup_X\lvert F\rvert
\]
\end{lem}

\begin{proof}
$\varphi$の最大値を達成する点$p\in X$をとる.
Theorem~\ref{thm-KE-uniqueness}と同様に
\begin{align*}
&(g_{j\bar k}+\frac{1}{2\pi}\partial_j\partial_{\bar k}\varphi)
\underalign{\text{(最大値)}}{\leq}
g_{j\bar k}
\text{ at $p$ }
\underalign{}{\Rightarrow }
\det(g_{j\bar k}+\frac{1}{2\pi}\partial_j\partial_{\bar k}\varphi)
\leq \det(g_{j\bar k})
\text{ at $p$ }
\\
&\underalign{}{\Rightarrow }
e^{F+\varphi}=
\frac{\left(\omega+\frac{\sqrt{-1}}{2\pi}\partial\bar\partial\varphi\right)^n}{\omega^n}
\underalign{\text{(上の式)}}{\leq} 1
\text{ at $p$ } \\
&\underalign{}{\Rightarrow }
F(p) + \varphi(p) \leq 0
\underalign{}{\Rightarrow }
-F(p) \geq \varphi(p) \underalign{\text{最大値}}{\geq}\varphi(x) 
\quad \forall x\in X
\end{align*}
同様にして, $\varphi$の最小値を達成する点$q\in X$において
$-F(q) \leq \varphi(q) $であるので, 二つを合わせて, $\sup_X\lvert\varphi\rvert\leq \sup_X\lvert F\rvert$である. 
\end{proof}


normal coordinateに関して復習する.
\begin{prop}[Normal coordinates]
\label{prop-normal-coordinate}
$g$がK\"ahler計量とすると, 
任意の点$p\in X$について, その近傍上の正則座標$z^1,\ldots,z^n$であって, 
点$p$における$g$の成分が
\[
g_{j\bar k}(p)=\delta_{jk}
\quad\text{and}\quad
\frac{\partial}{\partial z^i}g_{j\bar k}(p)
=
\frac{\partial}{\partial \bar z^i}g_{j\bar k}(p)=0,
%\tag{6}
\]
となるようにできる. 
%ここで $\delta_{jk}$ は恒等行列、すなわち $j\neq k$ なら $\delta_{jk}=0$、$j=k$ なら $\delta_{jk}=1$ である。
\end{prop}

\begin{proof}
K\"ahler formが
\[
\omega=\sqrt{-1}\sum_{j,k}\bigl(\delta_{jk}+O(\lvert z\rvert^2)\bigr)\,dz^j\wedge\,d\bar z^k,
%\tag{7}
\]
とできることを示せば良い. ($O(\lvert z\rvert^2)$は,$z^i,\bar z^i$に関して少なくとも2次の項とする)

まず線形変換することで, 
\[
\omega=\sqrt{-1}\sum_{j,k}\left(\delta_{jk}+\sum_l\bigl(a_{j\bar{k}l}w^l+a_{jk\bar{l}}\bar{w}^l\bigr)+O(\lvert w\rvert^2)\right)\,dw^j\wedge\,d\bar{w}^k.
%\tag{8}
\]
となる座標$w^i$がある. 
次に$b_{ijk}=b_{ikj}$を満たす係数$b_{ijk}$, $w^i$から$z^i$への座標変換を
\[
w^i\coloneqq z^i-\frac12\sum_{j,k}b_{ijk}z^jz^k,
\]
で定義すると, 
\[
dw^i=dz^i-\sum_{j,k}b_{ijk}z^jdz^k,
\]
であるので
\begin{align*}
\omega
&=\sqrt{-1}\sum_{j,k}\left(\delta_{jk}+\sum_l\bigl(a_{j\bar{k}l}z^l+a_{jk\bar{l}}\bar{z}^l-b_{klj}z^l-\bar{b}_{jlk}\bar{z}^l\bigr)+O(\lvert z\rvert^2)\right)\,dz^j\wedge\,d\bar{z}^k
\\
&=
\sqrt{-1}\sum_{j,k}\left(\delta_{jk}+
\sum_l\bigl(a_{j\bar{k}l}-b_{klj}\bigr)z^l
+ \sum_l\bigl(a_{jk\bar{l}}-\bar{b}_{jlk}\bigr)\bar{z}^l
+O(\lvert z\rvert^2)\right)\,dz^j\wedge\,d\bar{z}^k
\end{align*}
となる. $\omega$がK\"ahlerであれば,定義から$a_{j\bar{k}l}=a_{l\bar{k}j}$なので, 
$b_{klj}\coloneqq a_{j\bar{k}l}$と定れば良い. 
\end{proof}
\begin{rem}
Riemannian geometryでは,与えられた点で計量の一次微分が消えるような\underline{$C^\infty$級}な
normal coordinatesを常に取ることができる. Hermitian計量でも同様である. 

Proposition~\ref{prop-normal-coordinate}によって, 計量がK\"ahlerであれば, 計量の一次微分がある点で消えるような \underline{正則}座標系さえ見つけられる, ということである.逆にそのような正則なnormal coordinatesが存在すれば$d\omega=0$となり, 計量はK\"ahlerである.

\end{rem}

以下次のnotationを使う.
\begin{tcolorbox}[mybox]
\begin{itemize}
\item $g'_{j\bar k}\coloneqq g_{j\bar k}+\frac{1}{2\pi}\partial_j\partial_{\bar k}\varphi$
\item $\operatorname{tr}_g g' \coloneqq  g^{j\bar k}g'_{j\bar k}$. $\operatorname{tr}_{g'}g\coloneqq g'^{j\bar k}g_{j\bar k}$
\end{itemize}
\end{tcolorbox}
このnotationの上で次を得る. 
\[
g^{j\bar k}g'_{j\bar k}
\underalign{\text{定義}}{=}
g^{j\bar k}g_{j\bar k}+\frac{1}{2\pi}\underbrace{g^{j\bar k}\partial_j\partial_{\bar k}\varphi}_{\eqqcolon \Delta \varphi}
=n+\frac{1}{2\pi}\Delta\varphi.
\]
また$\Delta'$を計量$g'$に関するLaplacianとする.

\begin{lem}
[{\cite[Lemma 3.7]{Sze14}}]
\label{lem-logdelta-estimate}
$X$と$g$のみに依存する定数$B$が存在して
\[
\Delta'\log \operatorname{tr}_g g'\geq -B\,\operatorname{tr}_{g'}g-\frac{g^{j\bar k}R'_{j\bar k}}{\operatorname{tr}_g g'},
\]
が成り立つ. ここで$R'_{j\bar k}$は$g'$のRicci曲率である.
\end{lem}


\begin{proof}
$p \in X$を固定する. 
normal coordinateを用いる., 点$p$において
\begin{equation}
\label{eq-normal-cordinate-1}
g=\sum_{i=1}^{n}dz^i \wedge\,d\overline{z^i}, \quad 
\frac{\partial}{\partial z^i}g_{j\bar k}(p)
=
\frac{\partial}{\partial \bar z^i}g_{j\bar k}(p)=0,
\end{equation}
また任意のHermitian matrixはunitary transformationにより対角化できるので, $g'$は$p$で対角であると仮定よく, 下の形を得る. 
\begin{equation}
\label{eq-normal-cordinate-2}
g'=\sum_{i=1}^{n} g'_{i\bar{i}}dz^i \wedge\,d\overline{z^i}, \quad 
\operatorname{tr}_g g'=\sum_i g'_{i\bar i},
\quad
\operatorname{tr}_{g'}g=\sum_j g'^{j\bar j}=\sum_j\frac{1}{g'_{j\bar j}}
\text{ at $p$ }
\end{equation}
と仮定して良い. 
点$p$での以下の式が成り立つ. 
\begin{equation}
\begin{aligned}
\label{eq-delta-tr}
\Delta'\operatorname{tr}_g g'
&\underalign{\text{($\Delta'$の定義)}}{=}
\underbrace{g'^{p\bar q}\partial_p\partial_{\bar q}}_{\eqqcolon \Delta'}
\underbrace{(g^{j\bar k}g'_{j\bar k})}_{\eqqcolon \operatorname{tr}_g g'}
\\
&\underalign{\text{\eqref{eq-normal-cordinate-1}}}{=}
g'^{p\bar q}(\partial_p\partial_{\bar q}g^{j\bar k})g'_{j\bar k}
+ g'^{p\bar q}g^{j\bar k} \times
\underbrace{(\partial_p\partial_{\bar q}g'_{j\bar k})}_{=-R'_{j\bar k p\bar q}+g'^{a\bar b}(\partial_j g'_{p\bar b})(\partial_{\bar k}g'_{a\bar q}) \text{ by $R'_{i\bar{j}k\bar{l}}$の定義 in Sum. \ref{suma-brief-diff}}}
\\
&\underalign{}{=}
 \underbrace{g'^{p\bar q}(\partial_p\partial_{\bar q}g^{j\bar k})g'_{j\bar k}}_{\text{Term(1)}}
- \underbrace{g'^{p\bar q}g^{j\bar k}R'_{j\bar k p\bar q}}_{=g^{j\bar k}R'_{j\bar k} 
\text{ by $R'_{i\bar{j}}$の定義 in Sum. \ref{suma-brief-diff}}}
+ \underbrace{g'^{p\bar q}g^{j\bar k}g'^{a\bar b}(\partial_j g'_{p\bar b})(\partial_{\bar k}g'_{a\bar q})}_{=\sum_{p,j,a} g'^{p\bar p}g'^{a\bar a}\,\lvert\partial_j g'_{p\bar a}\rvert^2 \text{ by \eqref{eq-normal-cordinate-1} \& \eqref{eq-normal-cordinate-2}}}
\\
&\underalign{\text{(下のTerm(1)の評価 \eqref{eq-term1-ev})}}{\geq}
 -B(\operatorname{tr}_{g'}g)(\operatorname{tr}_g g')-g^{j\bar k}R'_{j\bar k}
+\sum_{p,j,a} g'^{p\bar p}g'^{a\bar a}\,\lvert\partial_j g'_{p\bar a}\rvert^2.
\end{aligned}
\end{equation}
ここで$B\coloneqq \inf_{x \in X} - R_{i \bar{j}k \bar{l}}$とする.  Term(1)の評価は以下のとおり:  
まず次が成り立つ.
\begin{equation}
\begin{aligned}
\label{eq-BSC-inf}
-\partial_p\partial_{\bar p}g^{j\bar j}
&\underalign{\text{\eqref{eq-normal-cordinate-1} }}{=}
 -\partial_i \partial_{\bar{j}} g_{k \bar{l}}
 \underalign{\text{($R'_{i\bar{j}k\bar{l}}$の定義 in Sum. \ref{suma-brief-diff})}}{=}
= - R_{i \bar{j}k \bar{l}} +g^{p\bar{q}}(\partial_i g_{k \bar{q}})(\partial_{\bar{j}}g_{p\bar{l}})
\\
&\underalign{\text{\eqref{eq-normal-cordinate-1} }}{=}
- R_{i \bar{j}k \bar{l}} 
\underalign{}{\geq}
 \inf_{x \in X} - R_{i \bar{j}k \bar{l}} \eqqcolon B
\end{aligned}
\end{equation}
以上より$g'$が対角であることを用いると
\begin{equation}
\label{eq-term1-ev}
\underbrace{g'^{p\bar q}(\partial_p\partial_{\bar q}g^{j\bar k})g'_{j\bar k}}_{\text{Term(1)}}
\underalign{\text{\eqref{eq-normal-cordinate-1} \& \eqref{eq-normal-cordinate-2} }}{=}
 \sum_{p,j} g'^{p\bar p}g'_{j\bar j}\,(\partial_p\partial_{\bar p}g^{j\bar j})
 \underalign{\eqref{eq-BSC-inf}}{\geq}
 -B\sum_{p,j} g'^{p\bar p}g'_{j\bar j}
\underalign{\text{\eqref{eq-normal-cordinate-2} }}{=}
 -B(\operatorname{tr}_{g'}g)(\operatorname{tr}_g g'),
\end{equation}
これを用いて$\Delta'\log \operatorname{tr}_g g'$を評価する. 
まず次の公式を認める. 
\begin{tcolorbox}[mybox]
\(u>0\)ならば
\[
\partial_p\partial_{\bar q}\log u
=
\partial_p\left(\frac{\partial_{\bar q}u}{u}\right)
=
\frac{\partial_p\partial_{\bar q}u}{u}
-
\frac{(\partial_p u)(\partial_{\bar q}u)}{u^2}
\]
特に\(
\Delta'
=
g'^{p\bar q}\partial_p\partial_{\bar q}
\)
について, 
\begin{equation}
\label{eq-logu-2nddiff}
\Delta'\log u
=
g'^{p\bar q}\partial_p\partial_{\bar q}\log u
=
g'^{p\bar q}
\left(
\frac{\partial_p\partial_{\bar q}u}{u}
-
\frac{(\partial_p u)(\partial_{\bar q}u)}{u^2}
\right)
=
\frac{\Delta' u}{u}
-
\frac{g'^{p\bar q}(\partial_p u)(\partial_{\bar q}u)}{u^2}
\end{equation}
\end{tcolorbox}

logを組み込んで, $u=\operatorname{tr}_g g'$とすれば
\begin{equation}
\begin{aligned}
\label{eq-logtr-lap}
&\Delta'\log \operatorname{tr}_g g'
\underalign{\eqref{eq-logu-2nddiff}}{=}
 \frac{\Delta'\operatorname{tr}_g g'}{\operatorname{tr}_g g'}
- \frac{g'^{p\bar q}(\partial_p\operatorname{tr}_g g')(\partial_{\bar q}\operatorname{tr}_g g')}{(\operatorname{tr}_g g')^2}
\\
&\underalign{\eqref{eq-delta-tr} \text{ \& } \eqref{eq-normal-cordinate-2}}{\geq}
\frac{1}{\operatorname{tr}_g g'}
\left( 
 -B(\operatorname{tr}_{g'}g)(\operatorname{tr}_g g')-g^{j\bar k}R'_{j\bar k}
+\sum_{p,j,a} g'^{p\bar p}g'^{a\bar a}\,\lvert\partial_j g'_{p\bar a}\rvert^2.
\right)
-\frac{1}{(\operatorname{tr}_g g')^2}\sum_{p,a,b} g'^{p\bar p}(\partial_p g'_{a\bar a})(\partial_{\bar p}g'_{b\bar b})
\\
&\underalign{\text{}}{=}
 -B\,\operatorname{tr}_{g'}g-\frac{g^{j\bar k}R'_{j\bar k}}{\operatorname{tr}_g g'}
+\frac{1}{\operatorname{tr}_g g'}\sum_{p,j,a} g'^{p\bar p}g'^{a\bar a}\,\lvert\partial_j g'_{p\bar a}\rvert^2
-\frac{1}{(\operatorname{tr}_g g')^2}\sum_{p,a,b} g'^{p\bar p}(\partial_p g'_{a\bar a})(\partial_{\bar p}g'_{b\bar b}).
%\tag{30}
\end{aligned}
\end{equation}

ここでCauchy--Schwarz inequality 
\[
\left(\sum_{p}x_py_p \right)^2 \leq 
\left(\sum_{p}x_{p}^2\right)\cdot
\left(\sum_{q}y_{q}^2 \right)
\]を2回用いると
\begin{equation}
\label{eq-trace-CS}
\begin{aligned}
&\sum_{p,a,b} g'^{p\bar p}(\partial_p g'_{a\bar a})(\partial_{\bar p}g'_{b\bar b})
\underalign{\text{(等式変形)}}{=}
\sum_{a,b}\sum_p 
\underbrace{\sqrt{g'^{p\bar p}}(\partial_p g'_{a\bar a})}_{\text{($x_p$とみる)}}
\underbrace{\sqrt{g'^{p\bar p}}(\partial_{\bar p}g'_{b\bar b})}_{\text{($y_p$とみる)}}
\\
&\underalign{\text{(Cauchy--Schwarz)}}{\leq}
 \sum_{a,b}\Bigl(\sum_p g'^{p\bar p}\lvert\partial_p g'_{a\bar a}\rvert^2\Bigr)^{1/2}
\Bigl(\sum_q g'^{q\bar q}\lvert\partial_q g'_{b\bar b}\rvert^2\Bigr)^{1/2}
\underalign{\text{(まとめる)}}{=}
\Bigl(\sum_a \Bigl(\sum_p g'^{p\bar p}\lvert\partial_p g'_{a\bar a}\rvert^2\Bigr)^{1/2}\Bigr)^2
\\
&\underalign{\text{(等式変形)}}{=}
\Bigl(\sum_a 
\underbrace{\sqrt{g'_{a\bar a}}}_{\text{($x_a$とみる)}}
\underbrace{\Bigl(\sum_p g'^{p\bar p}g'^{a\bar a}\lvert\partial_p g'_{a\bar a}\rvert^2\Bigr)^{1/2}
}_{\text{($y_a$とみる)}}\,
\Bigr)^2
\underalign{\text{(Cauchy--Schwarz)}}{\leq}
 \underbrace{\Bigl(\sum_a g'_{a\bar a}\Bigr)}_{=\operatorname{tr}_g g'}
 \Bigl(\sum_{b,p} g'^{p\bar p}g'^{b\bar b}\lvert\partial_p g'_{b\bar b}\rvert^2\Bigr).
\end{aligned}
\end{equation}
したがって
\begin{equation}
\label{eq-trace-CS-2}
\frac{1}{(\operatorname{tr}_g g')^2}\sum_{p,a,b} g'^{p\bar p}(\partial_p g'_{a\bar a})(\partial_{\bar p}g'_{b\bar b})
\underalign{\eqref{eq-trace-CS}}{\leq}
\frac{1}{\operatorname{tr}_g g'}\sum_{a,p} g'^{p\bar p}g'^{a\bar a}\lvert\partial_p g'_{a\bar a}\rvert^2
\underalign{\text{(非負の項を足す)}}{\leq}
\frac{1}{\operatorname{tr}_g g'}\sum_{a,j,p} g'^{p\bar p}g'^{a\bar a}\lvert\partial_p g'_{j\bar a}\rvert^2,
\end{equation}
以上より
\begin{equation*}
\begin{aligned}
&\Delta'\log \operatorname{tr}_g g'
\\
&\underalign{\eqref{eq-logtr-lap}}{\geq}
-B\,\operatorname{tr}_{g'}g-\frac{g^{j\bar k}R'_{j\bar k}}{\operatorname{tr}_g g'}
+\frac{1}{\operatorname{tr}_g g'}\sum_{p,j,a} g'^{p\bar p}g'^{a\bar a}\,
\underbrace{\lvert\partial_j g'_{p\bar a}\rvert^2}_{=\lvert\partial_p g'_{j\bar a}\rvert^2 \text{ by K\"ahler}}
-\frac{1}{(\operatorname{tr}_g g')^2}\sum_{p,a,b} g'^{p\bar p}(\partial_p g'_{a\bar a})(\partial_{\bar p}g'_{b\bar b}).
\\
&\underalign{\eqref{eq-trace-CS-2}}{\geq}
-B\,\operatorname{tr}_{g'}g-\frac{g^{j\bar k}R'_{j\bar k}}{\operatorname{tr}_g g'}
+
\underbrace{
\frac{1}{\operatorname{tr}_g g'}\sum_{p,j,a} g'^{p\bar p}g'^{a\bar a}\,
\lvert\partial_p g'_{j\bar a}\rvert^2
-
\frac{1}{\operatorname{tr}_g g'}\sum_{a,j,p} g'^{p\bar p}g'^{a\bar a}\lvert\partial_p g'_{j\bar a}\rvert^2
}_{=0}
\\
&\underalign{}{=}
-B\,\operatorname{tr}_{g'}g-\frac{g^{j\bar k}R'_{j\bar k}}{\operatorname{tr}_g g'}
\end{aligned}
\end{equation*}
\end{proof}

\begin{lem}
[{$C^2$ estimate \cite[Lemma 3.8]{Sze14}}]
\label{lem-C2-estimate}
$X,\omega,\sup_X\lvert F\rvert$および$\Delta F$の下界のみに依存する定数$C$が存在して,
\eqref{eq-MA-eq}の解$\varphi$は次を満たす. 
\[
C^{-1}(g_{j\bar k})<\left(g_{j\bar k}+\frac{1}{2\pi}\partial_j\partial_{\bar k}\varphi\right)<C(g_{j\bar k})
\]
\end{lem}

\begin{proof}

前と同様に$g'_{j\bar k}=g_{j\bar k}+\frac{1}{2\pi}\partial_j\partial_{\bar k}\varphi$と書く.
まず$(g_{j\bar k}+\frac{1}{2\pi}\partial_j\partial_{\bar k}\varphi)<C(g_{j\bar k})$を示す. 
これには, ${g'_{i\bar i}}$を$g$に関する$g'$の固有値として, 点$x \in X$によらずに${g'_{i\bar i}}$が定数で抑えられることを示せば良い. 
方程式\eqref{eq-MA-eq}から
\begin{align*}
\eqref{eq-MA-eq}
&\underalign{}{\Rightarrow }
\left(\omega+\frac{\sqrt{-1}}{2\pi}\partial\bar\partial\varphi\right)^n=e^{F+\varphi}\omega^n
\\
&\underalign{\text{($-\partial \bar\partial \log(\bullet)$をとる)}}{\Rightarrow }
-R'_{j\bar k}=\partial_j\partial_{\bar k}F+\partial_j\partial_{\bar k}\varphi-R_{j\bar k}
=\partial_j\partial_{\bar k}F+2\pi\underbrace{(g'_{j\bar k}-g_{j\bar k})}_{=\frac{1}{2\pi}\partial_j\partial_{\bar k}\varphi}
-R_{j\bar k}.
\end{align*}
となる. Lemma~\ref{lem-logdelta-estimate}
を用いると, 
$X$と$g$のみに依存する定数$B$が存在して
\begin{equation}
\label{eq-logtrace}
\Delta'\log \operatorname{tr}_g g'
\underalign{\text{(Lem. \ref{lem-logdelta-estimate})}}{\geq}
-B\,\operatorname{tr}_{g'}g-\frac{g^{j\bar k}R'_{j\bar k}}{\operatorname{tr}_g g'}
=
-B\,\operatorname{tr}_{g'}g+\frac{\Delta F+2\pi\operatorname{tr}_g g'-2\pi n-R}{\operatorname{tr}_g g'},
\end{equation}
を得る.ここで$R\coloneqq g^{j\bar k}R_{j\bar k}$は$g$のscalar curvatureである.
Cauchy--Schwarz inequalityから
\begin{equation}
\label{eq-trace-CS-3}
(\operatorname{tr}_g g')(\operatorname{tr}_{g'}g)
\underalign{\text{\eqref{eq-normal-cordinate-2} }}{=}
\left(\sum_i g'_{i\bar i}\right)
\cdot
\left(\sum_j\frac{1}{g'_{j\bar j}}\right)
\underalign{\text{(Cauchy--Schwarz)}}{\geq}
\left(\sum_i g'_{i\bar i}\cdot \frac{1}{g'_{i\bar i}}\right)^2
=
 n^2
\end{equation}
が成り立つ. 
そこで$-C'\coloneqq \min\{ -1, \Delta F-2\pi n-R\}$とすれば, $C'>0$で$C'$は$\Delta F$や$\omega$にしかよらず, 
\begin{equation}
\label{eq-logdelta-bc}
\begin{aligned}
\Delta'\log \operatorname{tr}_g g'
&\underalign{\text{\eqref{eq-logtrace}}}{\geq}
-B\,\operatorname{tr}_{g'}g+
\underbrace{\frac{\Delta F-2\pi n-R}{\operatorname{tr}_g g'}}
_{\geq -\frac{ C'}{n^2}\operatorname{tr}_{g'}g \text{ by \eqref{eq-trace-CS-3}}} 
+
2\pi
\\
&\underalign{}{\geq}
-B\,\operatorname{tr}_{g'}g
 -\underbrace{\frac{ C'}{n^2}}_{\eqqcolon C} \operatorname{tr}_{g'}g
\end{aligned}
\end{equation}
となる. 定義から
\[
\Delta'\varphi
\underalign{\text{($\Delta'$の定義)}}{=}g'^{j\bar k}\partial_j\partial_{\bar k}\varphi
\underalign{\text{($\varphi$の定義)}}{=}
2\pi g'^{j\bar k}(g'_{j\bar k}-g_{j\bar k})
\underalign{\text{($\operatorname{tr}$の定義)}}{=}
2\pi(n-\operatorname{tr}_{g'}g)
\]
であるので, 
 $A \coloneqq\frac{B+C+1}{2\pi}$とおけば
\begin{equation}
\label{eq-delta-logA}
\Delta'(\log \operatorname{tr}_g g'-A\varphi)
\underalign{\text{(\eqref{eq-logdelta-bc} \& 上の式)}}{\geq}
-B\,\operatorname{tr}_{g'}g-C\,\operatorname{tr}_{g'}g
-2\pi A(n-\operatorname{tr}_{g'}g)
\underalign{\text{($A$の定義)}}{=}
\operatorname{tr}_{g'}g-2\pi An
\end{equation}

さて$\log \operatorname{tr}_g g'-A\varphi$が最大値を達成する点$p\in X$をとる.
このとき
\begin{equation}
\label{eq-trace-estimate-gii}
0
\underalign{\text{(最大値)}}{\geq}
\ \Delta'(\log \operatorname{tr}_g g'-A\varphi)(p)
\underalign{\eqref{eq-delta-logA}}{\geq}\operatorname{tr}_{g'}g(p)-2\pi An
\Rightarrow
\underbrace{\operatorname{tr}_{g'}g(p)}_{\sum_{i}\frac{1}{g'_{i\bar i}}=\sum_{i}g'^{i\bar i}}\leq 2\pi An.
\end{equation}
また\eqref{eq-MA-eq}より点$p$において,
\begin{equation}
\label{eq-estimate-g_ii}
\eqref{eq-MA-eq}
\underalign{}{\Rightarrow }
\underbrace{\left(\omega+\frac{\sqrt{-1}}{2\pi}\partial\bar\partial\varphi\right)^n(p)
}_{=\prod_{i}g'_{i\bar i}}
=e^{F(p)+\varphi(p)}
\underbrace{\omega^n}_{=\prod_{i}g_{i\bar i}=1}
\underalign{}{\Rightarrow }
\prod_{i=1}^n g'_{i\bar i}=e^{F(p)+\varphi(p)}
\underalign{\text{(Lem. \ref{lem-sze-3.6})}}{\leq}
e^{2\sup\lvert F\rvert}\eqqcolon C_1
\end{equation}
よって$0 < g'_{1\bar{1}} \leq \cdots \leq g'_{n\bar{n}} $
とすれば
\begin{equation}
\label{eq_traceg'g-estimate}
\begin{aligned}
&2\pi An \underalign{\eqref{eq-trace-estimate-gii}}{\geq}
\underbrace{ \sum_{i}\frac{1}{g'_{i\bar i}}}_{=\operatorname{tr}_{g'}g(p)} \geq \frac{1}{g'_{1\bar 1}}
\underalign{}{\Rightarrow }
g'_{1\bar 1} \geq \frac{1}{2\pi An}
\\
&\underalign{}{\Rightarrow }
C_1 \underalign{\eqref{eq-estimate-g_ii}}{\geq}
\prod_{i=1}^n g'_{i\bar i} 
\geq g'_{n\bar n} (g'_{1\bar 1})^{n-1}
\geq   g'_{n\bar n}  \left(\frac{1}{2\pi An}\right)^{n-1}
\\
&\underalign{}{\Rightarrow }
\underbrace{C_1 \cdot  \left(2\pi An\right)^{n-1}}_{\eqqcolon C_2}
\geq   g'_{n\bar n}  
\geq \cdots  \geq g'_{1\bar{1}} >0
\\
&\underalign{}{\Rightarrow }
\operatorname{tr}_g g'(p)
\underalign{\text{(定義)}}{=}
\sum_{i=1}^n g'_{i\bar i}
\underalign{}{\leq}
 n g'_{n\bar{n}} 
\underalign{}{\leq}nC_2.
\end{aligned}
\end{equation}
よって
任意の$x \in X$について
\[
\log \operatorname{tr}_g g'(x)-A\varphi(x)
\underalign{\text{($p$が最大値)}}{\leq}
\log \operatorname{tr}_g g'(p)-A\varphi(p)
\underalign{\text{\eqref{eq_traceg'g-estimate}}}{\leq} \log(nC_2)-A\varphi(p)
\underalign{\text{(Lem. \ref{lem-sze-3.6})}}{\leq}
\underbrace{\log(nC_2)+A\sup \lvert F\rvert}_{\eqqcolon C_3}
\]
が成り立つ. 
これより
\begin{equation}
\label{eq-gnn-bound}
\sup_{x \in X} \log \operatorname{tr}_g g'(x)
\underalign{\text{(上の式)}}{\leq}
C_3 + A \sup \lvert\varphi\rvert
\underalign{\text{(Lem. \ref{lem-sze-3.6})}}{\leq}
\underbrace{C + A\sup \lvert F\rvert}_{\eqqcolon C_4}
\underalign{}{\Rightarrow}
(g'_{i\bar i})_{x} \leq C_4
\end{equation}
(ただし${g'_{i\bar i}}_{x}$は点$x$での$g$に関する$g'$の固有値とする.)
以上より$(g_{j\bar k}+\frac{1}{2\pi}\partial_j\partial_{\bar k}\varphi)<C(g_{j\bar k})$が言えた. 

逆側に関しては任意の$x \in X$について, normal coordinateをとると, $0 < (g'_{1\bar{1}})_x\leq \cdots \leq (g'_{n\bar{n}})_x$とすれば, 
\eqref{eq-estimate-g_ii}と同じ議論により, 
\[
\begin{aligned}
&\prod_{i=1}^n (g'_{i\bar i})_x
\underalign{\eqref{eq-estimate-g_ii}}{=}e^{F(x)+\varphi(x)}
\underalign{\text{(Lem. \ref{lem-sze-3.6})}}{\geq}
\underbrace{e^{\inf F - \sup |F|}}_{\eqqcolon C_5}
\\
&\underalign{}{\Rightarrow}
C_5 \leq \prod_{i=1}^n (g'_{i\bar i})_x \leq {g'_{1\bar{1}} }_x \left( {g'_{n\bar{n}}}_x\right)^{n-1}
\underalign{\eqref{eq-gnn-bound}}{\leq} C_{4}^{n-1} (g'_{1\bar{1}})_x  
\underalign{}{\Rightarrow}
\frac{C_5}{C_{4}^{n-1}} \leq  {g'_{1\bar{1}} }_x  
\end{aligned}
\]
となる. よって下からの評価もえて, $C^{-1}(g_{j\bar k})<(g_{j\bar k}+\frac{1}{2\pi}\partial_j\partial_{\bar k}\varphi)$となる. 
\end{proof}

\subsection{$C^{3}$評価}

%この節では、方程式 (28) を満たす $\varphi$ の三階微分に対する評価を導く。計算は Phong--Sesum--Sturm \cite{39} に従うが、これは \cite{55}, \cite{5} にある元の証明（あるいはその放物型の類似）をより簡潔にしたものである。より一般的な手法を用いて、前節で得た $\partial_{j}\partial_{\bar k}\varphi$ の評価から $C^{2,\alpha}$ 評価を得ることも可能である。すなわち Evans--Krylov の定理の複素版を用いる方法である（このアプローチについては \cite{7} または \cite{44} を参照）。

次のnotationを使う.
\begin{tcolorbox}[mybox]
\begin{itemize}
\item 固定された計量を$\hat g_{j\bar k}$とかく. (つまり$\hat g_{j\bar k}$が$g_0$に対応すると思って良い)
\item Monge--Amp\`ere方程式
の解の計量を
\(
g_{j\bar k}=\hat g_{j\bar k}+\frac{1}{2\pi}\partial_{j}\partial_{\bar k}\varphi
\)とする. 定義から次を満たす. 
\begin{equation*}
\left(\hat{\omega}+\frac{\sqrt{-1}}{2\pi}\partial\bar\partial\varphi\right)^n=e^{F+\varphi}\hat{\omega}^n,
\end{equation*}
\item 上の式からRicci曲率に関する方程式
\begin{equation}
\label{eq-ric-Tij}
-R_{j\bar k}
=\partial_j\partial_{\bar k}F+2\pi(g_{j\bar k}-\hat{g}_{j\bar k})
-\hat{R}_{j\bar k}
\Rightarrow 
R_{j\bar k}=-2\pi g_{j\bar k}+
\underbrace{2\pi\hat{g}_{j\bar k}-\partial_j\partial_{\bar k}F +\hat{R}_{j\bar k}}_{\eqqcolon T_{j\bar k}}
\end{equation}
を得る. ここで$R_{j\bar k}$は$g$のRicci曲率である.
$T_{j\bar k}$は固定された計量のテンソルなので, 上からも下からも抑えられる. 
\item Lemma~\ref{lem-C2-estimate}の$C^2$評価より, ある定数$\Lambda$が存在して
\begin{equation}
\label{eq-C2-estimate}
\Lambda^{-1}(\hat g_{j\bar k})<(g_{j\bar k})<\Lambda(\hat g_{j\bar k})
\end{equation}
ここでこの$C^2$評価の仮定からLemma \ref{lem-sze-3.6}の$C^0$評価も従う. 
\item Christoffel記号の差
\begin{equation}
\label{eq-Sijk-defn}
S^{i}_{jk}\coloneqq \Gamma^{i}_{jk}-\hat\Gamma^{i}_{jk}
% \tag{37}
\end{equation}
とおく. $\Gamma^{i}_{jk}$ (resp. $\hat\Gamma^{i}_{jk}$) は$g$ (resp. $\hat g$)のLevi--Civita接続のChristoffel記号とする. 
\end{itemize}
\end{tcolorbox}
%最後に関しては, $\varphi$ の混合三階微分を評価するために, Christoffel 記号\(\Gamma^{i}_{jk}=g^{i\bar\ell}\partial_{j}g_{k\bar\ell}\)を評価する. 
なぜか\cite[Subsection 3.3]{Sze14}はここで急に記号が変わり, $\hat g$が固定した計量になる. 
評価するのは$g$であることに注意する. 
\begin{lem}[{\cite[Lemma 3.9]{Sze14}}]
\label{lem-deltaS-estimate}
\eqref{eq-ric-Tij}と\eqref{eq-C2-estimate}の下で, $X,T,\hat g,\Lambda$のみに依存する定数$C$が存在して
\[
\Delta \lvert S\rvert_{g}^{2}\geq -C\lvert S\rvert_{g}^{2}-C
\]
が成り立つ. ここで$\lvert S\rvert$は計量$g$により測ったテンソル$S$のノルムで, $\Delta$は$g$に関するLaplacianである.
\end{lem}

\begin{proof}
以下$g$のnormal coordinateを用いる., 点$p$において
\begin{equation}
\label{eq-normal-cordinate-3}
g=\sum_{i=1}^{n}dz^i \wedge\,d\overline{z^i}, \quad 
\frac{\partial}{\partial z^i}g_{j\bar k}(p)
=
\frac{\partial}{\partial \bar z^i}g_{j\bar k}(p)=0,
\end{equation}
とする. 
すると
\[
\lvert S\rvert_{g}^{2}
\underalign{\text{(ノルムの定義)}}{=}
g^{j\bar k}g_{a\bar b}g^{p\bar q}S^{a}_{jp}\overline{S^{b}_{kq}}
\underalign{\eqref{eq-normal-cordinate-3}}{=}
\sum_{a,j,p}
S^{p}_{ja}\overline{S^{p}_{ja}},
\]
となる. 
すると, 
\begin{equation}
\label{eq-laplacian-S}
\begin{aligned}
\Delta\lvert S\rvert_{g}^{2}
&\underalign{\text{(Lem. \ref{lem-cpx-Laplace})}}{=}
\sum_{p, i,j,k} \nabla_{p}\nabla_{\bar p}\bigl(S^{i}_{jk}\overline{S^{i}_{jk}}\bigr) \\
&\underalign{\text{}}{=}
\sum_{p, i,j,k} 
(\nabla_{p}\nabla_{\bar p}S^{i}_{jk})\overline{S^{i}_{jk}}
+S^{i}_{jk}\,\overline{(\nabla_{\bar p}\nabla_{p}S^{i}_{jk})} 
+\underbrace{(\nabla_{p}S^{i}_{jk})\overline{(\nabla_{p}S^{i}_{jk})}}_{\geq 0}
+\underbrace{(\nabla_{\bar p}S^{i}_{jk})\overline{(\nabla_{\bar p}S^{i}_{jk})} }_{\geq 0}
\\
&\underalign{\text{}}{\geq}
\sum_{p, i,j,k} 
\left( (\nabla_{p}\nabla_{\bar p}S^{i}_{jk})\overline{S^{i}_{jk}}
+S^{i}_{jk}\,\overline{(\nabla_{\bar p}\nabla_{p}S^{i}_{jk})}\right)
%\tag{38}
\end{aligned}
\end{equation}
である. 
%\footnote{\eqref{eq-normal-cordinate-3}から$\nabla_{p}$と$\partial_p$は本質的に変わらないと思う. }
ここでいくつかRiemannian幾何学から由来する恒等式などを思い出しておく. 


\begin{tcolorbox}[mybox]
$(X,g)$をK\"ahler多様体とする. このとき次が成り立つ.
\begin{itemize}
\item 一般に共変微分の交換子において, 上付きindexには$+R$, 下付きindexには$-R$が作用する. すなわち, 例えばvector $V^i$ ($(1,0)$-tensor)とcovector $\alpha_j$ ($(0,1)$-tensor)に対して
\[
[\nabla_p,\nabla_{\bar q}]V^i
=
R^i{}_{mp\bar q}V^m,
\qquad
[\nabla_p,\nabla_{\bar q}]\alpha_j
=
-R^m{}_{jp\bar q}\alpha_m
\]
が成り立つ.

\item $(1,2)$-tensor $S=S^i{}_{jk}$に対して, 共変微分の交換公式
\[
\nabla_p\nabla_{\bar q}S^i{}_{jk}
-
\nabla_{\bar q}\nabla_pS^i{}_{jk}
=
[\nabla_p,\nabla_{\bar q}]S^i{}_{jk}
=
R^i{}_{mp\bar q}S^m{}_{jk}
-R^m{}_{jp\bar q}S^i{}_{mk}
-R^m{}_{kp\bar q}S^i{}_{jm}
\]
が成り立つ. 特にnormal coordinateにおいて$q=p$として$p$について和を取ると
\begin{equation}
\label{eq-covariant-formula}
\sum_p\nabla_{\bar p}\nabla_pS^i{}_{jk}
=
\sum_p\nabla_p\nabla_{\bar p}S^i{}_{jk}
+
R^m{}_jS^i{}_{mk}
+
R^m{}_kS^i{}_{jm}
-
R^i{}_mS^m{}_{jk},
\end{equation}
ここで添字を上げたRicciテンソルを
$R^{m}_{j}
\coloneqq g^{m\bar k}R_{j\bar k}$とする. 
すると以下が成り立つ. 
\[
\sum_{p, i,j,k}R^{m}_{jp\bar p}S^{i}_{mk}
=
\sum_{p, i,j,k}g^{p\bar{p}}R^{m}_{jp\bar p}S^{i}_{mk}
=
\sum_{p, i,j,k}R^{m}_{j}S^{i}_{mk}
\]

\item K\"ahler多様体上の第二Bianchi恒等式より
\[
\nabla_pR^i{}_{jk\bar l}
=
\nabla_kR^i{}_{jp\bar l}
\]
が成り立つ.特に$\bar l=\bar p$として$p$について縮約すると
\begin{equation}
\label{eq-Bianchi-formula}
\sum_p\nabla_pR^i{}_{jk\bar p}
=
\sum_p\nabla_kR^i{}_{jp\bar p}
=
\nabla_kR^i{}_j.
\end{equation}
\end{itemize}
\end{tcolorbox}
これを元に$\nabla_{p}\nabla_{\bar p}S^{i}_{jk}$を評価していく. 
まず
\begin{equation}
\label{eq-pbarp-tensor}
\begin{aligned}
\nabla_{p}\nabla_{\bar p}S^{i}_{jk}
&\underalign{\text{($S$の定義 \eqref{eq-Sijk-defn})}}{=}\nabla_{p}\partial_{\bar p}\bigl(\Gamma^{i}_{jk}-\hat\Gamma^{i}_{jk}\bigr) 
\underalign{\text{($R^{i}_{jk\bar{l}}$の定義 in Sum. \ref{suma-brief-diff})}}{=}
-\nabla_{p}\bigl(R^{i}_{jk\bar p}-\hat R^{i}_{jk\bar p}\bigr) \\
&\underalign{}{=}
-\underbrace{\nabla_{p}R^{i}_{jk\bar p}}_
{=\nabla_{k}R^{i}_{j} \text{ by \eqref{eq-Bianchi-formula}}}
+\hat{\nabla_{p}}\hat R^{i}_{jk\bar p}
+(\nabla_{p}-\hat\nabla_{p})\hat R^{i}_{jk\bar p}.
\end{aligned}
\end{equation}
%ここで Bianchi 恒等式 $\nabla_{p}R^{i}_{jk\bar p}=\nabla_{k}R^{i}_{jp\bar p}=\nabla_{k}R^{i}_{j}$ を用いた。また 
ここで$\hat\nabla,\hat R$は$\hat g$のLevi--Civita接続と曲率テンソルである. 
この右辺にある3つの項を評価すると, 
\begin{itemize}
\item $\nabla_{k}R^{i}_{j}$については, $T_{k \bar l}$は固定された計量のテンソルなので
\begin{align*}
&\underbrace{R_{j\bar k}=-2\pi g_{j\bar k}+
T_{j\bar k}}_{=\eqref{eq-ric-Tij}}
\underalign{}{\Rightarrow}
R_{j}^{i}
=g^{i \bar k}R_{j\bar k}
=g^{i \bar k}(-2\pi g_{j\bar k}+
T_{j\bar k})
=-2\pi \delta^{i}_{j} + g^{i \bar k}T_{j\bar k}
\\
&\underalign{}{\Rightarrow}
\nabla_{k}R^{i}_{j}=\nabla_{k}(g^{i \bar l}T_{j\bar l})
\underalign{\text{($\nabla_{k}g_{j\bar l}=0$ by K\"ahler 条件)}}{=}
g^{i \bar l}(\nabla_{k}T_{j\bar l})
\underalign{\text{(Rem. \ref{rem-BG13-formula})}}{=}
g^{i \bar l}(\partial_{k} -\Gamma^j_{ik})T_{j\bar l}
\\
&\underalign{}{\Rightarrow}
\lvert\nabla_{k}R^{i}_{j}\rvert_{g}^{2}
\leq 
\underbrace{\lvert g^{i \bar l}\partial_{k}T_{j\bar l}\rvert_{g}^{2}}_{\text{(上への一様評価 by $C^0$評価)}}
+
\underbrace{\lvert g^{i \bar l}\Gamma^j_{ik}T_{j\bar l}\rvert_{g}^{2}}_{\leq C \lvert S \rvert_g^2}
\underalign{}{\leq} C' + C \lvert S \rvert_g^2
\end{align*}

\[\nabla_{k}R^{i}_{j}=-2\pi\nabla_{k}g_{j\bar k}+\nabla_{k}T^{i}_{j}=\nabla_{k}T^{i}_{j}\]より, $\nabla_{k}R^{i}_{j}
$も評価できる. ($\nabla_{k}g_{j\bar k}=0$はkahler conditionから)

\item $\hat R^{i}_{jk\bar p}$は固定された計量のtensorであり, $\hat{\nabla_{p}}$も固定された計量の共変微分なので, $g$の$C^0$評価より
\[
\lvert \hat{\nabla_{p}}\hat R^{i}_{jk\bar p} \rvert_{g}^{2}
\leq C
\]
と上への一様評価がある. 
\item $(\nabla_{p}-\hat\nabla_{p})\hat R^{i}_{jk\bar p}$については, Remark~\ref{rem-BG13-formula}によって
\begin{align*}
&(\nabla_{p}-\hat\nabla_{p})\hat R^{i}_{jk\bar p}
\underalign{\text{(Rem. \ref{rem-BG13-formula})}}{=}
S^i{}_{pa}\hat R^a{}_{jk\bar q}-S^a{}_{pj}\hat R^i{}_{ak\bar q}
-S^a{}_{pk}\hat R^i{}_{ja\bar q}
\\
&\underalign{}{\Rightarrow}
\lvert(\nabla_{p}-\hat\nabla_{p})\hat R^{i}_{jk\bar p}\rvert_{g}^{2}
\leq C \lvert S \rvert_g^2 \cdot \underbrace{\lvert\hat R^{i}_{jk\bar p}\rvert_{g}^{2}}_{\text{(上への一様評価 by $C^0$評価)}}
\leq C' \lvert S\rvert_g^2
\end{align*}
%接続の差 $\nabla_{p}-\hat\nabla_{p}$ はChristoffel 記号の差$S^{i}_{jk}=\Gamma^{i}_{jk}-\hat\Gamma^{i}_{jk}$で評価できる. 
\end{itemize}
以上より上の3つの評価から, ある定数$C_{2},C_{3}$があって
\begin{equation}
\label{eq-nabla-pbarp-estimate}
\bigl\lvert\nabla_{p}\nabla_{\bar p}S^{i}_{jk}\bigr\rvert_{g}^{2}
\underalign{\eqref{eq-pbarp-tensor}}{\leq} 
\underbrace{\lvert\nabla_{p}R^{i}_{jk\bar p}\rvert_{g}^{2}}_{}
+
\lvert\hat{\nabla_{p}}\hat R^{i}_{jk\bar p}\rvert_{g}^{2}
+
\lvert(\nabla_{p}-\hat\nabla_{p})\hat R^{i}_{jk\bar p}\rvert_{g}^{2}
\underalign{}{\leq} C_{2}\lvert S\rvert_{g}^{2}+C_{3}
\end{equation}

次に$\nabla_{\bar p}\nabla_{p}S^{i}_{jk}$を評価する. 
これは共変微分の公式より, $\nabla_{p}\nabla_{\bar p}S^{i}_{jk}$に帰着できる. 実際, 
\begin{equation}
\label{eq-nabla-pbarp-estimate-2}
\begin{aligned}
&%\sum_{p, i,j,k} \nabla_{\bar p}\nabla_{p}S^{i}_{jk}\underalign{\eqref{eq-covariant-formula}}{=}
%\sum_{p, i,j,k}\left( \nabla_{p}\nabla_{\bar p}S^{i}_{jk}+R^{m}_{jp\bar p}S^{i}_{mk}+R^{m}_{kp\bar p }S^{i}_{jm}-R^{i}_{m p\bar p}S^{m}_{jk} \right)
%\\&=
%\sum_{p, i,j,k} \left( \nabla_{p}\nabla_{\bar p}S^{i}_{jk}+R^{m}_{j}S^{i}_{mk}+R^{m}_{k}S^{i}_{jm}-R^{i}_{m}S^{m}_{jk}
\sum_p\nabla_{\bar p}\nabla_pS^i{}_{jk}
 \underalign{\eqref{eq-covariant-formula}}{=}
\sum_p\nabla_p\nabla_{\bar p}S^i{}_{jk}
+
R^m{}_jS^i{}_{mk}
+
R^m{}_kS^i{}_{jm}
-
R^i{}_mS^m{}_{jk},
\\
&\underalign{}{\Rightarrow}
\lvert\nabla_{\bar p}\nabla_{p}S^{i}_{jk}\rvert_{g}^{2}
\underalign{}{\leq}
\underbrace{\lvert\sum_p\nabla_p\nabla_{\bar p}S^i{}_{jk}\rvert_{g}^{2}}
_{\leq C_{2}\lvert S\rvert_{g}^{2}+C_{3} \text{ by \eqref{eq-nabla-pbarp-estimate}}}
+
\underbrace{\lvert R^m{}_jS^i{}_{mk}\rvert_{g}^{2}}
_{\leq C\lvert S\rvert_{g}^{2} \text{ by $C^2$評価}}
+
\underbrace{\lvert R^m{}_kS^i{}_{jm}\rvert_{g}^{2}}
_{\leq C\lvert S\rvert_{g}^{2} \text{ by $C^2$評価}}
+
\underbrace{\lvert R^i{}_mS^m{}_{jk}\rvert_{g}^{2}}
_{\leq C\lvert S\rvert_{g}^{2} \text{ by $C^2$評価}}
\\
&\underalign{}{\Rightarrow}
\lvert\nabla_{\bar p}\nabla_{p}S^{i}_{jk}\rvert_{g}^{2}
\underalign{}{\leq} 
C_{4}\lvert S\rvert_{g}^{2}+C_{5}
\end{aligned}
\end{equation}
%\eqref{eq-ric-Tij}と\eqref{eq-C2-estimate}より$g$のRicci テンソルは有界なので,$C_1$を定数として\[
%\bigl\lvert\nabla_{\bar p}\nabla_{p}S^{i}_{jk}\bigr\rvert
%\le
% \bigl\lvert\nabla_{p}\nabla_{\bar p}S^{i}_{jk}\bigr\rvert
 %+C_{1}\lvert S\rvert
% \tag{39}
%\]
%となる. 
以上より
\begin{align*}
\Delta\lvert S\rvert^{2}
&\underalign{\eqref{eq-laplacian-S}}{\geq}
\sum_{p, i,j,k} 
\left( (\nabla_{p}\nabla_{\bar p}S^{i}_{jk})\overline{S^{i}_{jk}}
+S^{i}_{jk}\,\overline{(\nabla_{\bar p}\nabla_{p}S^{i}_{jk})}\right)
\\
&\underalign{\text{(Young不等式 as in \eqref{eq-Poincare-2})}}{\geq}
-C_6\sum_{} 
\left( \lvert \nabla_{p}\nabla_{\bar p}S^{i}_{jk}\rvert^{2}+
\lvert S^{q}_{rs}\rvert^{2}
+\lvert S^{q}_{rs} \rvert^{2}+
\lvert \nabla_{\bar p}\nabla_{p}S^{i}_{jk}\rvert^{2}
\right)
\underalign{\text{\eqref{eq-nabla-pbarp-estimate} \& \eqref{eq-nabla-pbarp-estimate-2}}}{\geq}
-C_{7}\lvert S\rvert^{2}-C_{8}.
\end{align*}
\end{proof}


\begin{lem}
[{\cite[Lemma 3.10]{Sze14}}]
\label{lem-Sze-3.10}
\eqref{eq-ric-Tij}と\eqref{eq-C2-estimate}の下で, $X,T,\hat g,\Lambda$のみに依存する定数$C$が存在して$\lvert S\rvert_g\leq C$が成り立つ.
\end{lem}
\begin{proof}
\eqref{eq-delta-tr}により, ある定数$\varepsilon,C_{1}>0$を選んで
\begin{equation}
\label{eq-trace-laplacial-estimate-2}
\Delta\operatorname{tr}_{\hat{g}} g
\underalign{\eqref{eq-delta-tr}}{\geq}
 \underbrace{
 -B(\operatorname{tr}_{g}{\hat{g}})(\operatorname{tr}_{\hat{g}}g)
 -{\hat{g}}^{j\bar k}R_{j\bar k}
 }_{ \geq -C_1 \text{ by \eqref{eq-ric-Tij} \& \eqref{eq-C2-estimate}}}
+
\sum_{p,j,a} g^{p\bar p}g^{a\bar a}\,\lvert\partial_j g_{p\bar a}\rvert^2
\underalign{}{\geq}
-C_{1}+\varepsilon \lvert S\rvert_g^{2}
\end{equation}
とできる. Lemma~\ref{lem-deltaS-estimate}における$C>0$について, 
$A\coloneqq \frac{C+1}{\varepsilon}>0$, $C_{2}\coloneqq C+AC_1>0$とすれば, 
\[
\Delta\bigl(\lvert S\rvert_g^{2}+A\operatorname{tr}_{\hat g}g\bigr)
\underalign{\text{(Lem. \ref{lem-deltaS-estimate}) \& \eqref{eq-trace-laplacial-estimate-2}}}{\geq} 
(-C+A\varepsilon) \lvert S\rvert_g^{2}-C-AC_1 
=
\lvert S\rvert_g^{2}-C_2
\]
となる. 
点$p\in X$で$\lvert S\rvert^{2}+A\operatorname{tr}_{\hat g}g$が最大値を達成するとすると, 
\begin{align*}
 &0\underalign{\text{(最大値)}}{\geq}\Delta\bigl(\lvert S\rvert^{2}+A\operatorname{tr}_{\hat g}g\bigr)
\geq 
 \lvert S\rvert^{2}(p)-C_{2}
\underalign{}{\Rightarrow} 
\lvert S\rvert^{2}(p)\leq C_{2}
\\
&\underalign{}{\Rightarrow} 
\lvert S\rvert^{2}(x)
\leq \lvert S\rvert^{2}(x)+A\operatorname{tr}_{\hat g}g(x)
\underalign{\text{($p$で最大値)}}{\leq} \lvert S\rvert^{2}(p)+A\operatorname{tr}_{\hat g}g(p)\leq C_{2}+C_{3}
\quad \forall x \in X
\end{align*}
\end{proof}

\subsection{Proposition~\ref{prop-uniform-estimate}の証明}

\begin{prop}[= Proposition~\ref{prop-uniform-estimate}]
$X,\omega_{0}$および$F$のみに依存する定数$C>0$が存在し,$\varphi_{t}$が
\[
\left(\omega_{0}+\frac{\sqrt{-1}}{2\pi}\partial\bar\partial\varphi_{t}\right)^{n}
=e^{tF+\varphi_{t}}\omega_{0}^{n}
\]
を満たすならば,ある$t\in[0,1]$に対して
\[
\left(g_{j\bar k}+\frac{1}{2\pi}\partial_{j}\partial_{\bar k}\varphi_{t}\right)>C^{-1}(g_{j\bar k}),
\quad \text{かつ} \quad 
\lVert\varphi_{t}\rVert_{C^{3,\alpha}(X)}\leq C
\]
が成り立つ.ここでH\"olderノルムは計量$\omega_{0}$に関して測られる.
\end{prop}

\begin{proof}
$(g_{j\bar k}+\frac{1}{2\pi}\partial_{j}\partial_{\bar k}\varphi_{t})>C^{-1}(g_{j\bar k})$に関してはLemma~\ref{lem-C2-estimate}の$C^2$評価による. 

$0<\alpha<1$を固定する. Lemma~\ref{lem-Sze-3.10}と$C^0$評価から
$\lVert\varphi_t\rVert_{C^{3}}\leq C$
を得る. 特に$\varphi$の$C^{2, \alpha}$評価を得る. よってTheorem~\ref{thm-open-2}と同じ議論で, $\lVert\varphi\rVert_{C^{3,\alpha}}$の評価を得る. 
\end{proof}


\subsection{$c_{1}(X)=0$の場合}

おそらくこの内容は\cite[Subsection 3]{Blo07}のレクチャーノートの方が詳しい.
\begin{thm}
[{\cite[Theorem 3.14]{Sze14}}]
$(X,\omega)$をコンパクトK\"ahler多様体とし,
$C^\infty$級関数$F\colon X\to\mathbb{R}$が
\[
\int_X e^F\omega^n=\int_X\omega^n
\]
を満たすとする.

このとき,定数和を除いて一意な$C^\infty$級関数$\varphi\colon X\to\mathbb{R}$が存在して,
\(
\omega+\frac{\sqrt{-1}}{2\pi}\partial\bar\partial\varphi
\)
はK\"ahler formであり,次のMonge--Amp\`ere方程式を満たす.
\[
\left(
\omega+\frac{\sqrt{-1}}{2\pi}\partial\bar\partial\varphi
\right)^n
=
e^F\omega^n.
\]
\end{thm}

$c_{1}(X)<0$の場合との証明の違いは$C^0$評価のみである. それ以外はほとんど同じなので$C^0$評価のみに焦点を当てる.

\begin{prop}
[{\cite[Proposition 3.15]{Sze14}}]
$(X,\omega)$をコンパクトK\"ahler多様体とし,$F,\varphi\colon X\to\mathbb{R}$を$C^\infty$級関数とする. $\omega-\frac{\sqrt{-1}}{2\pi}\partial\bar\partial\varphi$
がK\"ahler formであり,
\begin{equation}
\label{eq-Calabi-Yau-eq}
\left(
\omega-\frac{\sqrt{-1}}{2\pi}\partial\bar\partial\varphi
\right)^n
=
e^F\omega^n
\end{equation}
を満たすと仮定する.このとき,$(X,\omega)$と$\sup_XF$のみに依存する定数$C$が存在して次が成り立つ. 
\[
\sup_X\varphi-\inf_X\varphi<C
\]
\end{prop}
どうも\(\omega-\frac{\sqrt{-1}}{2\pi}\partial\bar\partial\varphi
\)
を用いているのは符号を簡単にするためらしい. 
\begin{proof}
\cite[Subsection 3]{Blo07}にあるようなMoser iterationに基づく. それは
\(
\lVert \varphi\rVert_p
\coloneqq 
\left(
\int_X|\varphi|^p\omega^n
\right)^{1/p}
\)
という$L^p$-normを, $p$を反復的に大きくしながら評価し, 最後に$p\to\infty$とするものである.

以下定数$C_1, C_2, \ldots, $は$(X,\omega)$と$\sup_XF$にのみ依存するものとする. 
また正規化して, 
\begin{equation}
\label{eq-normalization-omega}
\inf_X\varphi=1,
\quad \text{and} \quad
\int_X\omega^n=1
\end{equation}
とする. 特に$\varphi\geq1$なので, 
$p\leq q$に対して
\(
\lVert \varphi\rVert_p\leq\lVert \varphi\rVert_q
\)
である. 

$\Delta$を$\omega$に関するLaplacianとすると, 
\begin{equation}
\label{eq-2pin-laplacian}
\omega-\frac{\sqrt{-1}}{2\pi}\partial\bar\partial\varphi>0
\underalign{\text{(Def. \ref{defn-cpx-Laplace})}}{\Rightarrow} 
n-\frac{1}{2\pi}\Delta\varphi
=\operatorname{tr}(\omega-\frac{\sqrt{-1}}{2\pi}\partial\bar\partial\varphi)
>0
\underalign{}{\Rightarrow} 
2\pi n > \Delta\varphi 
\end{equation}
ここでGreen関数を復習する.
\begin{defn}
\(M\)をコンパクトRiemannian多様体とし,体積形式を\(dV\)と書く.
Green関数\(G(P,Q) \in C^{\infty}(M \times M \setminus \Delta_M)\) 
とは
\[
-\Delta_Q G(P,Q)
=
\delta_P-\frac{1}{\Vol(M)}
\]
を満たす関数である. ここで\(\delta_P\)はDirac measureである.
これは定数和を除いて一意に定まる.
%さらに正規化として\[\int_M G(P,Q)\,dV_Q=0\]を課す．この正規化により \(G(P,Q)\) は定数の不定性を除いて一意に定まる．

任意のsmooth function \(\varphi\in C^\infty(M)\)に対して
\begin{equation}
\label{eq-Green-1}
\varphi(P)
=
\frac{1}{\Vol(M)}\int_M\varphi\,dV
+
\int_M G(P,Q)(-\Delta\varphi)(Q)\,dV_Q
\end{equation}
が成り立つ.実際部分積分によって
\[
\begin{aligned}
\int_M G(P,Q)(-\Delta\varphi)(Q)\,dV_Q
&\underalign{\text{(部分積分)}}{=}
\int_M \varphi(Q)(-\Delta_QG(P,Q))\,dV_Q
\\
&\underalign{\text{(定義)}}{=}
\int_M \varphi(Q)
\left(
\delta_P-\frac{1}{\Vol(M)}
\right)dV_Q
=
\varphi(P)
-
\frac{1}{\Vol(M)}\int_M\varphi\,dV.
\end{aligned}
\]
これを移項して上の表示公式\eqref{eq-Green-1}を得る.
\end{defn}

$\varphi(p)=1$となる点$p\in X$を取り,$G(x,y)$を$\Delta$のGreen's functionとする.
定数を調整して$G\geq0$とすると
\begin{equation}
\label{eq-L1-bound-estimate}
\begin{aligned}
&\varphi(p)
\underalign{\text{\eqref{eq-normalization-omega} \& \eqref{eq-Green-1}}}{=}
\int_X\varphi\,\omega^n
-
\int_XG(x,p)\Delta\varphi(x)\,\omega^n(x)
\\
&\underalign{}{\Rightarrow}
1=\varphi(p)
\underalign{\eqref{eq-2pin-laplacian}}{\geq}
\int_X\varphi\,\omega^n
-
2\pi n\int_XG(x,p)\,\omega^n(x)
\geq
\int_X\varphi\,\omega^n-C_1
\\
&\underalign{}{\Rightarrow}
\lVert \varphi\rVert_1
=
\int_X\varphi\,\omega^n<C_2
\end{aligned}
\end{equation}
さて\(
\omega_\varphi
\coloneqq
\omega-\frac{\sqrt{-1}}{2\pi}\partial\bar\partial\varphi
\)とすると
\begin{equation}
\label{eq-omega-varphi-omega-1}
\begin{aligned}
\int_X\varphi(\omega_\varphi^n-\omega^n)
&\underalign{}{=}
\int_X
\varphi(\omega_\varphi-\omega)
\wedge
\left(
\omega_\varphi^{n-1}
+\omega_\varphi^{n-2}\wedge\omega
+\cdots
+\omega_\varphi \wedge\omega^{n-2}
+\omega^{n-1}
\right)\\
&\underalign{\text{(定義)}}{=}
-\frac{1}{2\pi}
\int_X
\varphi\sqrt{-1}\partial\bar\partial\varphi
\wedge
\left(
\omega_\varphi^{n-1}
+\omega_\varphi^{n-2}\wedge\omega
+\cdots
+\omega_\varphi\wedge\omega^{n-2}
+\omega^{n-1}
\right)\\
&\underalign{\text{(Stokesの定理)}}{=}
\frac{1}{2\pi}
\int_X
\sqrt{-1}\partial\varphi\wedge\bar\partial\varphi
\wedge
\left(
\omega_\varphi^{n-1}
+\omega_\varphi^{n-2}\wedge\omega
+\cdots
+\omega_\varphi\wedge\omega^{n-2}
+\omega^{n-1}
\right).
\end{aligned}
\end{equation}
ここでStokesの定理を用いた.各$0\leq k\leq n-1$に対して
\(
\sqrt{-1}\partial\varphi\wedge\bar\partial\varphi
\wedge
\omega_\varphi^k\wedge\omega^{n-1-k}
\geq 0\)
である. (これは各点で計算すれば良い). よって
\begin{equation}
\label{eq-L2-estimate-1}
\begin{aligned}
&\int_X\varphi
\underbrace{(\omega_\varphi^n-\omega^n)}_{=(e^F-1)\omega^n \text{ by \eqref{eq-Calabi-Yau-eq}}}
\underalign{\eqref{eq-omega-varphi-omega-1}}{\geq}
\frac{1}{2\pi}
\int_X
\sqrt{-1}\partial\varphi\wedge\bar\partial\varphi
\wedge\omega^{n-1}
=
\frac{1}{2\pi n}
\int_X|\partial\varphi|_\omega^2\omega^n
\\
&\underalign{}{\Rightarrow}
\frac{1}{2\pi n}
\int_X|\partial\varphi|_\omega^2\omega^n
\leq
\int_X\varphi(e^F-1)\omega^n
\underalign{\text{($F$を定数で抑える)}}{\leq}
C_3\int_X\varphi\,\omega^n
\leq C_4.
\\
&\underalign{}{\Rightarrow}
\int_X(\varphi-\lVert \varphi\rVert_1)^2\omega^n
=
\int_X
\left(
\varphi-\int_X\varphi\,\omega^n
\right)^2
\omega^n
\underalign{\text{(Poincar\'e inequality  by Thm. \ref{thm-poincare-inequality})}}{\leq}
C_P\int_X|\partial\varphi|_\omega^2\omega^n
\leq C_5
\\
&\underalign{\eqref{eq-L1-bound-estimate}}{\Rightarrow}
\lVert \varphi\rVert_2<C_6
\end{aligned}
\end{equation}
これを$p\geq2$に対して同様のこと行う. つまり
\begin{equation}
\label{eq-omega-varphi-omega-2}
\begin{aligned}
\int_X\varphi^{p-1}(\omega_\varphi^n-\omega^n)
&\underalign{\text{(定義)}}{=}
-\frac{1}{2\pi}
\int_X
\varphi^{p-1}\sqrt{-1}\partial\bar\partial\varphi
\wedge
\left(
\omega_\varphi^{n-1}
+\omega_\varphi^{n-2}\wedge\omega
+\cdots
+\omega_\varphi \wedge\omega^{n-2}
+\omega^{n-1}
\right)
\\
&\underalign{\text{(Stokesの定理)}}{=}
\frac{p-1}{2\pi}
\int_X
\sqrt{-1}\varphi^{p-2}
\partial\varphi\wedge\bar\partial\varphi
\wedge
\left(
\omega_\varphi^{n-1}
+\omega_\varphi^{n-2}\wedge\omega
+\cdots
+\omega_\varphi \wedge\omega^{n-2}
+\omega^{n-1}
\right)
\\
&\underalign{\text{}}{=}
\frac{4(p-1)}{2\pi p^2}
\int_X
\sqrt{-1}\partial\varphi^{p/2}
\wedge\bar\partial\varphi^{p/2}
\wedge
\left(
\omega_\varphi^{n-1}
+\cdots
+\omega^{n-1}
\right)\\
&\underalign{\text{(\eqref{eq-L2-estimate-1}と同じ)}}{\geq}
\frac{4(p-1)}{2\pi np^2}
\int_X
\left\lvert
\partial\varphi^{p/2}
\right\rvert_\omega^2
\omega^n.
\end{aligned}
\end{equation}
よって
\begin{equation}
\label{eq-L2-estimate-2}
\begin{aligned}
\frac{4(p-1)}{2\pi np^2}
\int_X
\left\lvert
\partial\varphi^{p/2}
\right\rvert_\omega^2
\omega^n
&\underalign{\eqref{eq-omega-varphi-omega-2}}{\leq}
\int_X\varphi^{p-1}(\omega_\varphi^n-\omega^n)
\underalign{\text{(\eqref{eq-L2-estimate-1}と同じ)}}{=}
\int_X\varphi^{p-1}(e^F-1)\omega^n\\
&\underalign{\text{(\eqref{eq-L2-estimate-1}と同じ)}}{\leq}
C_7\int_X\varphi^{p-1}\omega^n
\underalign{}{=}
C_7\lVert \varphi\rVert_{p-1}^{p-1}.
\end{aligned}
\end{equation}
これにより, $p\geq2$であるので
\begin{equation}
\label{eq-Sobolev-estimate-1}
\left\lVert
\partial\varphi^{p/2}
\right\rVert_2^2
\underalign{\eqref{eq-L2-estimate-2}}{\leq}
\frac{2\pi np}{4(p-1)} C_7 p \lVert \varphi\rVert_{p-1}^{p-1}
\underalign{(p\geq2)}{\leq}
C_8 p\lVert \varphi\rVert_{p-1}^{p-1}.
\end{equation}
$C_8$は$p$に依存しないことに注意する. 

Sobolev埋め込み定理(Theorem~\ref{thm-Sobolev})において, 
$N=2n, m=1, p=2$とすると$q=\frac{Np}{N-mp}=\frac{2n}{n-1}$とすれば, $mp<N$なので, 
\(
W^{m,p}(\Omega)=W^{1,2}(\Omega)\hookrightarrow L^{\frac{2n}{n-1}}(\Omega) =L^q(\Omega) 
\)
が連続埋め込みとなる. よって, $X$はコンパクトであることより, ある定数$C_S$が存在して任意の$f$に対して
\begin{equation}
\label{eq-Sobolev-estimate-2}
\lVert f\rVert_{L^{\frac{2n}{n-1}}}^2
\leq
C_S
\lVert f\rVert_{W^{1,2}}^2
=
C_S
\left(
\lVert f\rVert_2^2+\lVert \partial f\rVert_2^2
\right)
\end{equation}
が成り立つ.これを$f=\varphi^{p/2}$に適用すると, 任意の$p \geq 2$について
\begin{equation}
\label{eq-iteration-1}
\begin{aligned}
&\lVert \varphi\rVert_{\frac{np}{n-1}}^p
=
\left\lVert
\varphi^{p/2}
\right\rVert_{\frac{2n}{n-1}}^2
\underalign{\eqref{eq-Sobolev-estimate-2}}{\leq}
C_S
\left(
\left\lVert \varphi^{p/2}\right\rVert_2^2
+
\left\lVert \partial\varphi^{p/2}\right\rVert_2^2
\right)
\underalign{\eqref{eq-Sobolev-estimate-1}}{\leq}
C_S
\left(
\lVert \varphi\rVert_p^p
+
C_8p\lVert \varphi\rVert_{p-1}^{p-1}
\right).
\\
&\underalign{\text{($\varphi\geq1$ \& $\lVert \varphi\rVert_{p-1}\leq\lVert \varphi\rVert_p$)}}{\Rightarrow}
\lVert \varphi\rVert_{\frac{np}{n-1}}^p
\leq 
C_S
\left(
\lVert \varphi\rVert_p^p
+
C_8p\lVert \varphi\rVert_{p-1}^{p}
\right)
\leq 
C_S
\left(
\lVert \varphi\rVert_p^p
+
C_8p\lVert \varphi\rVert_{p}^{p}
\right)
\leq 
C_9p\lVert \varphi\rVert_{p}^{p}
\\
&\underalign{}{\Rightarrow}
\lVert \varphi\rVert_{\frac{np}{n-1}}
\leq
(Cp)^{1/p}\lVert \varphi\rVert_p.
\end{aligned}
\end{equation}
そこで
\(
p_k
\coloneqq 
2\left(
\frac{n}{n-1}
\right)^k 
\)
とおく.$p_k = \frac{n}{n-1}p_{k-1}$であることに注意すると
\begin{equation}
\label{eq-iteration-2}
\begin{aligned}
\lVert \varphi\rVert_{p_k}
\underalign{\eqref{eq-iteration-1}}{\leq}
(Cp_{k-1})^{1/p_{k-1}}
\lVert \varphi\rVert_{p_{k-1}}
\leq
\cdots\leq
\lVert \varphi\rVert_2\cdot
\prod_{i=0}^{k-1}
(Cp_i)^{1/p_i}
\leq
\lVert \varphi\rVert_2 \cdot
\prod_{i=0}^{\infty}
(Cp_i)^{1/p_i}
\end{aligned}
\end{equation}
ここで
\[
\sum_{i=0}^{\infty}\frac{1}{p_i}<+\infty, 
\sum_{i=0}^{\infty}\frac{\log p_i}{p_i}<+\infty
\underalign{}{\Rightarrow}
\log \prod_{i=0}^{\infty}
(Cp_i)^{1/p_i}
=
(\log C)\sum_{i=0}^{\infty}\frac1{p_i}
+
\sum_{i=0}^{\infty}\frac{\log p_i}{p_i}
<+\infty
\]
以上より$k\to\infty$とすれば
\[
\sup_X\varphi
=
\lVert \varphi\rVert_\infty
\underalign{\eqref{eq-iteration-2}}{\leq}
C'\lVert \varphi\rVert_2
\underalign{\eqref{eq-L2-estimate-1}}{\leq}
C''
\underalign{}{\Rightarrow}
\sup_X\varphi-\inf_X \varphi
\underalign{\eqref{eq-normalization-omega}}{=}\sup_X\varphi-1
<C''-1
\]
よって欲しい式が得られた.
\end{proof}


\newpage
\section{Yauの定理の応用 1 K\"ahler--Einstein計量の存在とMiyaoka-Yau不等式}
\label{sec-Yau-App1}
以後Yauの定理とその応用を調べる. Chern類などは\cite{Iwa-Chern}参照のこと.

\subsection{直線束のpositivityと$\partial \bar\partial$補題.}
直線束や曲率の定義は\cite[Section 3~5]{Dem12}もしくは\cite[Chapters 1 and 2]{Iwa-Chern}参照. 必要なものだけ定義しておく. 
\begin{defn}[{\cite[Section 5]{Dem12}}]
$(X, \omega)$をコンパクトK\"ahler多様体とし, $L$を直線束とする. 
\begin{itemize}
\item $L$がample (positive)であるとは, \(L\)の\(C^\infty\) Hermite計量\(h\)が存在して,
\(
 \sqrt{-1}F_{L,h}>0
\)
を満たすこと. 
\item $L$がsemipositiveであるとは, \(L\)の\(C^\infty\) Hermite計量\(h\)が存在して,
\(
 \sqrt{-1}F_{L,h}\geq 0
\)
を満たすこと. 
\item $L$がnefであるとは, 任意の\(\varepsilon>0\)に対して, \(L\)の\(C^\infty\) Hermite計量\(h_\varepsilon\)が存在して,
\(
 \sqrt{-1}F_{L,h_\varepsilon}\geq -\varepsilon\omega
\)
を満たすこと. 
\end{itemize}
\end{defn}
また次もよく使うので書いておく. 
%\begin{lem}[$\partial \overline{\partial}$補題]
%\label{lem-ddc-lem-2}
%$X$ がコンパクト K\"ahler 多様体で 
%$\alpha$が$d$-exactならば, $\alpha= \partial \overline{\partial} \psi$となる$\psi$が存在する. 
%\end{lem}

\begin{lem}[$\partial \bar\partial$補題]
\label{lem-ddc-lem-2}
$X$をコンパクトK\"ahler多様体とし, $\alpha\in A^{1,1}(X)$が$d$-exactであるとする.
このときある$\psi\in C^\infty(X)$が存在して
\(
\alpha=\partial\bar\partial\psi
\)
となる.
\end{lem}

\begin{proof}
Theorem~\ref{thm-harmonic-integral}より
直交分解
\[A^{1,1}(X)= \mathbf{H}^{1,1}(X)\oplus \bar\partial A^{1,0}(X) \oplus \bar\partial^{*}A^{1,2}(X)
\]
が存在する. 
$\alpha$は$d$-exactなので任意の調和形式に直交する.
したがってTheorem~\ref{thm-harmonic-integral}(2) より,
ある
$\varphi\in A^{1,1}(X)$が存在して
\(
\alpha=\square_{\bar\partial}\varphi
\)
となる. よって$\bar\partial\varphi=\partial\varphi=0$が示せれば, 
 K\"ahler identity
\begin{equation}
\label{eq-Kahler-indentity}
[\Lambda,\partial]=\sqrt{-1}\bar\partial^{*}
\iff
\bar\partial^*
=
-\sqrt{-1}
(\Lambda\partial-\partial\Lambda)
%\quad\text{and}\quad
\end{equation}
によって
\[
\alpha
=\underbrace{\square_{\bar\partial}}_{\coloneqq\bar\partial\bar\partial^{*}+\bar\partial^{*}\bar\partial}\varphi
=\bar\partial\bar\partial^{*}\varphi
\underalign{\eqref{eq-Kahler-indentity}}{=} -\sqrt{-1}
\bar\partial(\Lambda\partial-\partial\Lambda)\varphi
\underalign{\text{(by $\partial\varphi=0$)}}{=}-\sqrt{-1}\,\partial\bar\partial\Lambda\varphi.
\]
よって$\psi=-\sqrt{-1}\Lambda\varphi$とすればよい.

$\bar\partial\varphi=0$については, $d\alpha=\partial\alpha+\bar\partial\alpha=0$と
$\square_{\bar\partial}$と$\bar \partial$が(K\"ahlerでなくても)可換より, 
\[
\square_{\bar\partial}\bar\partial\varphi=
\bar\partial\square_{\bar\partial}\varphi
=\bar\partial\alpha=0
\] 
なので, $\bar\partial\varphi$と$\partial\varphi$は調和かつexactなので, $\bar\partial\varphi=0$になる.  (直交分解を用いても良いし随伴による積分からも言える.)
 
 $\partial\varphi=0$については, 上と同じく$\square_{\bar\partial}\partial\varphi=0$を示せば良い. 
 (ここにK\"ahlerがいる). K\"ahler identityから
 \[
\partial\bar\partial^*+\bar\partial^*\partial
\underalign{\eqref{eq-Kahler-indentity}}{=}
-\sqrt{-1}
\left\{
\partial(\Lambda\partial-\partial\Lambda)
+
(\Lambda\partial-\partial\Lambda)\partial
\right\}
=
-\sqrt{-1}
\left(
\partial\Lambda\partial
-\partial^2\Lambda
+\Lambda\partial^2
-\partial\Lambda\partial
\right)
=0
 \]
よってこれから$
\square_{\bar\partial}\partial
=
\partial\square_{\bar\partial}
 $がいえ, 上と同じ議論から$\square_{\bar\partial}\partial\varphi=0$である. 
\end{proof}


\subsection{Yauの定理とK\"ahler--Einstein計量の存在}
Section~\ref{Sec-MA-eq}で示した定理をYauの定理としてまとめておく.

\begin{thm}[Yauの定理]
\label{thm-Yau-theorem}
$(X, \omega)$をコンパクトK\"ahler多様体, $f$を$C^{\infty}$級関数, $\lambda\geq 0$とする. 
さらに$\lambda=0$の時は, 正規化条件
\(\int_X e^f\omega^n=\int_X\omega^n\)をかす.

この時$C^{\infty}$級関数\(\varphi\)であって
\[
\left(\omega+\frac{\sqrt{-1}}{2\pi}\partial\bar\partial\varphi\right)^n=e^{\lambda\varphi+f}\omega^n,
\qquad
\omega+\frac{\sqrt{-1}}{2\pi}\partial\bar\partial\varphi>0.
\]
となるものが存在する. 
\end{thm}

\begin{cor}
\label{cor-Yau-thm}
$(X, \omega_0)$をコンパクトK\"ahler多様体とする. 
\begin{enumerate}
\item $K_X$ ampleかつ$\{\omega_0\} \in c_1(K_X)$ならば, $\{ \omega\}=\{\omega_0\}$となるK\"ahler--Einstein計量$\Ric(\omega)=-\omega$が存在する
\item $c_1(K_X) =  0$の時, $\{ \omega\}=\{\omega_0\}$となるK\"ahler--Einstein計量$\Ric(\omega)=0$が存在する
\item $-K_X$ ampleの時, ある$\varepsilon>0$があって, $\{ \omega\}=\{\omega_0\}$かつ$\Ric(\omega) \geq \varepsilon \omega$なるK\"ahler計量が存在する.
%\item $-K_X$ nefの時, 任意の$\varepsilon>0$について$\{ \omega\}=\{\omega_0\}$かつ$\Ric(\omega) \ge - \varepsilon \omega$となるケーラー計量が存在する.
%\item $K_X$ nefの時, 任意の$\varepsilon>0$について$\{ \omega\}=\{\omega_0\}$かつ$\Ric(\omega) \le \varepsilon \omega_0$なるケーラー計量が存在する
\end{enumerate}
\end{cor}

K\"ahler--Einstein計量の定義をしておく. 
\begin{defn}[K\"ahler--Einstein計量]
$X$のK\"ahler計量$g$について, ある定数$\lambda \in \R$があって
\[
\Ric(\omega)=\lambda \omega
\Leftrightarrow
\underbrace{-\frac{\sqrt{-1}}{2\pi}\partial_i \partial_{\overline{j}} \log \det(g)}_{=\frac{\sqrt{-1}}{2\pi}R_{i \overline{j}}}
= \sqrt{-1}\lambda g_{i \overline{j}}
\Leftrightarrow
\frac{1}{2\pi}
R_{i \overline{j}}=\lambda g_{i \overline{j}}
\]
となる時$g$をK\"ahler--Einstein計量という. 
\end{defn}



\begin{proof}[Proof of Corollary~\ref{cor-Yau-thm}]
まず次の簡単なclaimを示す.
%%%%%%%%%%%%%
\begin{comment}
まず次を思い出す


\begin{tcolorbox}[mybox]
Ricci curvature tensor を
$$
R_{i\bar{j}}\coloneqq g^{k \bar{l}}R_{i \bar{j}k \bar{l}} 
$$
とすると次が成り立つ
$$
R_{i \bar{j}} = -\partial_i \partial_{\bar{j}} \log \det(g)
$$
Ricci 曲率は以下のように定める. 
$$
\Ric(\omega)\coloneqq  \frac{\sqrt{-1}}{2 \pi} R_{i\bar{j}}dz^i \wedge d\bar{z}^{\bar{j}}
\in c_1(X) \in H^{2}(X, \R)
$$
\end{tcolorbox}
%%%%%%%%%%%%%%%
\end{comment}
\begin{claim}
\label{claim-ric-det-formula}
2つのK\"ahler form $\omega, \omega'$について
\[
\Ric(\omega')-
\Ric (\omega)
= -\frac{\sqrt{-1}}{2\pi}\partial\bar\partial \log \frac{\omega'^n}{\omega^n}
\]
ここでlocal chartで
\[
\omega'^n\coloneqq h'\sqrt{-1}^n  n! dz^1 \wedge d \bar{z}^1 \wedge \cdots \wedge d z^n\wedge\,d\bar{z}^{\bar{n}} \quad 
\omega^n \coloneqq h \sqrt{-1}^n  n! dz^1 \wedge d \bar{z}^1 \wedge \cdots \wedge d z^n\wedge\,d\bar{z}^{\bar{n}}
\]
と書けている時, 
$
\frac{\omega'^n}{\omega^n}\coloneqq \frac{h'}{h}
$
という$C^\infty$級関数を表す. これはlocal chartの取り方によらない.
\end{claim}

\begin{proof}
$\omega=\sqrt{-1} g_{i \bar{j}}dz^i \wedge\,d\bar{z}^{\bar{j}}$, $\omega' =\sqrt{-1} g'_{i \bar{j}}dz^i \wedge\,d\bar{z}^{\bar{j}}$とすると
\[
\omega'^n
=\det(g'_{i \bar{j}}) \sqrt{-1}^n  n! dz^1 \wedge d \bar{z}^1 \wedge \cdots \wedge d z^n\wedge\,d\bar{z}^{\bar{n}}
\Rightarrow 
 \frac{\omega'^n}{\omega^n}=\frac{\det(g'_{i \bar{j}})}{\det(g_{i \bar{j}})}
\]
であるので, 
\begin{align*}
-\frac{\sqrt{-1}}{2\pi}\partial\bar\partial \log \frac{\omega'^n}{\omega^n}
 &\underalign{\text{上を代入}}{=}
 \underbrace{-\frac{\sqrt{-1}}{2\pi}\partial\bar\partial \log\det(g'_{i \bar{j}})}_{=\Ric(\omega') \text{上の定義}}
- \underbrace{(-\frac{\sqrt{-1}}{2\pi}\partial\bar\partial \log\det(g_{i \bar{j}}))}_{=\Ric(\omega) \text{上の定義}}
\\
\end{align*}
\end{proof}

以上よりMonge--Amp\`ere方程式$\left(\omega+\frac{\sqrt{-1}}{2\pi}\partial\bar\partial\varphi\right)^n=e^{\lambda\varphi+f}\omega^n$
に関して上のClaimを適応すると
\begin{equation}
\label{eq-MA-equation-ricci}
\begin{aligned}
\Ric\left(\omega+\frac{\sqrt{-1}}{2\pi}\partial\bar\partial\varphi\right)-\Ric (\omega)
&\underalign{\text{(Claim~\ref{claim-ric-det-formula})}}{=}
 -\frac{\sqrt{-1}}{2\pi}\partial\bar\partial \log \frac{\left(\omega+\frac{\sqrt{-1}}{2\pi}\partial\bar\partial\varphi\right)^n}{\omega^n}
\\
&\underalign{\text{(Monge--Amp\`ere方程式の定義)}}{=}
 -\frac{\sqrt{-1}}{2\pi}\partial\bar\partial \log e^{\lambda\varphi+f}\underalign{}{=}
 - \lambda \frac{\sqrt{-1}}{2\pi}\partial\bar\partial \varphi
 -\frac{\sqrt{-1}}{2\pi}\partial\bar\partial f
\end{aligned}
\end{equation}


(1). 仮定より$\{\omega_0\} \in c_1(K_X)$であり, 
$-\Ric(\omega_0) \in -c_1(X)=c_1(K_X)$なので, 
 \[
\{-\Ric(\omega_0) - \omega_0 \}\in c_1(0)
\underalign{\text{($\partial\bar\partial$補題 by Lem. \ref{lem-ddc-lem-2})}}{\Rightarrow} 
\exists F \in C^{\infty}(X) \text{ s.t. }
-\Ric(\omega_0) - \omega_0 = \frac{\sqrt{-1}}{2\pi}\partial\bar\partial F 
 \]
 Yauの定理(Theorem~\ref{thm-Yau-theorem})を用いれば
 \begin{align*}
& \varphi \in C^{\infty}(X) \text{ s.t. }
  \left(\omega_0+\frac{\sqrt{-1}}{2\pi}\partial\bar\partial\varphi\right)^n=
  e^{\varphi-F}\omega_0^n
  \\
&\underalign{\eqref{eq-MA-equation-ricci}}{\Rightarrow} 
  \Ric\left(\omega_0+\frac{\sqrt{-1}}{2\pi}\partial\bar\partial\varphi\right)-\Ric (\omega_0)
   = -\frac{\sqrt{-1}}{2\pi}\partial\bar\partial \varphi
 +\frac{\sqrt{-1}}{2\pi}\partial\bar\partial F
 \\
 &\underalign{}{\Rightarrow} 
  \Ric\left(\omega_0+\frac{\sqrt{-1}}{2\pi}\partial\bar\partial\varphi\right)
   =\underbrace{\Ric (\omega_0) +\frac{\sqrt{-1}}{2\pi}\partial\bar\partial F}_{=-\omega_0}
   - \frac{\sqrt{-1}}{2\pi}\partial\bar\partial \varphi
   =-(\omega_0 + \frac{\sqrt{-1}}{2\pi}\partial\bar\partial \varphi)
 \end{align*}
となり言えた.

(2). 
$-\Ric(\omega_0) \in -c_1(X)=0$なので, 
\[
\{\Ric(\omega_0) \}\in c_1(0)
\underalign{\text{($\partial\bar\partial$補題 by Lem. \ref{lem-ddc-lem-2})}}{\Rightarrow} 
\exists F \in C^{\infty}(X) \text{ s.t. }
\Ric(\omega_0)  = \frac{\sqrt{-1}}{2\pi}\partial\bar\partial F 
 \]
 さらに$F+c$と定数和を考えることで, 正規化条件\(\int_X e^F\omega_0^n=\int_X\omega_0^n\)を満たすと仮定する. 
 すると
 Yauの定理(Theorem~\ref{thm-Yau-theorem})を用いれば
 \begin{align*}
& \varphi \in C^{\infty}(X) \text{ s.t. }
  \left(\omega_0+\frac{\sqrt{-1}}{2\pi}\partial\bar\partial\varphi\right)^n=
  e^{F}\omega_0^n
  \\
&\underalign{\eqref{eq-MA-equation-ricci}}{\Rightarrow} 
  \Ric\left(\omega_0+\frac{\sqrt{-1}}{2\pi}\partial\bar\partial\varphi\right)-\Ric (\omega_0)
   =  -\frac{\sqrt{-1}}{2\pi}\partial\bar\partial F
 \\
 &\underalign{}{\Rightarrow} 
  \Ric\left(\omega_0+\frac{\sqrt{-1}}{2\pi}\partial\bar\partial\varphi\right)
   =\underbrace{\Ric (\omega_0) -\frac{\sqrt{-1}}{2\pi}\partial\bar\partial F}_{=0}
   =0
 \end{align*}

(3). K\"ahler form $\theta$で$\theta \in c_1(-K_X)$となるものをとる. 
$\Ric(\omega_0) \in c_1(X)=c_1(-K_X)$なので, 
 \[
\{\Ric(\omega_0) -\theta  \}\in c_1(0)
\underalign{\text{($\partial\bar\partial$補題 by Lem. \ref{lem-ddc-lem-2})}}{\Rightarrow} 
\exists F \in C^{\infty}(X) \text{ s.t. }
\Ric(\omega_0) - \theta= \frac{\sqrt{-1}}{2\pi}\partial\bar\partial F 
 \]
  さらに$F+c$と定数和を考えることで, 正規化条件\(\int_X e^F\omega_0^n=\int_X\omega_0^n\)を満たすと仮定する. 
 Yauの定理(Theorem~\ref{thm-Yau-theorem})を用いれば
 \begin{align*}
& \varphi \in C^{\infty}(X) \text{ s.t. }
  \left(\omega_0+\frac{\sqrt{-1}}{2\pi}\partial\bar\partial\varphi\right)^n=
  e^{F}\omega_0^n
  \\
&\underalign{\eqref{eq-MA-equation-ricci}}{\Rightarrow} 
  \Ric\left(\omega_0+\frac{\sqrt{-1}}{2\pi}\partial\bar\partial\varphi\right)-\Ric (\omega_0)
   = -\frac{\sqrt{-1}}{2\pi}\partial\bar\partial F
 \\
 &\underalign{}{\Rightarrow} 
  \Ric\left(\omega_0+\frac{\sqrt{-1}}{2\pi}\partial\bar\partial\varphi\right)
   =\underbrace{\Ric (\omega_0) -\frac{\sqrt{-1}}{2\pi}\partial\bar\partial F}_{=\theta}
   =\theta
  \underalign{\text{($X$コンパクト)}}{>} 
  \varepsilon \left(\omega_0+\frac{\sqrt{-1}}{2\pi}\partial\bar\partial\varphi\right)
 \end{align*}
最後は$X$がコンパクトを用いてうまく$\varepsilon>0$をとる. 
 \end{proof}
 
 \subsection{Miyaoka-Yau不等式}
 
 
Chen--Ogiue \cite{CO75}による恒等式を思い出そう. 今回は\cite{Hisa24}を参照した.
\begin{prop}[{\cite{CO75, Hisa24}}] 
\label{prop-miyaoka-yau-curvature}
\((X,\omega)\)を\(n\)次元コンパクトK\"ahler多様体とする.
\[
\widetilde{\Ric}(\omega)
\coloneqq 
\Ric(\omega)
-
\frac{R(\omega)}{n}\omega
\quad 
\widetilde{\operatorname{Rm}}_{i\overline j k\overline l}(\omega)
\coloneqq 
\frac{1}{2\pi}\operatorname{R}_{i\overline j k\overline l}(\omega)
-
\frac{R(\omega)}{n(n+1)}
\left(
g_{i\overline j}g_{k\overline l}
+
g_{i\overline l}g_{k\overline j}
\right)
\]
と定義する. 
このとき
\[
\begin{aligned}
&\left\{
2(n+1)c_2(X)-nc_{1}^{2}(X)
\right\}\{\omega\}^{n-2}
\\
&=
\frac{1}{n(n-1)}
\int_X
\left\{
(n+1)\left\lvert \widetilde{\operatorname{Rm}}(\omega)\right\rvert_{\omega}^2
-
(n+2)\left\lvert \widetilde{\Ric}(\omega)\right\rvert_{\omega}^2
\right\}
\omega^n .
\end{aligned}
\]
\end{prop}
証明は\cite[Proposition 4.17]{Iwa-Chern}参照. 

\begin{thm}[Miyaoka--Yau不等式 {\cite{CO75, Yau77}}]
\label{thm-MY-curv}
$X$がK\"ahler--Einstein計量$\omega$を持つ時, 次のMiyaoka--Yau不等式が成り立つ. 
\[
\left\{
2(n+1)c_2(X)-nc_{1}^{2}(X)
\right\}
\{\omega\}^{n-2}
\geq 0
\]

さらに等号が成立するときは$X$の普遍被覆が$\C\mathbb{P}^n$, $\C^n$, $\C^n$のunit ballの三つに限られる.
\end{thm}


\begin{proof}
$\omega$はK\"ahler--Einstein計量なので, 
$\Ric(\omega) = \lambda \omega$となる$\lambda$があり, 
traceをとると$R(\omega) = n\lambda$なので, 
\[
\widetilde{\Ric}(\omega)
\underalign{\text{(定義)}}{=}
\Ric(\omega)
-
\frac{R(\omega)}{n}\omega=
 \lambda \omega- \lambda \omega=0
\]
よってProposition~\ref{prop-miyaoka-yau-curvature}より
\[
\begin{aligned}
&\left\{
2(n+1)c_2(X)-nc_{1}^{2}(X)
\right\}\{\omega\}^{n-2}
\\
&\underalign{\text{(Prop. \ref{prop-miyaoka-yau-curvature})}}{=}
\frac{1}{n(n-1)}
\int_X
\left\{
(n+1)\left\lvert \widetilde{\operatorname{Rm}}(\omega)\right\rvert_{\omega}^2
-
(n+2)\left\lvert \widetilde{\Ric}(\omega)\right\rvert_{\omega}^2
\right\}
\omega^n 
\\
&\quad =
\frac{1}{n(n-1)}
\int_X
\left\{
(n+1)\left\lvert \widetilde{\operatorname{Rm}}(\omega)\right\rvert_{\omega}^2
\right\}
\omega^n 
\geq 0
\end{aligned}
\]

等号成立するときは$\widetilde{\operatorname{Rm}}(\omega)=0$である. 
よって
\[
\frac{1}{2\pi}\operatorname{R}_{i\overline j k\overline l}(\omega)
=
\frac{R(\omega)}{n(n+1)}
\left(
g_{i\overline j}g_{k\overline l}
+
g_{i\overline l}g_{k\overline j}
\right)
\]
となり, 定数holomorphic sectional curvatureをもつ. 古典的な結果\footnote{かの有名な微分幾何の本``Kobayashi--Nomizu''によるとHawleyとIgusaが独立に示したとのことである.}によって$X$の普遍被覆が$\C\mathbb{P}^n$, $\C^n$, $\C^n$のunit ballの三つに限られる.
\end{proof}

よって次の定理を得る. 
\begin{thm}[{\cite{Yau77}}]
$K_X$ampleならば
\[
\left\{
2(n+1)c_2(X)-nc_{1}^{2}(X)
\right\}c_1(K_X)^{n-2}
\geq 0
\]
であり, 等号が成立するときは$X$の普遍被覆が$\C^n$のunit ballになる. 
\end{thm}
 \begin{proof}
Corollary~\ref{cor-Yau-thm}よりK\"ahler--Einstein計量が存在するため, Theorem~\ref{thm-MY-curv}より従う. 
 \end{proof}

\begin{thm}[{\cite{Yau77}}]
$c_1(K_X)=0$ならば, 任意のK\"ahler form $\omega$について
\[
c_2(X)
\{\omega\}^{n-2}
\geq 0
\]
であり, 等号が成立するとき, $X$の有限不分岐被覆が複素トーラスになる. 
\end{thm}
 \begin{proof}
Corollary~\ref{cor-Yau-thm}よりK\"ahler--Einstein計量が存在するため, Theorem~\ref{thm-MY-curv}より従う. 
 等号が成立するとき$X$の普遍被覆が$\C^n$になるが, Bieberbachの古典的定理により, 複素トーラスによって被覆される.
 \end{proof}

$-K_X$ampleの時はK\"ahler--Einstein計量を持つとは限らないので次の形でしか書けない. 
\begin{thm}[{=Theorem~\ref{thm-MY-curv}}]
$-K_X$ampleかつK\"ahler--Einstein計量を持てば
\[
\left\{
2(n+1)c_2(X)-nc_{1}^{2}(X)
\right\}c_1(X)^{n-2}
\geq 0
\]
であり, 等号が成立するときは$X \cong \mathbb{C}\mathbb{P}^n$となる.
\end{thm}

実際Miyaoka--Yauが成り立たないFanoの例が\cite{GKP22}により知られている. 
\cite{Bando, SW16, GKP22, DGP24}より次がわかっている. 
%\footnote{もしかすると$X$がsmoothな場合はこの結果よりも前に知られていたかもしれない. なお\cite{Bando, SW16, GKP22, DGP24}はKLT多様体で示している.}

\begin{thm}[{\cite{Bando, SW16, GKP22, DGP24}}]
$-K_X$ampleかつ, K-semistableならば
\[
\left\{
2(n+1)c_2(X)-nc_{1}^{2}(X)
\right\}c_1(X)^{n-2}
\geq 0
\]
であり, 等号が成立するときは$X \cong \mathbb{C}\mathbb{P}^n$となる.
\end{thm}


\subsection{Miyaoka-Yau不等式の歴史}
  
  例えば曲面の場合のMiyaoka--Yau不等式は
\[
3c_2 \geq c_1(X)^2
\]
と同値である. この形はBogomolov-Miyaoka-Yauの不等式と呼ばれる.
歴史的には以下の通り.
\begin{itemize}
\item Theorem~\ref{thm-MY-curv}はChen--Ogiue \cite{CO75}で示された. なので, 個人的にこの不等式を述べる際には\cite{CO75}も引用すべきな気がする. 
\item \cite{Miy77}が曲面の場合に$3c_2 \geq c_1(X)^2$を示した. この場合は$K_X$ bigでもよい. 
\item \cite{Yau77}でYauがK\"ahler--Einstein計量の存在を示して, $K_X$ ampleの場合のMiyaoka-Yau不等式を示した. $c_1(K_X)=0$の場合もYauの定理である. 
\item KLTの拡張としては$c_1(K_X)=0$の場合は\cite{GKP16, LT18}により示された. maximal-quasi-\'etale coverをつかう. $K_X$ ampleの場合は\cite{GKPT19a, GKPT19b}により示された. Higgs束などを使う.
\item $-K_X$ ampleかつ, K-semistableの場合は, smoothな場合は実質的には\cite{Bando}が初めらしいこれは久本さんに教えてもらった. \cite{SW16}も近いことをやっていて引用すべきな気もする. \cite{Tian92}によるcanonical extension sheafを使った別証明もある. 
\item $-K_X$ ampleかつ, K-semistable KLT多様体は\cite{GKP22, DGP24}による. Tianのcanonical extension sheafをつかう. 
\end{itemize}

 \newpage
 
 \section{Yauの定理の応用2 Enoki's Generic nefness theoremとChern類の不等式}
 
  \subsection{Enoki's generic nefness theorem}
  
 \begin{cor}
\label{cor-Yau-thm-2}
$(X, \omega_0)$をコンパクトK\"ahler多様体とする. 
\begin{enumerate}
\item $-K_X$ nefの時, 任意の$\varepsilon>0$について$\{ \omega\}=\{\omega_0\}$かつ$\Ric(\omega) \geq - \varepsilon \omega$となるK\"ahler計量が存在する.
\item $K_X$ nefの時, 任意の$\varepsilon>0$について$\{ \omega\}=\{\omega_0\}$かつ$\Ric(\omega) \leq \varepsilon \omega_0$なるK\"ahler計量が存在する
\end{enumerate}
\end{cor}
\begin{proof}

 (1). 
$\varepsilon>0$を固定する. 
$-K_X$がnefなので, ある$-K_X$の計量$h_{\varepsilon}$で
\[
\frac{\sqrt{-1}}{2 \pi}
\underbrace{F_{h_{\varepsilon}}}_{\text{Chern 曲率}}
\geq - \varepsilon \omega_0
\quad
\text{and}
\quad
\left\{
\frac{\sqrt{-1}}{2 \pi}F_{h_{\varepsilon}}
\right\}=c_1(-K_X)
\]
となるものが存在する. 
$\Ric(\omega_0) \in c_1(X)=c_1(-K_X)$なので, 
 \[
\{\Ric(\omega_0) - \frac{\sqrt{-1}}{2 \pi}F_{h_{\varepsilon}} \}\in c_1(0)
\underalign{\text{($\partial\bar\partial$補題 by Lem. \ref{lem-ddc-lem-2})}}{\Rightarrow} 
\exists G_{\varepsilon} \in C^{\infty}(X) \text{ s.t. }
\Ric(\omega_0) -  \frac{\sqrt{-1}}{2 \pi}F_{h_{\varepsilon}}  = \frac{\sqrt{-1}}{2\pi}\partial\bar\partial G_{\varepsilon}
 \]
 Yauの定理(Theorem~\ref{thm-Yau-theorem})を用いれば
 \begin{align*}
& \varphi \in C^{\infty}(X) \text{ s.t. }
  \left(\omega_0+\frac{\sqrt{-1}}{2\pi}\partial\bar\partial\varphi\right)^n=
  e^{\varepsilon\varphi+G_{\varepsilon}}\omega_0^n
  \\
&\underalign{\eqref{eq-MA-equation-ricci}}{\Rightarrow} 
  \Ric\left(\omega_0+\frac{\sqrt{-1}}{2\pi}\partial\bar\partial\varphi\right)-\Ric (\omega_0)
   = -\varepsilon\frac{\sqrt{-1}}{2\pi}\partial\bar\partial \varphi
 -\frac{\sqrt{-1}}{2\pi}\partial\bar\partial G_{\varepsilon}
 \\
 &\underalign{}{\Rightarrow} 
  \Ric\left(\omega_0+\frac{\sqrt{-1}}{2\pi}\partial\bar\partial\varphi\right)
   =\underbrace{\Ric (\omega_0) -\frac{\sqrt{-1}}{2\pi}\partial\bar\partial G_{\varepsilon}}_{=\frac{\sqrt{-1}}{2 \pi}F_{h_{\varepsilon}} \geq - \varepsilon \omega_0}
   - \varepsilon\frac{\sqrt{-1}}{2\pi}\partial\bar\partial \varphi
  \geq -\varepsilon(\omega_0 + \frac{\sqrt{-1}}{2\pi}\partial\bar\partial \varphi)
 \end{align*}

 (2). 
$\varepsilon>0$を固定する. 
$K_X$がnefなので, ある$K_X$の計量$h_{\varepsilon}$で
\[
\frac{\sqrt{-1}}{2 \pi}
\underbrace{F_{h_{\varepsilon}}}_{\text{Chern 曲率}}
\geq - \varepsilon \omega_0
\quad
\text{and}
\quad
\left\{
\frac{\sqrt{-1}}{2 \pi}F_{h_{\varepsilon}}
\right\}=c_1(K_X)
\]
となるものが存在する. 
$\Ric(\omega_0) \in c_1(X)=c_1(-K_X)$なので, 
 \[
\{-\Ric(\omega_0) - \frac{\sqrt{-1}}{2 \pi}F_{h_{\varepsilon}} \}\in c_1(0)
\underalign{\text{($\partial\bar\partial$補題 by Lem. \ref{lem-ddc-lem-2})}}{\Rightarrow} 
\exists G_{\varepsilon} \in C^{\infty}(X) \text{ s.t. }
-\Ric(\omega_0) -  \frac{\sqrt{-1}}{2 \pi}F_{h_{\varepsilon}}  = \frac{\sqrt{-1}}{2\pi}\partial\bar\partial G_{\varepsilon}
 \]
  さらに$G_\varepsilon+c$と定数和を考えることで, 正規化条件\(\int_X e^{-G_\varepsilon}\omega_0^n=\int_X\omega_0^n\)を満たすと仮定する. 
 Yauの定理(Theorem~\ref{thm-Yau-theorem})を用いれば
 \begin{align*}
& \varphi \in C^{\infty}(X) \text{ s.t. }
  \left(\omega_0+\frac{\sqrt{-1}}{2\pi}\partial\bar\partial\varphi\right)^n=
  e^{-G_{\varepsilon}}\omega_0^n
  \\
&\underalign{\eqref{eq-MA-equation-ricci}}{\Rightarrow} 
  \Ric\left(\omega_0+\frac{\sqrt{-1}}{2\pi}\partial\bar\partial\varphi\right)-\Ric (\omega_0)
   = 
 \frac{\sqrt{-1}}{2\pi}\partial\bar\partial G_{\varepsilon}
 \\
 &\underalign{}{\Rightarrow} 
  \Ric\left(\omega_0+\frac{\sqrt{-1}}{2\pi}\partial\bar\partial\varphi\right)
   =\underbrace{\Ric (\omega_0) +\frac{\sqrt{-1}}{2\pi}\partial\bar\partial G_{\varepsilon}}_{=-\frac{\sqrt{-1}}{2 \pi}F_{h_{\varepsilon}} \leq  \varepsilon \omega_0}
  \leq \varepsilon \omega_0 
 \end{align*}
\end{proof}


 
 %%%%%%%%%%%%%
 \begin{comment}

\xr{2026/05/08 M.I. せっかくだしなんか追加しましょう}

 \xr{2026/05/09 以降の議論を整理する.多分$\sqrt{-1}\Theta(T_{X},h)\wedge \omega_0^{n-1}$をRicを使って簡単にしないといけないと思う. }
\begin{lem}
 \label{lem-sharp-kahler}
 $\omega, \omega_0$をK\"ahler 計量とする. 
 $T_X$
$\End(T_X)$値$(n, n)$-formの式として次が成り立つ.
$$
n\cdot \sqrt{-1}\Theta(T_{X},h)\wedge \omega^{n-1}
= \sharp^{\omega}\Ric(\omega) \cdot \omega^{n} 
$$
\end{lem}
\begin{proof}
\begin{align*}
n\cdot \sqrt{-1}\Theta(T_{X},h)\wedge \omega^{n-1}
&\underalign{\text{公式}}{=}  
\operatorname{tr}_{\omega}(\Theta(T_{\widehat{X}},h)) \omega^{n} \\
&\underalign{\text{curvature tensor}}{=}  
\frac{1}{2 \pi} g^{k \bar{l}}R^{j}_{i k \bar{l}} \frac{\partial}{\partial z^{j}}\otimes dz^i\cdot \omega^{n}\\
&\underalign{\text{Riemann curvature}}{=}  
\frac{1}{2 \pi} g^{j \bar{p}}g^{k \bar{l}}R_{i \bar{p} k \bar{l}} \frac{\partial}{\partial z^{j}}\otimes dz^i\cdot \omega^{n}\\
&\underalign{\text{Ricci}}{=} 
\frac{1}{2 \pi} g^{j \bar{p}}R_{i \bar{p}} \frac{\partial}{\partial z^{j}}\otimes dz^i \cdot \omega^{n}\\
&\underalign{\text{$\sharp$とRicciの定義}}{=} 
\sharp^{\omega}\Ric(\omega) \cdot \omega^{n} \\
\end{align*}
\end{proof}

 \end{comment}
 %%%%%%%%%%%%%%%%%%%%
 
 
 

 \begin{lem}
 \label{lem-sharp-kahler}
 $\omega$をK\"ahler計量とする.
 ある$\lambda\in \R$とsemipositive form $u$があって
 \[
 \Ric(\omega)=\lambda \omega \pm u
 \]
 となると仮定する. 
 
 $\omega$によって$T_X$に誘導されるHermite計量を$h$とするとき
 \[
 n\frac{\sqrt{-1}}{2 \pi} F_{h}\wedge \omega^{n-1}
= \lambda Id_{T_X}\omega^{n}
\pm \sharp^{\omega}u \omega^{n}
\quad
\text{in $\\End(T_X)$-値$(n,n)$-form}
 \]
 \end{lem}
\begin{proof}
Lemma~\ref{lem-sharp-formula}の公式
 \[
 n\cdot \frac{\sqrt{-1}}{2 \pi} F_{h}\wedge \omega^{n-1}
= \sharp^{\omega}\Ric(\omega) \cdot \omega^{n} 
 \]
 を用いると
 \begin{align*}
  n\frac{\sqrt{-1}}{2 \pi} F_{h}\wedge \omega^{n-1}
  &=\sharp^{\omega}\Ric(\omega) \cdot \omega^{n} 
\underalign{\text{(代入)}}{=}
  \sharp^{\omega}(\lambda \omega \pm u)\cdot \omega^{n}
\underalign{}{=} 
    \lambda Id_{T_X}\omega^{n}
\pm \sharp^{\omega}u \omega^{n}
 \end{align*}


\end{proof}


 
 \begin{thm}[Enoki's generic nefness theorem {\cite{Eno88, Cao13}}]
 \label{thm-Enoki-genericneff}
$X$をコンパクトK\"ahler多様体とする. 
\begin{enumerate}
\item $-K_X$ nefの時, 任意のquotient torsion-free shaef
$T_X \twoheadrightarrow \mathcal{Q}$と任意のK\"ahler form$\omega_0$について, $c_1(\mathcal{Q}) \{ \omega_0\}^{n-1} \geq 0$である. 
\item $K_X$ nefの時, 任意のquotient torsion-free shaef
$\Omega_X \twoheadrightarrow \mathcal{Q}$と任意のK\"ahler form $\omega_0$について, $c_1(\mathcal{Q}) \{ \omega_0\}^{n-1} \geq 0$である. 
\end{enumerate}
 \end{thm}
 \begin{proof}
 (1). $\varepsilon>0$を固定する.
 するとCorollary~\ref{cor-Yau-thm-2}よりK\"ahler form $\omega$と
 あるsemipositive form $u$があって
 \begin{equation}
 \label{eq-antinef-kahler}
 \Ric(\omega)
 =- \varepsilon \omega + u
 \quad 
\{\omega\}=\{\omega_0\}
 \end{equation}
 となるものがある. 

 $T_X \twoheadrightarrow \mathcal{Q}$をrank $r$ torsion-free quotient sheafとし
 $W$をlocally free locus, $Q=\mathcal{Q}\vert_{W}$を$W$上のベクトル束とする. 
 $b$を$T_X \twoheadrightarrow \mathcal{Q}$と$T_X$の計量に関するsecond fundamental formとする. (\cite[Section 11]{Dem12}や\cite[Subsection 5.7]{Iwa-Chern}参照)
 これはあるベクトル値を持つ$(1,0)$-formである. 
 
すると次のようになる. 
\begin{align*}
& c_1(\mathcal{Q}) \cdot \{\omega_0\}^{n-1}
\underalign{\eqref{eq-antinef-kahler}}{=}
c_1(\mathcal{Q}) \cdot \{\omega\}^{n-1} 
\underalign{\text{(Chern-Weil)} }{=}
\int_{W} \frac{\sqrt{-1}}{2\pi} \operatorname{tr}F_{Q} \wedge \omega^{n-1}
\\
&\underalign{\text{\cite[Chapter 11]{Dem12}}}{=} 
\int_{W} 
 \underbrace{
 \frac{\sqrt{-1}}{2\pi} \operatorname{tr}\!\Bigl(\mathrm{pr}_Q
\bigl(F_{T_X,h}\vert_Q\bigr)\Bigr)\wedge \omega^{n-1}
}_{=\frac{1}{n}\operatorname{tr}(\sharp^{\omega}(- \varepsilon \omega + u))\omega^n \text{by Lemma~\ref{lem-sharp-kahler}}}
+
\underbrace{\int_{W} \frac{\sqrt{-1}}{2\pi} \operatorname{tr}\!\bigl(b\wedge b^{*}\wedge \omega^{n-1}\bigr)}_{\geq 0 \text{ by \cite[Chapter 11]{Dem12}}}
 \\
&\underalign{\text{(Lem. \ref{lem-sharp-kahler})}}{\geq}
\frac{\varepsilon}{n} \int_{W} 
\underbrace{- \operatorname{tr}\!\Bigl(\mathrm{pr}_Q\bigl(Id_{T_X})\vert_Q\bigr)
}_{\geq - \operatorname{tr}(Id_{Q}) \text{ (by 射影の性質) }}
\Bigr)\omega^{n}
 +
 \underbrace{\frac{ 1}{n}\int_{W}
 \operatorname{tr}\!\Bigl(\mathrm{pr}_Q(\sharp^{\omega}u)\vert_{Q} \Bigr)\omega^{n} }_{\geq 0}\\
 &\underalign{}{\geq}
- \frac{\varepsilon}{n} \int_{W} 
 \operatorname{tr}\!\bigl(Id_{Q}\bigr)\omega^{n} 
\underalign{\eqref{eq-antinef-kahler}}{=}
- \varepsilon  \left( \frac{ \rank (\mathcal{Q})}{n} \int_{X}  \omega_{0}^{n}\right)
%=- \rank (Q) \frac{\varepsilon}{n} c_1(X)^n
\end{align*}
よって$\frac{ \rank (\mathcal{Q})}{n} \int_{X}  \omega_{0}^{n}$は$\varepsilon$によらないので, $\varepsilon \to 0$とすれば$c_1(\mathcal{Q}) \cdot \{\omega_0\}^{n-1} \geq 0$をえる.

(2). 
双対をとって, 任意のrank $r$ torsion-free subsheaf $\mathcal{F} \subset T_X $について$ c_1(\mathcal{F}) \cdot \{\omega_0\}^{n-1} \leq 0$を示せば良い. 
$\varepsilon>0$を固定する.
Corollary~\ref{cor-Yau-thm-2}よりあるsemipositive form $u$があって
 \begin{equation}
\label{eq-nef-kahler}
 \Ric(\omega)
 =\varepsilon \omega_0 - u
 \quad 
\{\omega\}=\{\omega_0\}
 \end{equation}
 となるものがある. 

 $\mathcal{F} \subset T_X $をrank $r$ torsion-free subsheafとし
 $W$をlocally free locus, $F=\mathcal{F}\vert_{W}$を$W$上のベクトル束とする. \footnote{下の式変形の\cite[Chapter 11]{Dem12}の部分は微妙な点がある.  $T_X/\mathcal{F}$がtorsionを持つ場合は"$=$"ではなく"$\leq$"が正しい. ただ今slopeの不等式の評価しかしていないので, 問題はない. 面倒ならsaturationをとって$\mathcal{F}$が$ T_X$内でsaturatedであるとして良い.}
 $b$を$\mathcal{F} \subset T_X $と$T_X$の計量に関するsecond fundamental formとする. 
これはあるベクトル値を持つ$(1,0)$-formである. 
すると次のようになる. 
\begin{align*}
& c_1(\mathcal{F}) \cdot \{\omega_0\}^{n-1} 
\underalign{\eqref{eq-nef-kahler}}{=} c_1(\mathcal{F}) \cdot \{\omega\}^{n-1}
\underalign{\text{(Chern-Weil)}}{=}
\int_{W} \frac{\sqrt{-1}}{2\pi} \operatorname{tr} F_{F} \wedge \omega^{n-1} 
\\
&\underalign{\text{\cite[Chapter 11]{Dem12}}}{=} 
\int_{W} 
 \underbrace{
 \frac{\sqrt{-1}}{2\pi} \operatorname{tr}\!\Bigl(\mathrm{pr}_F
\bigl(F_{T_X,h}\vert_F\bigr)\Bigr)\wedge \omega^{n-1} 
}_{=\frac{1}{n}\operatorname{tr}(\sharp^{\omega}(\varepsilon \omega_0 - u))\omega^n
 \text{ by Lemma~\ref{lem-sharp-kahler}}}
+
\underbrace{\int_{W} \frac{\sqrt{-1}}{2\pi} \operatorname{tr}\!\bigl(-b\wedge b^{*}\wedge \omega^{n-1}\bigr)}_{\leq 0 \text{ by \cite[Chapter 11]{Dem12}}}
 \\
&\underalign{\text{(Lem. \ref{lem-sharp-kahler})}}{\leq}
\frac{\varepsilon}{n} \int_{W} 
\underbrace{\operatorname{tr}\!\Bigl(\mathrm{pr}_F\bigl(\sharp^{\omega}\omega_0)\vert_F\bigr)
\Bigr)\omega^{n}
}_{\leq  \operatorname{tr}(\sharp^{\omega}\omega_0) \text{(by 射影の性質)}}
 +
 \underbrace{-\frac{ 1}{n}\int_{W}
 \operatorname{tr}\!\Bigl(\mathrm{pr}_F(\sharp^{\omega}u)\vert_{F} \Bigr)\omega^{n} }_{\leq0}\\
 &\underalign{}{\leq}
\frac{\varepsilon}{n} \int_{W} 
\underbrace{
\operatorname{tr}\!\bigl(\sharp^{\omega}\omega_0\bigr)
\omega^{n}
}_{=n \omega_0 \wedge \omega^{n-1}}
\underalign{\eqref{eq-nef-kahler}}{=} 
\varepsilon \int_{X} \omega_{0}^{n}
\end{align*}
 よって$\int_{X} \omega_{0}^{n}$は$\varepsilon$によらないので, $\varepsilon \to 0$とすれば$ c_1(\mathcal{F}) \cdot \{\omega_0\}^{n-1} \leq 0$をえる.
 \end{proof}
 
 
 \subsection{Chern類の不等式}
Generic nefness theorem (Theorem~\ref{thm-Enoki-genericneff})を使えば\cite{Miy87}の計算を用いて次を示すことができる. 
\begin{thm}[{\cite{Miy87, Cao13, Ou23, IM22, IMM24}}]
$(X, \omega)$をコンパクトK\"ahler多様体とする
\begin{itemize}
\item \(-K_X\)がnefならば, 
\[
c_2(T_X) (c_1(-K_X) + \varepsilon\{ \omega\})^{n-2}\geq 0
\]
が十分に小さい$\varepsilon >0$で成り立つ. 
\item \(K_X\)がnefならば
\[
\bigl(3c_2(\Omega_X)-c_1(X)^2\bigr)(c_1(K_X) + \varepsilon \{\omega\})^{n-2} \geq 0
\]
が十分に小さい$\varepsilon >0$で成り立つ. 
\end{itemize}
\end{thm}
 この辺りは\cite[Subsection 6.2, 6.3]{Iwa-Chern}を参照のこと. 
 次は\cite[Subsection 6.2, 6.3]{Iwa-Chern}の方が強いことが言えている. 

 
 \subsection{Generic nefness theoremの歴史}
まずProjectiveとコンパクトK\"ahlerによってだいぶ違う. 
Projectiveの場合は, \cite{Miy87}が初だと思う. 
\begin{thm}[{\cite[Corollary 6.4]{Miy87}}]
$X$を$n$次元のnormal projective varietyとし$H$をample Cartier divisorとする. $\pi_{*} \colon \widetilde{X} \to X$を特異点解消とする.
$X$がnon-uniruledならば, 任意のtorsion-free quotient sheaf $\pi_{*}\Omega_{\widetilde{X}} \twoheadrightarrow \mathcal{Q}$について$c_1(\mathcal{Q})\cdot H^{n-1} \geq 0$である.
\end{thm}
証明の概略は次のとおり. 
\begin{enumerate}
\item $c_1(\mathcal{Q})\cdot H^{n-1} <0$となる$\pi_{*}\Omega_{\widetilde{X}} \twoheadrightarrow \mathcal{Q}$があるとする. すると$\mathcal{F} \subset T_{\widetilde{X}}$でminimal slopeが0より大きいものが作れる. 
\item この$\mathcal{F}$がalgebraic foliation with rationally connected leavesになる. 特に$\mathcal{F}$からfibration $f \colon \widetilde{X} \dashrightarrow Y$でgeneral fiberのZariski closureがRationally connectedであるものが誘導される. 
\item しかしこれはrational curveがいっぱいあるので, $X$がnon-uniruledに矛盾する. 
\end{enumerate}

この方向性の研究は歴史的には以下の通り.
\begin{itemize}
\item \cite{BDPP13}より$X$がsmoothならばnon-uniruledと$K_X$ psefが同値なので, 上はcanonical singularityで$K_X$ psefで成り立つ. 
\item Foliationに関しては\cite{BM16, KST07}など色々あるが, 最終的にmovable classまで拡張したのは\cite{CP19}だと思う. \cite{Cla16}のサーベイに詳しい. \cite{Lazicnote}も読みやすいと思う. 
\item \cite{CKT21}よりKLTでも成り立つことがわかった
\item $-K_X$nefのgeneric nefnessは\cite{Ou23}で完全解決. 
\end{itemize}

 コンパクトK\"ahlerの場合は上のfoliationの議論が成り立つかどうか最近まで不明だった. (2025年に解決) ただそれまで時間がかなりあったので, コンパクトK\"ahlerの場合は独自の進化があった
\begin{itemize}
\item generic nefnessは\cite{Eno88}がおそらく初め. その後\cite{Cao13}が出たのだが, 多分結果が被ったようだと思う(\cite[Remark 2]{Cao13}参照)
\item \cite{Gue16}でKLTやLCで拡張された. 多分SLCまでできたと思う
\item  \cite{Ou25}でコンパクトK\"ahlerの場合に上のfoliationの議論が成り立つことがわかった. ブレイクスルーと言っても良い. 
\item $-K_X$がnefの場合は\cite{Ou25}を使って\cite{IJZ25, MWWZ25, FGSW26}で解決. 意外と最近のことである. 
\end{itemize}

Chern類の不等式に関しては\cite{Miy87, Cao13, Ou23, IM22}でsmoothな場合, singularな場合は\cite{IMM24}である.
等号成立もきっちりとしらべたのが\cite{Cao13, Ou23, IM22, IMM24}である. \cite{IMM24}で全て解決という感じである. (今読んでも\cite{Miy87, Cao13, Ou23}で学ぶことは多い)

最後に"generic nef"についてだが, これは多分\cite{Pet12}による. 元々\cite{Miy87,Eno88}ではgeneric semipositiveという言い方だったが, \cite{Ou23}などでgeneric nefに変わったと思う. 

 
 \newpage
  \section{Yauの定理の応用3 Enoki's  stability theorem}
 
 \subsection{Enoki's stability theorem}
\begin{cor}
\label{cor-Yau-thm-3}
$(X, \omega_0)$をコンパクトK\"ahler多様体とする. 
$K_X$ nefの時, 任意の$\varepsilon>0$について, $\omega \in c_1(K_X) + \varepsilon\{\omega_0\}$かつ$\Ric(\omega) = -\omega + \varepsilon \omega_0$なるK\"ahler計量が存在する. 
\end{cor}

\begin{proof}
$\varepsilon>0$を固定する. 
$K_X$がnefなので, ある$K_X$の計量$h_{\varepsilon}$で
\[
\frac{\sqrt{-1}}{2 \pi}
\underbrace{F_{h_{\varepsilon}}}_{\text{Chern 曲率}}
\geq - \frac{\varepsilon}{2} \omega_0
\quad
\text{and}
\quad
\left\{
\frac{\sqrt{-1}}{2 \pi}F_{h_{\varepsilon}}
\right\}=c_1(K_X)
\]
となるものが存在する. そこで
\[
\chi_{\varepsilon}\coloneqq  \frac{\sqrt{-1}}{2 \pi}F_{h_{\varepsilon}}
+ \varepsilon \omega_0
\]
とする. 
$\chi_{\varepsilon}$はK\"ahler formで$\{ \chi_{\varepsilon}\}=c_1(K_X) + \varepsilon\{\omega_0\}$である. 
$\Ric(\chi_{\varepsilon}) \in c_1(X)=c_1(-K_X)$なので, 
 \[
\{-\Ric(\chi_{\varepsilon}) - \chi_{\varepsilon} + \varepsilon \omega_0\}\in c_1(0)
\underalign{\text{($\partial\bar\partial$補題 by Lem. \ref{lem-ddc-lem-2})}}{\Rightarrow} 
\exists G_{\varepsilon} \in C^{\infty}(X) \text{ s.t. }
-\Ric(\chi_{\varepsilon}) - \chi_{\varepsilon} + \varepsilon \omega_0 = \frac{\sqrt{-1}}{2\pi}\partial\bar\partial G_{\varepsilon}
 \]
 Yauの定理(Theorem~\ref{thm-Yau-theorem})を用いれば
 \begin{align*}
& \varphi \in C^{\infty}(X) \text{ s.t. }
  (\chi_{\varepsilon}+\frac{\sqrt{-1}}{2\pi}\partial\bar\partial\varphi)^n=
  e^{\varphi -G_{\varepsilon}}\chi_{\varepsilon}^n
  \\
&\underalign{\eqref{eq-MA-equation-ricci}}{\Rightarrow} 
  \Ric(\chi_{\varepsilon}+\frac{\sqrt{-1}}{2\pi}\partial\bar\partial\varphi)-\Ric (\chi_{\varepsilon})
   = 
-    \frac{\sqrt{-1}}{2\pi}\partial\bar\partial \varphi 
 +\frac{\sqrt{-1}}{2\pi}\partial\bar\partial G_{\varepsilon}
 \\
 &\underalign{}{\Rightarrow} 
  \Ric(\chi_{\varepsilon}+\frac{\sqrt{-1}}{2\pi}\partial\bar\partial\varphi)
   =
   \Ric (\chi_{\varepsilon})
   -    \frac{\sqrt{-1}}{2\pi}\partial\bar\partial \varphi 
    +\frac{\sqrt{-1}}{2\pi}\partial\bar\partial 
    G_{\varepsilon}
=
 - \left(\chi_{\varepsilon} +\frac{\sqrt{-1}}{2\pi}\partial\bar\partial \varphi \right)
    + \varepsilon \omega_0
  \end{align*}
 よって$\omega\coloneqq \chi_{\varepsilon} +\frac{\sqrt{-1}}{2\pi}\partial\bar\partial \varphi$とおけば良い. 
\end{proof}

 \begin{thm}[{\cite[Theorem 1.1]{Eno88}}]
 \label{thm-Enoki-stability}
$(X, \omega_0)$をコンパクトK\"ahler多様体とする. 
$K_X$ nefの時, 任意のrank $r$ torsion freesheaf $\mathcal{F} \subset T_X$について
\begin{equation}
\label{eq-Enoki-stability}
\begin{aligned}
&\frac{1}{r}
c_1(\mathcal{F}) \cdot \left(c_1(K_X) + \varepsilon \{ \omega_0\} \right)^{n-1}
\\
&\leq
\frac{1}{n}
c_1(-K_X) \cdot \left(c_1(K_X) + \varepsilon \{ \omega_0\} \right)^{n-1}
+
\varepsilon
\left( \frac{1}{r} - \frac{1}{n}\right)
\cdot \{ \omega_0\}
 \left(c_1(K_X) + \varepsilon \{ \omega_0\} \right)^{n-1}
\end{aligned}
\end{equation}
 \end{thm}

証明の前にすぐにわかるcorollaryを書いておく.
\begin{cor}[{\cite[Corollary 1.2, 1.3]{Eno88}}]
$(X, \omega_0)$を$n$次元コンパクトK\"ahler多様体とする. $K_X$ nefとし, 数値的小平次元を
\[
\nu \coloneqq  \max \left\{ 0 \leq k \leq n\;\middle\vert\; c_1(K_X)^k \not \equiv 0 \right\}.
\]
と定義する. この時次が成り立つ.
\begin{enumerate}
\item $\nu=0$の時, $T_X$は$\{\omega_0\}^{n-1}$-semistableである. 
\item  $\nu=n$の時, つまり$K_X$がbigの時, $T_X$は$K_X^{n-1}$-semistableである. 
\item $0<\nu<n$の時, 任意のrank $r$ torsion-free sheaf $\mathcal{F} \subset T_X$について
\[
c_1(\mathcal{F}) \cdot c_1(K_X)^{\nu} \cdot \{ \omega_0\}^{n-\nu-1}
\leq 0
\]
である. さらにこの等号が成立する時
\[
\nu  c_1(\mathcal{F}) \cdot c_1(K_X)^{\nu-1} \cdot \{ \omega_0\}^{n-\nu}
\leq 
(r+\nu - n)
c_1(-K_X) \cdot c_1(K_X)^{\nu-1} \cdot \{ \omega_0\}^{n-\nu}
\]
となる. 
特に$T_X$は$K_X^{\nu}\cdot \{ \omega_0\}^{n-\nu-1}$-semistableである. 
\end{enumerate}
\end{cor}
\begin{proof}
rank $r$ torsion-free sheaf $\mathcal{F} \subset T_X$とする. 
$\nu <n$ならば 
 \eqref{eq-Enoki-stability}の$\varepsilon^{n-\nu-1}$の項のみを展開すると,
\begin{equation}
\label{eq-Enoki-stability-1}
\binom{n-1}{\nu}\varepsilon^{n-\nu-1}
\cdot 
\frac{1}{r}
c_1(\mathcal{F}) 
c_1(K_X)^{\nu}
\{ \omega_0\}^{n-\nu-1}
\leq
\binom{n-1}{\nu}\varepsilon^{n-\nu-1}
\cdot 
\frac{1}{n}
c_1(-K_X) c_1(K_X)^{\nu}
\{ \omega_0\}^{n-\nu-1}
+O(\varepsilon^{n-\nu})
\end{equation}
となる. ここで$O(\varepsilon^{n-\nu})$はランダウの記号とする. 
$\nu =n$ならば\eqref{eq-Enoki-stability}の$\varepsilon$を含まない定数項を見れば
\begin{equation}
\label{eq-Enoki-stability-2}
\frac{1}{r}
c_1(\mathcal{F}) 
c_1(K_X)^{n-1}
\leq
\frac{1}{n}
c_1(-K_X) c_1(K_X)^{n-1}
+O(\varepsilon)
\end{equation}
よって(1), (2)は\eqref{eq-Enoki-stability-1}, \eqref{eq-Enoki-stability-2}を$\varepsilon \to 0$すればすぐに従う.(厳密には$\varepsilon$の何乗かで割ってから極限を取る.)

(3)については,  一つ目の主張は$c_1(-K_X) c_1(K_X)^{\nu}
\{ \omega_0\}^{n-\nu-1}=0$なので, \eqref{eq-Enoki-stability-1}からすぐに従う. 
二つ目の主張について, \eqref{eq-Enoki-stability}の$\varepsilon^{n-\nu}$の項を展開すれば, 
\begin{align*}
&
\binom{n-1}{\nu-1}\varepsilon^{n-\nu}
\frac{1}{r}
 c_1(\mathcal{F}) \cdot c_1(K_X)^{\nu-1} \cdot \{ \omega_0\}^{n-\nu}
\\
&\leq
\binom{n-1}{\nu-1}\varepsilon^{n-\nu}\frac{1}{n}
c_1(-K_X) \cdot c_1(K_X)^{\nu-1} \cdot \{ \omega_0\}^{n-\nu}
+
\binom{n-1}{\nu}\varepsilon^{n-\nu}
\left( \frac{1}{r} - \frac{1}{n}\right)
\cdot
c_1(K_X)^{\nu} \cdot \{ \omega_0\}^{n-\nu}
 +O(\varepsilon^{n-\nu+1})
\end{align*}
よって$\varepsilon^{n-\nu}$の項を整理すれば,　欲しい式を得る. 
\end{proof}
\begin{proof}[Proof of Theorem~\ref{thm-Enoki-stability}]

$\varepsilon>0$を固定する.
Corollary~\ref{cor-Yau-thm-3}より
 \begin{equation}
\label{eq-nef-kahler-Enoki}
\Ric(\omega) = -\omega + \varepsilon \omega_0
 \quad 
\{\omega \} = c_1(K_X) + \varepsilon\{\omega_0\}
 \end{equation}
 となるものがある. 
 

 $\mathcal{F} \subset T_X $をrank $r$ torsion-free subsheafとし
 $W$をlocally free locus, $F=\mathcal{F}\vert_{W}$を$W$上のベクトル束とする. 
\footnote{ここも前と同じく微妙な問題があって, 下の式変形の\cite[Chapter 11]{Dem12}の部分$T_X/\mathcal{F}$がtorsionを持つ場合は"$=$"ではなく"$\leq$"が正しい.}
 $b$を$\mathcal{F} \subset T_X $と$T_X$の計量に関するsecond fundamental formとする. 
これはあるベクトル値を持つ$(1,0)$-formである. 
すると次のようになる. 
\begin{align*}
& c_1(\mathcal{F}) \cdot 
\left(c_1(K_X) + \varepsilon \{ \omega_0\} \right)^{n-1}
\underalign{\eqref{eq-nef-kahler-Enoki}}{=} c_1(\mathcal{F}) \cdot \{\omega\}^{n-1}
\underalign{\text{Chern-Weil}}{=}
\int_{W} \frac{\sqrt{-1}}{2\pi} \operatorname{tr} F_{F} \wedge \omega^{n-1} 
\\
&\underalign{\text{\cite[Chapter 11]{Dem12}}}{=} 
\int_{W} 
 \underbrace{
 \frac{\sqrt{-1}}{2\pi} \operatorname{tr}\!\Bigl(\mathrm{pr}_F
\bigl(F_{T_X,h}\vert_F\bigr)\Bigr)\wedge \omega^{n-1} 
}_{=\frac{1}{n}\operatorname{tr}(\sharp^{\omega}( -\omega + \varepsilon \omega_0))\omega^n \text{ by Lemma~\ref{lem-sharp-kahler}}}
+
\underbrace{\int_{W} \frac{\sqrt{-1}}{2\pi} \operatorname{tr}\!\bigl(-b\wedge b^{*}\wedge \omega^{n-1}\bigr)}_{\leq 0 \text{ by \cite[Chapter 11]{Dem12}}}
 \\
&\underalign{\text{(Lem. \ref{lem-sharp-kahler})}}{\leq}
-\frac{1}{n} \int_{W} 
\operatorname{tr}\!
\bigl(
Id_{F}
\bigr)\omega^{n}
 +
 \underbrace{\frac{\varepsilon}{n}\int_{W}
 \operatorname{tr}\!\Bigl(\mathrm{pr}_F(\sharp^{\omega}\omega_0))\vert_{F} \Bigr)\omega^{n} }_ {\leq  \operatorname{tr}(\sharp^{\omega}\omega_0) \text{ (by 射影の性質)  }}
 \\
 &\underalign{}{\leq}
-\frac{r}{n} \int_{X} \omega^{n}
+
\frac{\varepsilon}{n}
\underbrace{\int_{X}
 \operatorname{tr}(\sharp^{\omega}\omega_0) \omega^{n} 
 }_{= n \int_{X} \omega_0 \wedge \omega^{n-1}}
\underalign{\eqref{eq-nef-kahler-Enoki}}{=}
  -\frac{r}{n} \left(c_1(K_X) + \varepsilon \{ \omega_0\} \right)^{n}
  +\varepsilon \{\omega_0\} \cdot  \left(c_1(K_X) + \varepsilon \{ \omega_0\} \right)^{n-1}
\end{align*}
よってあとは式変形をすれば欲しい式が得られる. 
\end{proof}





 \begin{thm}[{\cite[Theorem 1.1]{Eno88}}]
 \label{thm-Enoki-stability-tensor}
$(X, \omega_0)$をコンパクトK\"ahler多様体とし, $m \geq 1$となる整数とする. 
$K_X$ nefの時, 任意のrank $r$ torsion-free sheaf $\mathcal{F} \subset T_X^{\otimes m}$について
\begin{equation}
\label{eq-Enoki-stability-tensor}
\begin{aligned}
&\frac{1}{r}
c_1(\mathcal{F}) \cdot \left(c_1(K_X) + \varepsilon \{ \omega_0\} \right)^{n-1}
\\
&\leq
\frac{1}{n^m}
c_1(T_{X}^{\otimes m}) \cdot \left(c_1(K_X) + \varepsilon \{ \omega_0\} \right)^{n-1}
+
\varepsilon
m n^{m-1}\left( \frac{1}{r} - \frac{1}{n^m}\right)
\cdot \{ \omega_0\}
 \left(c_1(K_X) + \varepsilon \{ \omega_0\} \right)^{n-1}
\end{aligned}
\end{equation}
 \end{thm}
 \begin{proof}
 $\varepsilon>0$を固定する.
Corollary~\ref{cor-Yau-thm-3}より
 \begin{equation}
\label{eq-nef-kahler-Enoki-tensor}
\Ric(\omega) = -\omega + \varepsilon \omega_0
 \quad 
\{\omega \} = c_1(K_X) + \varepsilon\{\omega_0\}
 \end{equation}
 となるものがある. 
 
 $h$を$\omega$によって誘導される$T_X$のHermite計量とし, 
 $T_X^{\otimes m}$に$h^{m}$という計量を入れる. 
 するとtensorと曲率の公式から, 
 \[
F_{T_X^{\otimes m},h^m}
=
\sum_{j=1}^m
\Id^{\otimes(j-1)}
\otimes
F_{T_X,h}
\otimes
\Id^{\otimes(m-j)}
 \]
 となる. 特に
 以上より
 \begin{equation}
\label{eq-multi-tensor}
  \frac{\sqrt{-1}}{2\pi} 
  F_{T_X^{\otimes m},h^m} \wedge \omega^{n-1}
  \underalign{\text{(Lem. \ref{lem-sharp-kahler} \& }\eqref{eq-nef-kahler-Enoki-tensor})}{=}
 -\frac{m}{n}\Id_{T_X^{\otimes m}}
 \cdot \omega^{n}
+
\frac{\varepsilon }{n}
\sum_{j=1}^m
\left(
\Id^{\otimes(j-1)}
\otimes \sharp^\omega\omega_0
\otimes
\Id^{\otimes(m-j)} 
\right)
 \cdot \omega^{n}
 \end{equation}

 
 $\mathcal{F} \subset T_X^{\otimes m}$をrank $r$ torsion-free subsheafとし
 $W$をlocally free locus, $F=\mathcal{F}\vert_{W}$を$W$上のベクトル束とする. 
\footnote{ここも前と同じく微妙な問題があって, 下の式変形の\cite[Chapter 11]{Dem12}の部分$T_X/\mathcal{F}$がtorsionを持つ場合は"$=$"ではなく"$\leq$"が正しい.}
 $b$を$\mathcal{F} \subset T_X^{\otimes m}$と$T_X^{\otimes m}$の計量に関するsecond fundamental formとする. 
これはあるベクトル値を持つ$(1,0)$-formである. 
すると次のようになる. 
\begin{align*}
& c_1(\mathcal{F}) \cdot 
\left(c_1(K_X) + \varepsilon \{ \omega_0\} \right)^{n-1}
\underalign{\eqref{eq-nef-kahler-Enoki}}{=} c_1(\mathcal{F}) \cdot \{\omega\}^{n-1}
\underalign{\text{Chern-Weil}}{=}
\int_{W} \frac{\sqrt{-1}}{2\pi} \operatorname{tr} F_{F} \wedge \omega^{n-1} 
\\
&\underalign{\text{\cite[Chapter 11]{Dem12}}}{=} 
\int_{W} 
 \frac{\sqrt{-1}}{2\pi} \operatorname{tr}\!\Bigl(\mathrm{pr}_F
\bigl(F_{T_{X}^{\otimes m},h^{m}}\vert_F\bigr)\Bigr)\wedge \omega^{n-1} 
+
\underbrace{\int_{W} \frac{\sqrt{-1}}{2\pi} \operatorname{tr}\!\bigl(-b\wedge b^{*}\wedge \omega^{n-1}\bigr)}_{\leq 0 \text{ by \cite[Chapter 11]{Dem12}}}
 \\
&\underalign{\eqref{eq-multi-tensor}}{\leq}
 -\frac{m}{n}
 \int_{W} 
\operatorname{tr}\!
\bigl(
Id_{F}
\bigr)\omega^{n}
 +
 \frac{\varepsilon }{n}
\sum_{j=1}^m
\int_{W}
 \operatorname{tr}\!\Bigl(\mathrm{pr}_F(
 \Id^{\otimes(j-1)}
\otimes \sharp^\omega\omega_0
\otimes
\Id^{\otimes(m-j)})\vert_{F} \Bigr)\omega^{n}
 \\
 &\underalign{}{\leq}
-\frac{mr}{n} \int_{X} \omega^{n}
+
\frac{\varepsilon}{n}
\sum_{j=1}^m
n^{m-1}
\underbrace{\int_{X}
 \operatorname{tr}(\sharp^{\omega}\omega_0) \omega^{n} 
 }_{= n \int_{X} \omega_0 \wedge \omega^{n-1}}
 \\
&\underalign{\eqref{eq-nef-kahler-Enoki}}{=}
  -\frac{mr}{n} \left(c_1(K_X) + \varepsilon \{ \omega_0\} \right)^{n}
  +\varepsilon m n^{m-1} \{\omega_0\} \cdot  \left(c_1(K_X) + \varepsilon \{ \omega_0\} \right)^{n-1}
\end{align*}
よってあとは
$\frac{1}{n^m}
c_1(T_{X}^{\otimes m}) = \frac{m}{n}c_1(-K_X)$を用いて
式変形をすれば欲しい式が得られる. 

 
 \end{proof}


\begin{cor}[{cf. \cite[Subsection 5]{CPa16}}]
$(X, \omega_0)$を2次元以上のコンパクトK\"ahler多様体とし, $m \geq 2$となる整数とする. torsion-free sheafによる完全系列
\[
0 \longrightarrow L \longrightarrow \left(\Omega_{X}^{1}\right)^{\otimes m} \longrightarrow \mathcal{Q} \longrightarrow 0
\]
であって, $\rk L =1$となるものを考える. $K_X$がnefかつ$c_1(\mathcal{Q})=0$が成り立つならば, $K_X \equiv 0$である. 
\end{cor}
実は\cite{CPa16}により弱い仮定$c_1(\mathcal{Q}) \cdot \{ \omega_{0}\}^{n-1}=0$でもいけるが, それにはより強い主張を示す必要がある.


\begin{proof}
$\rank \mathcal{Q}=n^m-1$であり,
\(
\mathcal{Q}^{\vee}\subset T_X^{\otimes m},
\) かつ\(
c_1(\mathcal{Q}^{\vee})=-c_1(\mathcal{Q})=0
\)
である. Theorem~\ref{thm-Enoki-stability-tensor}を$\mathcal F=\mathcal{Q}^{\vee}$に適用すると,
\begin{align*}
0
&\leq
-\frac{m}{n}c_1(K_X)\cdot\left(c_1(K_X)+\varepsilon\{\omega_0\}\right)^{n-1}
+\frac{\varepsilon m}{n(n^m-1)}
\{\omega_0\}\cdot\left(c_1(K_X)+\varepsilon\{\omega_0\}\right)^{n-1}.
\end{align*}
したがって
\[
0\geq
\left((n^m-1)c_1(K_X)-\varepsilon\{\omega_0\}\right)
\cdot
\left(c_1(K_X)+\varepsilon\{\omega_0\}\right)^{n-1}.
\]

$c_1(K_X)\not\equiv0$と仮定し, $\nu\coloneqq\nu(K_X)\geq1$とする. すると\begin{align*}
&\left((n^m-1)c_1(K_X)-\varepsilon\{\omega_0\}\right)
\cdot
\left(c_1(K_X)+\varepsilon\{\omega_0\}\right)^{n-1} =
\binom{n-1}{\nu-1}
\frac{\nu n^m-n}{\nu}
\varepsilon^{n-\nu}
c_1(K_X)^\nu\cdot\{\omega_0\}^{n-\nu}
+O(\varepsilon^{n-\nu+1}).
\end{align*}
$m\geq2$, $n\geq2$, $\nu\geq1$より
\(
\frac{\nu n^m-n}{\nu}>0.
\)である. 
したがって十分小さい$\varepsilon>0$に対して上式の右辺は正となり矛盾する. よって$\nu(K_X)=0$であり, $K_X$はnefなので
\(
K_X\equiv0
\)となる. 
\end{proof}






 \subsection{Enoki's stability theoremの歴史}
 時系列としてはこんな感じである. 
\begin{itemize}
\item K\"ahler--Einstein計量があれば$T_X$がpolystabilityはわかるので(\cite[Section 5]{Iwa-Chern}参照),  $K_X$がampleや$c_1(K_X)=0$の場合は, polystabilityに関しては\cite{Yau77}が初めだと思う.
\item $K_X$がnefの場合は\cite{Eno88}がおそらく初め. この論文, $X$がterminalかつ$K_X$がnef + effectiveと書ける場合をやっている. 時代を先取りした研究である. 
\item \cite{Gue16}でKLTやLCで拡張された. 多分SLCまでできたと思う. これは\cite{Eno88}の方法を精密化したものだと思う. 
\item \cite{CPa16}のlog pairの応用もある. これも\cite{Eno88}の方法を精密化したものだと思う. 
\item 実はLCだと$c_1(K_X)=0$でもpolystabilityは成り立たない. \cite{BFPT24}による. 
\end{itemize}
 
 
 \newpage
 
  \section{Yauの定理の応用4 $-K_X$がnefな多様体の基本群}
 以下Riemannian幾何学, とりわけ比較定理に関する定理を用いる. 
 これに関しては\cite{Zhu97}のサーベイを引用する
 \subsection{Fanoは単連結}
\begin{thm}
Fano多様体は単連結である. 
つまり$-K_X$ ampleならば単連結である. 
\end{thm}
\begin{proof}
Corollary~\ref{cor-Yau-thm}から, ある$\varepsilon>0$があって, $\Ric(\omega) \geq \varepsilon \omega$なるK\"ahler計量が存在する.
ここでMyersの定理を思い出す.

\begin{thm}[Myersの定理 {\cite[Theorem 3.4]{Zhu97}}]
$n$次元完備Riemannian多様体$(M,g, \omega)$について, $H>0$となる定数があって, 
$\Ric(\omega) \geq (n-1)H \omega$が成り立つならば, 
$\diam(M) \leq \frac{\pi}{\sqrt{H}}$であり, 基本群$\pi_1(M)$は有限である. 
特に$M$はコンパクトである. 
\end{thm}

よって$X$の基本群は有限である. 
そこで普遍被覆$\pi \colon \widetilde{X} \to X$をとる. 基本群の有限性より
$\pi$はfinite etaleであり, $\widetilde{X}$もまたFanoである. 
$\pi$の次数を$d$とすると
\[
d \cdot \chi(X, \mathcal{O}_X) =  \chi( \widetilde{X}, \mathcal{O}_{\widetilde{X}})
\]
である. 今$X$も$\widetilde{X}$もFanoなので, 
$ \chi( \widetilde{X}, \mathcal{O}_{\widetilde{X}}) = \chi(X, \mathcal{O}_X) =1$である.
よって$d=1$であり, $\pi$は恒等写像となる. 特に$X$の基本群は自明である.  
\end{proof}


 \subsection{$-K_X$がsemipositiveならば基本群はalmost Abelian}
 
 \subsubsection{群論の用語の定義}
 以下群に関して次の用語を定義しておく.
 \begin{defn}
 $G$を群とする.
 \begin{enumerate}
 \item $G$がalmost nilpotent (or virtually nilpotent)であるとは, 有限指数の正規部分群$H \triangleleft G$があって, $H$がnilpotentであること. 
 \item $G$がalmost Abelian (or virtually Abelian)であるとは, 有限指数の正規部分群$H \triangleleft G$があって, $H$がAbelianであること. 
 \end{enumerate}
 \end{defn}
 
 almost nilpotentに関しては有限生成であれば次の有用な定理がある. 
 そのために次の準備をする. 
 
\begin{defn}[{\cite[Definition 3.9]{Zhu97}}]
$G$を有限生成群とし,
\(
G=\langle g_1,g_2,\ldots,g_k\rangle
\)
とする. 
\begin{itemize}
\item generatorsの集合
\(
S=\{g_1,g_2,\ldots,g_k\}
\)
とする. $S \cup S^{-1}$に関する$g \in G$のword lengthを
\[
\lvert g\rvert_{S}\coloneqq 
\min\{ m \geq 0 \mid g=t_1 \cdots t_m \quad t_i \in S \cup S^{-1} \}
\]
と定める. 
\item $r$-neighborhoodを
\(
U_{S}(r)
\coloneqq 
\{g\in G\mid \lvert g\rvert_{S} \leq r\}
\)
によって定義する.
\item $G$がpolynomial growthであるとは, generatorsの集合$S$と正の数$s$が存在して,十分大きな$r$に対して
\(
\lvert U_{S}(r)\rvert\leq r^s
\)
が成立することとする. 
\end{itemize}
\end{defn}
\begin{thm}[Gromov 81 {\cite[Theorem 3.10]{Zhu97}}]
\label{thm-Gromov-thm}
有限生成群$G$がpolynomial growthをもつことは, almost nilpotentと同値. 
\end{thm}

\subsubsection{ $-K_X$がsemipositiveの基本群}

\begin{thm}[{\cite{DPS96}}]
コンパクトK\"ahler多様体$X$について, $-K_X$がsemipositiveならば基本群はalmost Abelianである. 
\end{thm}
この証明は二つの定理に基づく. 
\begin{thm}
\label{thm-Milnor-thm}
$X$をコンパクトK\"ahler多様体とする. 
$-K_X$ semipositiveならば, 
基本群はalmost nilpotentである. 
\end{thm}

\begin{thm}
\label{thm-Campana-Special}
$X$をコンパクトK\"ahler多様体とする. 
$-K_X$がnefかつ基本群はalmost nilpotentならば, 
基本群はalmost Abelianである. 
\end{thm}

以下これを示していく. 

\subsubsection{Proof of Theorem~\ref{thm-Milnor-thm}}
こっちはMilnorの定理(cf. \cite[Theorem 3.11]{Zhu97})と呼ばれているものである. 
Bishop-Gromovの不等式の応用である. 
\begin{proof}
$-K_X$ semipositiveならば, $\Ric(\omega) \geq 0$となるK\"ahler計量が存在する.
これはCorollary~\ref{cor-Yau-thm-2}と同じである. 実際
$-K_X$がsemipositiveなので, ある$-K_X$の計量$h$で
\[
\frac{\sqrt{-1}}{2 \pi}
\underbrace{F_{h}}_{\text{Chern 曲率}}
\geq 0
\quad
\text{and}
\quad
\left\{
\frac{\sqrt{-1}}{2 \pi}F_{h}
\right\}=c_1(-K_X)
\]
となるものが存在する. 
$\Ric(\omega_0) \in c_1(X)=c_1(-K_X)$なので, 
 \[
\{\Ric(\omega_0) - \frac{\sqrt{-1}}{2 \pi}F_{h} \}\in c_1(0)
\underalign{\text{($\partial\bar\partial$補題 by Lem. \ref{lem-ddc-lem-2})}}{\Rightarrow} 
\exists G_{\varepsilon} \in C^{\infty}(X) \text{ s.t. }
\Ric(\omega_0) -  \frac{\sqrt{-1}}{2 \pi}F_{h}  = \frac{\sqrt{-1}}{2\pi}\partial\bar\partial G
 \]
  さらに$G+c$と定数和を考えることで, 正規化条件\(\int_X e^G\omega_0^n=\int_X\omega_0^n\)を満たすと仮定する. 
 Yauの定理(Theorem~\ref{thm-Yau-theorem})を用いれば
 \begin{align*}
& \varphi \in C^{\infty}(X) \text{ s.t. }
  \left(\omega_0+\frac{\sqrt{-1}}{2\pi}\partial\bar\partial\varphi\right)^n=
  e^{G}\omega_0^n
  \\
&\underalign{\eqref{eq-MA-equation-ricci}}{\Rightarrow} 
  \Ric\left(\omega_0+\frac{\sqrt{-1}}{2\pi}\partial\bar\partial\varphi\right)-\Ric (\omega_0)
   = 
 -\frac{\sqrt{-1}}{2\pi}\partial\bar\partial G
 \\
 &\underalign{}{\Rightarrow} 
  \Ric\left(\omega_0+\frac{\sqrt{-1}}{2\pi}\partial\bar\partial\varphi\right)
   =\Ric (\omega_0) -\frac{\sqrt{-1}}{2\pi}\partial\bar\partial G
   =
\frac{\sqrt{-1}}{2 \pi}F_{h} \geq 0
 \end{align*}
 
以降の議論は\cite[Theorem 3.11]{Zhu97}に同じである. 
普遍被覆$ \pi \colon \widetilde{X} \to X$を考え, $\widetilde{X}$上のRiemannian計量$\widetilde{g}$を, $\Ric(\omega) \geq 0$となるK\"ahler計量から誘導する. 
 $n = \dim_{\C}X$とすると, 
$\widetilde{X}$は複素$n$次元の複素多様体であるとともに, 実$2n$次元完備Riemann多様体で$\Ric(\widetilde{g}) \geq 0$となる. 

ここでBishop-Gromovの定理(で下からの曲率が0のもの)を思い出そう.
\begin{thm}[Bishop--Gromovの定理 {\cite[Corollary 3.3]{Zhu97}}]
実$2n$次元完備Riemann多様体$(M,g, \omega)$が, $\Ric(\omega) \geq 0$を満たすならば, 
任意の$p\in M$と$0<r\leq R$に対して
\begin{equation}
\label{eq-Bishop-Gromov}
\frac{\vol(B_p(r))}{\vol(B_p(R))}
\geq
\left(\frac{r}{R}\right)^{2n}.
\end{equation}
が成り立つ. ここで
\(
B_p(r)\coloneqq \{x\in M\mid d_g(p,x)<r\}
\)
は$p$を中心とする半径$r$の測地球であり,
\(
\vol(B_p(r))
\coloneqq 
\int_{B_p(r)}dV_g
=
\int_{B_p(r)}\frac{\omega^n}{n!}
\)
は$g$に関するRiemann体積である.

特に次が成り立つ. 
\[
\vol(B_p(r))\leq \frac{\pi^n}{n!} r^{2n}
\]
\end{thm}

示すべきことは$\pi_1(X)$がalmost nilpotentであること, 同値的にTheorem~\ref{thm-Gromov-thm}からpolynomial growthであることである. 
$G\coloneqq \pi_1(X, p)$とし, $\widetilde{X}$に$G$が作用しているとして良い. 
$\pi(\widetilde{p})=p$となる$\widetilde{p} \in \widetilde{X}$を一つ固定しておく. 

$G$は有限生成なので, $S =\{ g_1, \ldots, g_k\}$とgeneratorを一つ固定する. 
このとき,各$g_i \in \pi_1(X, p)$は$p \in X$内のloopの同値類, つまり
\[
\sigma_i \colon [0, l_i] \to X
\quad
\sigma_i(0)=\sigma_i(l_i)=p
\]
とみなせる. 
今$\sigma_i$のliftは
\[
\widetilde{\sigma_i} \colon [0, l_i] \to \widetilde{X}
\quad
\widetilde{\sigma_i}(0)=\widetilde{p}
\quad
\widetilde{\sigma_i}(l_i)=g_i(\widetilde{p})
\]
となる. $G$は$\widetilde{X}$にisometryに作用しているので
\begin{equation}
\label{eq-length-evaluate}
d_{\widetilde{g}}(\widetilde{p}, g_i \widetilde{p}) \leq 
l_i \leq \max\{l_1,\ldots,l_k\} \eqqcolon l
\end{equation}
である. 
次に$\varepsilon>0$で相異なる任意の$g_1,g_2\in G$に対して,
\begin{equation}
\label{eq-disjoint-h}
h_1(B_{\widetilde p}(\varepsilon))
\cap
h_2(B_{\widetilde p}(\varepsilon))
=
\varnothing
\end{equation}
となるようなものをとる. 
これは$\pi \colon \widetilde{X} \to X$がcovering mapなので, 
$p \in V$という近傍で$\pi^{-1}(V)=\sqcup_{g \in G}V_g$となるものが存在する. 
そこで$\widetilde{p} \in B_{\widetilde p}(\varepsilon) \subset V_g$となるように$\varepsilon$を小さく取れば良い. 

この時長さ$r$となる$h \in U(r) \subset G$について
\begin{equation}
\label{eq-h-length}
h(B_{\widetilde p}(\varepsilon))) \subset B_{\widetilde p}(rl+\varepsilon)
\end{equation}
である. 実際$z \in h(B_{\widetilde p}(\varepsilon))
$について, 
$z=h \cdot x$となる$x \in B_{\widetilde p}(\varepsilon)$をとると
\[
d_{\widetilde{g}}(\widetilde{p}, z) 
\leq 
d_{\widetilde{g}}(\widetilde{p}, h\cdot \widetilde{p}) 
+ d_{\widetilde{g}}(h\cdot \widetilde{p}, h \cdot x) 
=
d_{\widetilde{g}}(\widetilde{p}, h\cdot \widetilde{p}) 
+ d_{\widetilde{g}}(\widetilde{p}, x) 
\underalign{\eqref{eq-length-evaluate}}{\leq}  lr + \varepsilon
\]
となるためである. 以上をまとめると
\[
\begin{aligned}
\lvert U(r)\rvert\cdot\operatorname{vol}(B_{\widetilde p}(\varepsilon))
\underalign{\eqref{eq-disjoint-h}}{=}
\operatorname{vol}
\left(
\bigsqcup_{h\in U(r)}
h(B_{\widetilde p}(\varepsilon))
\right)
\underalign{\eqref{eq-h-length}}{\leq}
\operatorname{vol}(B_{\widetilde p}(rl+\varepsilon))
\end{aligned}
\]
よってBishop Gromovの定理を使って
\[
\begin{aligned}
\lvert U(r)\rvert
&\leq
\frac{\operatorname{vol}(B_{\widetilde p}(rl+\varepsilon))}
{\operatorname{vol}(B_{\widetilde p}(\varepsilon))}
\underalign{\eqref{eq-Bishop-Gromov}}{\leq}
\frac{(rl+\varepsilon)^{2n}}{\varepsilon^{2n}}
=
cr^{2n}
\end{aligned}
\]
となるように定数$c$を$r$によらず取れる. 
よって定義から$G$はpolynomial growthである. 
\end{proof}


\subsubsection{Proof of Theorem~\ref{thm-Campana-Special}}
この辺りは\cite[Section 7]{Iwa-Cam}にまとめている. 
要するにCampanaのSpecial多様体の話である. 
鍵となるのはSpecial多様体などでも現れる"Albanese mapが全射かつファイバー連結"であることである. 

\begin{proof}
finite etale coverをとって, $\pi_1(X)$をnilpotentかつtorsion-freeであるとして良い. 
ここで$-K_X$がnefならば, Campanaの意味でspecialである.\footnote{\cite[Remark 2.22]{Iwa-Cam}参照. Theorem~\ref{thm-Enoki-genericneff}のCorollaryである.}よって\cite[Subsection 5.2]{Iwa-Cam}からAlbanese mapが全射かつファイバー連結
である. (下のRemark~\ref{rem-cao-nefanti}の別証明も参照)

このあとは\cite{Cam95a}の議論による. \cite[Theorem 3.23]{Voi}を参考にした. 
(\cite[Subsection 7.2]{Iwa-Cam}そのまま写した)
$\alpha \colon X \to A$をAlbanese mapとする. 
 \begin{itemize}
 \item \(\alpha_* \colon \pi_1(X) \to \pi_1(A)\)は全射になる. これは\( \phi \colon X \to A\)がファイバー連結なので (Lemma~\ref{lem-Voi-1.2}(1)参照)
\item  \( \alpha_* \colon H_1(X, \mathbb{Z}) \to H_1(A, \mathbb{Z}) \)は全射かつ, finite kernelである. これはAlbanese mapの性質から.
\item \( \alpha^* \colon H^2(A, \mathbb{Q}) \to H^2(X, \mathbb{Q}) \)単射である. (Lemma~\ref{lem-Voi-1.2}(2)参照)
\end{itemize}

\( H \coloneqq  \pi_1(A) =\Z^{2k}\)とおく.
\[
n\colon H^2(H, \mathbb{Q}) =H^2(\Z^{2k}, \Q)\to H^2(A, \mathbb{Q})
\]
という(全)単射が存在することを示す. 
\(E_H \to B_H = E_H / H\)を$H$の分類空間とする.\footnote{\url{http://pantodon.jp/index.rb?body=BG}参照. 任意の位相空間$X$上の主$H$束$E \to X$について, ある連続写像$h \colon X \to B_H$が存在して, $h^{*}E_H$と$E$は同相になる. 例えば$H=\Z$なら$\R\to S^1$がそれに当たる. } 
ここで\(E_H\)は可縮である. 
一方普遍被覆
\(
\C^k\to A=\C^k/H
\)
も主$H$束であり, $\C^k$は可縮である. 
したがって$\C^k\to A$自身が$H$の普遍主束のmodelとなり, 
$A$は$B_H$のmodelとみなすことができる. 
特にclassifying map
\(
u\colon A\to B_H
\)
はhomotopy equivalenceである. 

そこで$n$を以下で定める.\footnote{ここは\cite[Theorem 3.23]{Voi}に合わせた. \cite[Theorem 3.23]{Voi}ではもっと一般の状況を考えている.}
\[
n\coloneqq u^{*} \colon H^2(B_H, \mathbb{Q}) \to H^2(A, \mathbb{Q}).
\]
$u$はhomotopy equivalenceなので,
これは全単射である. 
また
\(
H^2(H,\Q)\coloneqq H^2(B_H,\Q)
\)
であり, $H\simeq\Z^{2k}$なので,
\[
H^2(H,\Q)
=
H^2(\Z^{2k},\Q)
\overset{n}{\simeq}
H^2(A,\Q)
\]
を得る.
%%%%%%%%%%%%%%
\begin{comment}

\[ n\colon H^2(H, \mathbb{Q}) =H^2(\Z^{2k}, \Q)\to H^2(A, \mathbb{Q}) \]
という(全)単射が存在することを示す. 
  \( E_H \to B_H = E_H / H \)を$H$の分類空間とする.\footnote{ \url{http://pantodon.jp/index.rb?body=BG}参照. 任意の位相空間$X$上の主$H$束$E \to X$について, ある連続写像$h \colon X \to B_H$が存在して, $h^{*}E_H$と$E$は同相になる. 例えば$H=\Z$なら$\R\to S^1$がそれに当たる. } 
   \( E_H \)は可縮である. 
すると普遍被覆\( \C^k \to A \)が主$H$束なので, 
$A \underset{\text{homeo}}{\cong}(\C^k\times E_H) / H $である.
そこで 
\[ u \colon A \underset{\text{(上の同型)}}{\to} (\C^k\times E_H)/H 
\underset{\text{(射影)}}{\to} E_H / H = B_H
\]
とおく. 
 そして$n$を以下で定める.\footnote{ここは\cite[Theorem 3.23]{Voi}に合わせた. \cite[Theorem 3.23]{Voi}ではもっと一般の状況を考えている.}
\[
n\coloneqq u^{*} \colon H^2(B_H, \mathbb{Q}) \to H^2(A, \mathbb{Q}) = H^2((\C^k\times E_H)/H, \mathbb{Q})
\]
今の場合$ E_H =\R^{2k}\to B_H=(S^1)^{2k}$であるので, $A \to B_H$は$\C^{k}/\Z^{2k} \to (S^1)^{2k}$のように複素torusの複素構造を忘れて, $(S^1)^{2k}$とみなした写像(恒等写像)に等しい.
よって全単射となる.
\end{comment}
%%%%%%%%%%%%%%
よって\( \alpha_* \colon H_1(X, \mathbb{Q}) \to H_1(A, \mathbb{Q}) \)は単射なので上と合わせて
 \[  
 H_2(X, \mathbb{Q}) \underset{\alpha_*}{\to} 
  H_2(A, \mathbb{Q}) \underset{n=u^{*}}{\to} 
 H_2(H, \mathbb{Q}) 
 \]
 は全射の合成で全射になる. 
 \( G = \pi_1(X) \)とすると以下の可換図式が成り立つ:
\[
\begin{array}{ccc}
H_2(X, \mathbb{Q}) & \underset{\alpha_*}{\longrightarrow}& H_2(A, \mathbb{Q}) \\
\downarrow & & \downarrow \\
H_2(G, \mathbb{Q}) & \underset{\alpha_*}{\longrightarrow} & H_2(H, \mathbb{Q})
\end{array}
\]
よって%$ \pi_1(X) =G \twoheadrightarrow H　=\pi_1(A)$によって
  \(H_2(G, \mathbb{Q}) \to H_2(H, \mathbb{Q}) \)も全射である. 
以上より
\begin{itemize}
\item  \( \alpha_* \colon H_1(G, \mathbb{Q}) \to H_1(H, \mathbb{Q}) \)同型. 
\item   \( \alpha_* \colon H_2(G, \mathbb{Q}) \to H_2(H, \mathbb{Q}) \)全射. 
\end{itemize}
であるので,$\pi_1(X),  \pi_1(A)$ともにnilpotentであるから,  Lemma~\ref{lem-Voi-3.24}より
\[G=\pi_1(X) \to \pi_1(A)=H
\]
は有限のkernelを持つ.(全射であることは前から言えてる)
\end{proof}

使った定理をまとめておく. 
\begin{lem}[{\cite[Lemma 1.2]{Voi}}]
\label{lem-Voi-1.2}
 \(  f \colon X \to Y \ \)はproper surjectiveで次を満たす. 
\begin{enumerate}
\item \( f_{*} \colon \pi_1(X) \to \pi_1(Y) \)について, $ f_{*}(\pi_1(X)) \subset \pi_1(Y)$はfinite indexを持つ. さらにファイバー連結ならば,  $f_{*}(\pi_1(X)) =  \pi_1(Y)$
\item  \( f^* \colon H^i(Y,\mathbb{Q}) \to H^i(X,\mathbb{Q}) \)はHodge構造における単射である. 
\item  一般ファイバー\( X_s \)が連結かつ, 任意の\( i>0 \)で\( H^0(X_s,\Omega^i_{X_s}) = 0 \)
ならば, \( f^* \colon H^0(Y,\Omega^i_Y) \to H^0(X,\Omega^i_X) \)は\( i \geq 0 \)で同型である. 
\end{enumerate}
\end{lem}

\begin{lem}[{\cite[Lemma 3.24]{Voi}, \cite[Proposition 6.2]{Wang22}}]
\label{lem-Voi-3.24}
%\label{lem:3.24}
\( \alpha \colon G \to H \)を有限生成群の準同型とする. 
\( H_1(G, \mathbb{Q}) \simeq H_1(H, \mathbb{Q}) \)同型かつ, 
\( H_2(G, \mathbb{Q}) \twoheadrightarrow  H_2(H, \mathbb{Q}) \)全射ならば, 
任意の\( n \in \N \)について,  
\[ G/G_n \to H/H_n \]
はfinite kernel and cokernelを持つ. 
%ここで， 群$\Gamma$について, $n$-th lower central seriesを\( \Gamma_n = [\Gamma, \Gamma_{n-1}] \)とする. 
ここで群$\Gamma$について, $n$-th lower central seriesを\( \Gamma_n = [\Gamma, \Gamma_{n-1}] \)とする. 
\end{lem}


\begin{rem}
\label{rem-cao-nefanti}
$-K_X$ nefならばAlbanese mapが全射かつファイバー連結に関しては\cite{Cao13}による別証明もある. こっちも見ておく. 
\begin{proof}
まず$0\neq\tau\in H^0(X,\Omega_X^1)$は$X$上のどの点でも消えないならば, 
Albanese map $\alpha \colon X \to A$がsubmersionであることを示そう.
任意の$x\in X$について
\[
(d\alpha_x)^*\colon T_{\alpha(x)}^*A\longrightarrow T_x^*X
\]
を考える.$A$上のholomorphic $1$-formはtranslation invariantであり,
\[
\alpha^*\colon H^0(A,\Omega_A^1)\xrightarrow{\sim}H^0(X,\Omega_X^1)
\]
である.もし$(d\alpha_x)^*$がinjectiveでなければ,ある$0\neq\eta\in H^0(A,\Omega_A^1)$について$(\alpha^*\eta)(x)=0$となる. 
よって仮定から, 任意の$x\in X$について$(d\alpha_x)^*$はinjectiveであり,
\[
d\alpha_x
\colon \Omega_{X, x}^{1}
\longrightarrow 
 \Omega_{A, \alpha(x)}^{1}
\]はsurjectiveである.ゆえに$\alpha$はsubmersionである.特に$\alpha(X)$は$A$のopen subsetであるが,$X$はcompactなので$\alpha(X)$はclosedでもある.$A$はconnectedであるから$\alpha(X)=A$,すなわち$\alpha$は全射である.

$\alpha$のStein factorizationを
\[
X\xrightarrow{\beta}Y\xrightarrow{\nu}A
\]
とする.ここで$\beta$は連結ファイバーをもち,$\nu$はfinite surjectiveである.
$\nu$が\'etaleでないと仮定する.
branch locusのpurityにより$\nu$はramification divisorをもつ.
その一般点$y$では$d\nu_y$はsurjectiveでない.
$x\in\beta^{-1}(y)$を取ると
\(
d\alpha_x=d\nu_y\circ\,d\beta_x
\)
であるから$d\alpha_x$もsurjectiveでないことになり, $\alpha$がsubmersionであることに矛盾する.
したがって$\nu$はfinite \'etaleである.
よって$Y$もcomplex torusである.
原点を適当に選べば, Albanese mapのuniversal propertyから$\nu$は同型となり, $\alpha$は連結ファイバーをもつ.

よって$0\neq\tau\in H^0(X,\Omega_X^1)$がどの点でも消えないことを示す.\footnote{一般には$-K_X$がnefなだけではこれは言えない. $-K_X$がnefで言えるのは$0\neq\tau\in H^0(X,\Omega_X^1)$がcodimension 2のanalytic setを除いてどの点でも消えないことである. しかしこの場合でも, $\alpha$が全射連結ファイバーを持つのは言える. またcodimension 2のanalytic setを除いてsubmersionも言える. }
$T_X$の$\omega^{n-1}$に関するHarder--Narasimhan filtrationをとったのちJordan filtrationをとって
\[
0=E_0\subset E_1\subset\cdots\subset E_s=T_X
\]
というfiltrationで
$Q_i\coloneqq E_i/E_{i-1}$は$\omega$-stableかつ$\mu_\omega(Q_i)$が$i$に関して非増加となるものが取れる. Theorem~\ref{thm-Enoki-genericneff}より$\mu_\omega(Q_i)\geq0$である.
ここで以下簡単のための各$E_i$はlocally freeであり, subbundleによるfiltrationであると仮定する.\footnote{この仮定の場合は$0\neq\tau\in H^0(X,\Omega_X^1)$がどの点でも消えないことが言える.そうでない場合にcodimension 2のanalytic setの議論が出てくる. }

今,$Q_i\coloneqq E_i/E_{i-1}$は$\omega$-stableなので, 
Hermite-Einstein計量$h_i$があって
\[
\Lambda_{\omega}\left(\frac{\sqrt{-1}}{2 \pi}F_{h_i}\right)
=c_i Id_{Q_i}
\quad \text{here}
\quad
c_i \coloneqq  \frac{n}{\{ \omega\}^{n}}\mu_{\omega}(Q_i)
\]
となるものがある.  (例えば\cite[Proposition 4.7]{Iwa-Chern}参照)
特にTheorem~\ref{thm-Enoki-genericneff}から$c_i \geq 0$である. 

すると下のLemma~\ref{lem-Cao13-2.6}によって, ある$C>0$があって, 
任意の$\varepsilon>0$に対して$T_X$上のHermite計量$h_\varepsilon$で
\begin{equation}
\label{eq-hepsilon-cutvature}
\Lambda_{\omega}\left(\frac{\sqrt{-1}}{2 \pi} F_{h_\varepsilon}\right)
\geq-\varepsilon\operatorname{id}_{T_X}
\quad \text{and}
\quad
\lVert\sqrt{-1} F_{h_\epsilon}\rVert_{L^\infty}
\leq
C
\end{equation}
を満たすものを取ることができる. ここで次を思い出す
\begin{tcolorbox}[mybox]
$\rho \colon X \to \R, \theta \colon \R \to \R$について
$\chi(x)\coloneqq \theta(\rho(x))$とすると
\[
i\partial\bar\partial\chi=\theta''(\rho)\, i\partial\rho\wedge \bar\partial\rho + \theta'(\rho)\, i\partial\bar\partial\rho.
\]

特に
$u$を実数値関数として
\[
\partial\bar\partial \log u=
\frac{\partial\bar\partial u}{u}
- \frac{\partial u \wedge \bar\partial u}{u^2}
\]

特に$(E, H)$ベクトル束とHermite計量のpairと
$\tau \in H^0(X, E)$について

\begin{equation}
\label{eq-Bochner-formula}
\begin{aligned}
\sqrt{-1}\partial \bar\partial \log (\lvert\tau\rvert^{2}_{H})
&\underalign{\text{上の公式}}{=}
\frac{\sqrt{-1}\partial \overline{\partial } \lvert\tau\rvert^{2}_{H}}{\lvert\tau\rvert^{2}_{H}}
- \frac{\sqrt{-1} \partial \lvert\tau\rvert^{2}_{H} \wedge \overline{\partial } \lvert\tau\rvert^{2}_{H}}{(\lvert\tau\rvert^{2}_{H})^2}
\\
&\underalign{}{=}\frac{\sqrt{-1} \{ D'\tau, D'\tau\}  }{(\lvert\tau\rvert^{2}_{H})}
- \frac{\sqrt{-1}  \{ D'\tau, \tau\} \wedge  \{ \tau, D'\tau\}}{(\lvert\tau\rvert^{2}_{H})^2}
- \frac{\{\tau, \sqrt{-1}F_{H} \tau\}}{(\lvert\tau\rvert^{2}_{H})}
\end{aligned}
\end{equation}

ここで$\{\cdot,\cdot\}$は次のcanonical sesquilinear pairingを表す：
\[
C^\infty\!\bigl(\wedge^p T_X^*\otimes E\bigr)
\times
C^\infty\!\bigl(\wedge^q T_X^*\otimes E\bigr)
\longrightarrow
C^\infty\!\bigl(\wedge^{p+q}T_X^*\bigr).
\quad
\{ \alpha \otimes u, \beta \otimes v\}\coloneqq \alpha \wedge \overline{\beta} H(u, v)
\]
上の等式は
\[
\partial \lvert\tau\rvert^{2}_{H} = \{ D'\tau, \tau\}
\quad 
\bar\partial \lvert\tau\rvert^{2}_{H} = \{ \tau, D'\tau\}
\quad
\partial \bar\partial \lvert\tau\rvert^{2}_{H}
= \{ D'\tau, D'\tau\} - \{\tau, \Theta_{E, H} \tau\}
\]
から従う. 
\end{tcolorbox}
今$0\neq\tau\in H^0(X,\Omega_X^1)$を取り, $h_\varepsilon^*$をdual metricとする.
$E=\Omega_{X}^{1}, H=h_\varepsilon^*$に上の公式\eqref{eq-Bochner-formula}を用いれば
\begin{equation}
\label{eq-curvature-formula-1}
\begin{aligned}
\sqrt{-1}\partial \bar\partial \log (\lvert\tau\rvert^{2}_{h_\varepsilon^*})
&\underalign{\eqref{eq-Bochner-formula}}{=}
\underbrace{
\frac{\sqrt{-1} \{ D'\tau, D'\tau\}  }{(\lvert\tau\rvert^{2}_{h_\varepsilon^*})}
- \frac{\sqrt{-1}  \{ D'\tau, \tau\} \wedge  \{ \tau, D'\tau\}}{(\lvert\tau\rvert^{2}_{h_\varepsilon^*})^2}
}_{\geq 0 \text{ by local computation}}  
- \frac{\{\tau, \sqrt{-1}F_{h_\varepsilon^*} \tau\}}{(\lvert\tau\rvert^{2}_{h_\varepsilon^*})}
\end{aligned}
\end{equation}
以上よりこれを合わせて
\begin{equation}
\label{eq-curvature-formula-2}
\begin{aligned}
\sqrt{-1}\partial\bar\partial\log\lvert\tau\rvert_{h_\varepsilon^*}^2\wedge\omega^{n-1}
&\underalign{}{=}
\frac{1}{n}
\Lambda_{\omega}\left(\left(\sqrt{-1}\partial\bar\partial\log\lvert\tau\rvert_{h_\varepsilon^*}^2 \right)\right)\cdot\omega^{n}
\\
&\underalign{\eqref{eq-curvature-formula-1}}{\geq}
- \frac{\{\tau, \Lambda_{\omega}\sqrt{-1}F_{h_\varepsilon^*} \tau\}}{n (\lvert\tau\rvert^{2}_{h_\varepsilon^*})}
\cdot\omega^{n}
\underalign{\eqref{eq-hepsilon-cutvature}}{\geq}
-\frac{2 \pi}{n}\varepsilon\omega^n
\end{aligned}
\end{equation}
これを用いてある$x\in X$で$\tau(x)=0$であるとして矛盾を示す. 
\[
T_\varepsilon \coloneqq  \sqrt{-1}\partial \bar\partial \log (\lvert\tau\rvert^{2}_{h_\varepsilon^*})
\]
とする. 
すると\eqref{eq-hepsilon-cutvature}にある$C>0$を用いて
\[
T_\varepsilon \coloneqq  
\sqrt{-1}\partial \bar\partial \log (\lvert\tau\rvert^{2}_{h_\varepsilon^*})
\underalign{\eqref{eq-curvature-formula-1}}{\geq}
- \frac{\{\tau, \sqrt{-1}F_{h_\varepsilon^*} \tau\}}{(\lvert\tau\rvert^{2}_{h_\varepsilon^*})}
\underalign{\eqref{eq-hepsilon-cutvature}}{\geq}
-C\omega
\]
特に$T_\varepsilon+C\omega$はpositivecurrentで$T_\varepsilon$はexactなので, 
そのmass$\int_{X}(T_\varepsilon+C\omega) \wedge \omega^{n-1}$は有界である. 
よってpositivecurrentのweak compactness\footnote{positivecurrentの列のmassが一様有界なら, 部分列がcurrentの意味で収束する. これはBanach-Alaogluの定理から. }
から部分列を取り替えて
\[
T_\varepsilon+C\omega \longrightarrow S
\]
とclosed positivecurrent $S$にcurrentの意味で収束する.
$T\coloneqq S - C\omega$とする. 

今$x\in X$で$\tau(x)=0$なので, 
$\log (\lvert\tau\rvert^{2}_{h_\varepsilon^*})$は$x$でanalytic singularityを持つ. 特にLelong numberは1以上である. 
よって
\[
1 \leq\limsup_{\varepsilon \to 0}\nu(T_\varepsilon+C\omega , x) 
\underalign{\text{(weak compactness)}}{\leq}
\nu(T+C\omega , x) 
\]
である. Lelong数の定義を思い出すと, $\omega$にしか寄らない定数$D>0$があって
\[
\nu(T+C\omega , x) 
= 
D 
\lim_{r \to 0^+}
\frac{1}{r^{2n-2}}
\int_{B_x(r)}
(T+C\omega) \wedge \omega^{n-1}
\]
である. よって$\nu(T+C\omega , x) \geq 1$なので, 定数$D$を取り替えて, 
十分小さな$r>0$について
\begin{equation}
\label{eq-Br-TComega-1}
\int_{B_x(r)}
(T+C\omega) \wedge \omega^{n-1}
\geq D r^{2n-2}
\end{equation}
が成り立つとして良い. 
以上から$r$を十分小さくすれば, 定数$D'>0$があって
\begin{equation}
\label{eq-Br-TComega-2}
\int_{B_x(r)}
T\wedge \omega^{n-1}
=
\int_{B_x(r)}
(T+C\omega)\wedge \omega^{n-1}
-
\int_{B_x(r)}
C\omega^{n}
\underalign{\eqref{eq-Br-TComega-1}}{\geq} D r^{2n-2} + O(r^{2n-1})
\geq D' r^{2n-2}
\end{equation}
 とできる.
今ほとんど全ての$r>0$について
\[
\int_{B_x(r)}(T_{\varepsilon}+C\omega) \wedge \omega^{n-1}
\underset{\varepsilon \to 0}{\longrightarrow}
\int_{B_x(r)}(T+C\omega) \wedge \omega^{n-1}
\Rightarrow 
\int_{B_x(r)}T_{\varepsilon} \wedge \omega^{n-1}
\underset{\varepsilon \to 0}{\longrightarrow}
\int_{B_x(r)}T\wedge \omega^{n-1}
\]
である. 
以上より
\[
\int_{X}T_{\varepsilon} \wedge \omega^{n-1}
=
\underbrace{\int_{B_x(r)}T\wedge \omega^{n-1}}_{\geq D' r^{2n-2} \text{ by \eqref{eq-Br-TComega-2}}}
+
\underbrace{\int_{B_x(r)}(T_\varepsilon - T)\wedge \omega^{n-1}}_{\to 0 \text{ as $\varepsilon \to 0$}}
+
\underbrace{\int_{X \setminus B_x(r)}T_{\varepsilon} \wedge \omega^{n-1}}
_{\geq -\frac{2 \pi}{n}\varepsilon \int_{X} \omega^{n} \text{ by\eqref{eq-curvature-formula-2}}}
\Rightarrow 
\liminf_{\varepsilon \to 0} \int_{X}T_{\varepsilon} \wedge \omega^{n-1}>0
\]
しかし$T_\varepsilon$はexactで$ \log (\lvert\tau\rvert^{2})$はlocally $L^1$なので, Stokesの定理から$\int_{X}T_{\varepsilon} \wedge \omega^{n-1}=0$である.\footnote{$ \log (\lvert\tau\rvert^{2})$はlocally $L^1$なので$\partial \bar\partial \log (\lvert\tau\rvert^{2})$がwell definedなcurrentになり, integral by partsを使えるから.}これは矛盾である. 
\end{proof}

\begin{lem}[{\cite[Lemma 2.6]{Cao13}}]
\label{lem-Cao13-2.6}
$(X,\omega)$をコンパクトK\"ahler多様体とする.
$E$を二つのベクトル束$E_1,E_2$のextension
\[
0\to E_1\to E\to E_2\to 0
\]
とする.$E_1,E_2$上に二つの滑らかな計量$h_1,h_2$が存在して, 各点で
\begin{equation}
\label{eq:lemma26-7}
\Lambda_{\omega}\left(\frac{\sqrt{-1}}{2 \pi}F_{h_1}\right)
\geq c_1\cdot\operatorname{Id}_{E_1}
\quad\text{and}\quad
\Lambda_{\omega}\left(\frac{\sqrt{-1}}{2 \pi}F_{h_2}\right)
\geq c_2\cdot\operatorname{Id}_{E_2}
\end{equation}
が成り立つと仮定する.このとき,任意の$\epsilon>0$に対して,$E$上の滑らかな計量$h_\epsilon$が存在して
\begin{equation}
\label{eq:lemma26-7-2}
\Lambda_{\omega}\left(\frac{\sqrt{-1}}{2 \pi}F_{h_\epsilon}\right)
\geq \bigl(\min(c_1,c_2)-\epsilon\bigr)\cdot\operatorname{Id}_E
\end{equation}
を満たし,さらに$\epsilon$に依存しないある一様な定数$C$に対して
\begin{equation}
\label{eq:lemma26-8}
\lVert\sqrt{-1} F_{h_\epsilon}(E)\rVert_{L^\infty}
\leq
C\cdot
\left(
\lVert \sqrt{-1} F_{h_1}(E_1)\rVert_{L^\infty}
+
\lVert\sqrt{-1} F_{h_2}(E_2)\rVert_{L^\infty}
\right)
\end{equation}
が成り立つ.
\end{lem}

\begin{proof}
$[E]\in H^1(X,\operatorname{Hom}(E_2,E_1))$をextension groupにおいて$E$を表す元とする.\cite[Section 11]{Dem12}から
これは$\operatorname{Hom}(E_2,E_1)$値$(0,1)$-formであるsecond fundamental formの$\beta$の代表元とも見れる. 
そこで$E_1$と$E_2$の別のextensionで$s\in\mathbb{C}^*$について
\(
[E_s]=s\cdot[E]
\)
を満たすものとする.これは$s \beta$を考えることに同じである. 
このとき,これら二つのベクトル束の間にはisomorphismが存在する. 
このisomorphismを
\[
\varphi_s\colon E\to E_s
\]
と書く.\cite[Equation (11.6)]{Dem12}の式と
\eqref{eq:lemma26-7}から, $\lvert s\rvert$が$\epsilon$に比べて十分小さければ, $E_s$上の滑らかな計量$h_s$で
\begin{equation}
\label{eq:lemma26-9}
\Lambda_{\omega}\left(\frac{\sqrt{-1}}{2 \pi}F_{h_s}\right)
\geq
\bigl(\min(c_1,c_2)-\epsilon\bigr)\cdot\operatorname{Id}_{E_s}
\end{equation}
を満たし,さらに一様な定数$C$に対して
\begin{equation}
\label{eq:lemma26-10}
\lVert\sqrt{-1} F_{h_s}(E_s)\rVert_{L^\infty}
\leq
C\cdot
\left(
\lVert\sqrt{-1} F_{h_1}(E_1)\rVert_{L^\infty}
+
\lVert\sqrt{-1} F_{h_2}(E_2)\rVert_{L^\infty}
\right)
\end{equation}
を満たすものを取れる. 
そこで$h=\varphi_s^*(h_s)$を$E$上に誘導されるmetricとする.
このとき,任意の$\alpha\in E$に対して
\begin{equation}
\label{eq:lemma26-11}
\begin{aligned}
\frac{\langle\sqrt{-1} F_h(E)\alpha,\alpha\rangle_h}{\langle\alpha,\alpha\rangle_h}
&=
\frac{\langle\varphi_s^{-1}\circ \sqrt{-1} F_{h_s}(E_s)\varphi_s(\alpha),\alpha\rangle_h}
{\langle\alpha,\alpha\rangle_h}
=
\frac{\langle \sqrt{-1} F_{h_s}(E_s)\varphi_s(\alpha),\varphi_s(\alpha)\rangle_{h_s}}
{\langle\varphi_s(\alpha),\varphi_s(\alpha)\rangle_{h_s}}
\end{aligned}
\end{equation}
以上より
\begin{equation*}
\begin{aligned}
\langle \Lambda_{\omega}\frac{\sqrt{-1}}{2\pi} F_h(E)\alpha,\alpha\rangle_h
&\underalign{\eqref{eq:lemma26-11}}{=}
\frac{\langle \Lambda_{\omega}\frac{\sqrt{-1}}{2\pi} F_{h_s}(E_s)\varphi_s(\alpha),\varphi_s(\alpha)\rangle_{h_s}}
{\langle\varphi_s(\alpha),\varphi_s(\alpha)\rangle_{h_s}}
\cdot\langle\alpha,\alpha\rangle_h
\\
&\underalign{\eqref{eq:lemma26-9}}{\geq}
\bigl(\min(c_1,c_2)-\epsilon\bigr)\cdot
\frac{\langle \varphi_s(\alpha),\varphi_s(\alpha)\rangle_{h_s}}
{\langle\varphi_s(\alpha),\varphi_s(\alpha)\rangle_{h_s}}
\cdot\langle\alpha,\alpha\rangle_h
=\bigl(\min(c_1,c_2)-\epsilon\bigr)\cdot\langle\alpha,\alpha\rangle_h
\end{aligned}
\end{equation*}
を得る. よって\eqref{eq:lemma26-7-2}が成り立つ. 
\eqref{eq:lemma26-8}も\eqref{eq:lemma26-10}と\eqref{eq:lemma26-11}からも従う.
\end{proof}

\end{rem}



 \subsection{$-K_X$がnefならば基本群はalmost Abelian }
 $-K_X$がsemipositiveの仮定は弱めることができる. つまり以下が成り立つ
\begin{thm}[{\cite{Paun98}}]
コンパクトK\"ahler多様体$X$について, $-K_X$がnefならば基本群はalmost Abelianである. 
\end{thm}
Theorem~\ref{thm-Campana-Special}によって, 以下を示せば良い. 
\begin{thm}[{\cite{Paun98}}]
\label{thm-Paun-thm}
$X$をコンパクトK\"ahler多様体とする. 
$-K_X$ nefならば, 
基本群はalmost nilpotentである. 
\end{thm}
\begin{proof}
次のGeneralized Margulis lemmaを認める. 
\begin{thm}[Generalized Margulis lemma {\cite[Theorem 1]{KW11}}]
\label{thm-Generalized-Margulis-lem}
任意の整数$m\geq 2$に対して定数$\delta(m)>0$と$C(m)>0$が存在し, 次を満たす: 

$(M,g)$を完備$m$次元Riemann多様体とし, $p\in M$とする. $B_g(p,1)$上で
\[
\Ric(g)>-(m-1)g
\]
が成り立つと仮定する. このとき自然な準同型
\[
\pi_1(B_g(p,\delta(m)),p)\longrightarrow \pi_1(B_g(p,1),p)
\]
の像は, indexが高々$C(m)$であるnilpotent subgroupを含む. 特にこの像はalmost nilpotentである.
\end{thm}

$-K_X$がnefであるから, Corollary~\ref{cor-Yau-thm-2}から任意の$\varepsilon>0$に対して
\(
\{\omega_\varepsilon\}=\{\omega_0\}\) かつ\(\Ric(\omega_\varepsilon)\geq-\varepsilon\omega_\varepsilon
\)
を満たすK\"ahler計量$\omega_\varepsilon$が存在する. 
$G\coloneqq \pi_1(X)$とおき, 有限個のgenerator
\(
G=\langle\gamma_1,\ldots,\gamma_k\rangle
\)
を固定する.
\begin{claim}[{\cite[Lemma 1.3]{DPS96}}]
\label{claim-uniformD}
$\widetilde{X} \to X$を普遍被覆とし, $G=\langle\gamma_1,\ldots,\gamma_k\rangle$が$\widetilde{X}$に作用しているとみなす.
ある定数$D>0$があって, 任意の$\varepsilon$についてある$p_\varepsilon\in\widetilde X$があって
\[
d_{\pi^{*}\omega_\varepsilon}(p_\varepsilon,\gamma_j(p_\varepsilon))\leq D,\qquad j=1,\ldots,k
\]
\end{claim}

\begin{proof}
次のclaimを後に示す.  
\begin{claim}[{\cite[Lemma 1.3]{DPS96}}]
\label{claim-DPS-lem1.3}
任意のコンパクト集合 $U\subset\widetilde X$と任意の$\delta>0$に対して, 閉部分集合$U_{\varepsilon,\delta}\subset U$が存在して次を満たすようにできる. 
\begin{enumerate}
\item $\Vol_{\widetilde{\omega_0}}(U\setminus U_{\varepsilon,\delta})<\delta$
\item 任意の$x,y\in U_{\varepsilon,\delta}$は$\widetilde{\omega_\varepsilon}$-lengthが高々$C\delta^{-1/2}$であるpathによって結ぶことができる
\end{enumerate}
ここで$C>0$は$\varepsilon$および$\delta$に依らないようにできる.
\end{claim}
この部分はテクニカルな上に重要かどうか不明なので後に回す. 
さて基本領域$E\subset\widetilde X$とし, $E \subset U$となるコンパクト集合を
\[
\operatorname{Int}(U)\cap\gamma_j(\operatorname{Int}(U))\neq\emptyset,\qquad j=1,\ldots,k
\]
となるように大きくとる. 特に$j=1,\ldots,k$について
\(
\Vol_{\widetilde{\omega_0}}(U\cap\gamma_j(U))>0
\)
なので, 
\[
2\delta<\min_{1\leq j\leq k}\Vol_{\widetilde{\omega_0}}(U\cap\gamma_j(U))
\]
となる$\delta>0$を取ることができる. 

そこでClaim~\ref{claim-DPS-lem1.3}の(1), (2)を満たすような閉部分集合$U_{\varepsilon,\delta}\subset U$をとる. 
このとき各$j$について
\(
U_{\varepsilon,\delta}\cap\gamma_j(U_{\varepsilon,\delta})\neq\emptyset
\)
である.実際,もしこの交わりが空ならば
\[
U\cap\gamma_j(U)\subset
(U\setminus U_{\varepsilon,\delta})
\cup
\gamma_j(U\setminus U_{\varepsilon,\delta})
\]
となる. $\gamma_j$は$\widetilde{\omega_0}=\pi^*\omega_0$に関するisometryなので
\(
\Vol_{\widetilde{\omega_0}}
\bigl(\gamma_j(U\setminus U_{\varepsilon,\delta})\bigr)
=
\Vol_{\widetilde{\omega_0}}
(U\setminus U_{\varepsilon,\delta})
<\delta.
\)であり, これから
\(
\Vol_{\widetilde{\omega_0}}(U\cap\gamma_j(U))<2\delta
\)
となり, $\delta$の取り方に矛盾する.
そこで各$j$に対して
\[
q_j\in U_{\varepsilon,\delta},\quad
\gamma_j(q_j)\in U_{\varepsilon,\delta}
\]
となる$q_j$を取る.

$p_\varepsilon\in U_{\varepsilon,\delta}$を固定する.すると任意の$j=1,\ldots, k$について
\begin{itemize}
\item $p_\varepsilon,\gamma_j(q_j)\in U_{\varepsilon,\delta}$より,  (2)からこの二つを結ぶ$\widetilde{\omega_\varepsilon}$-lengthが高々$C\delta^{-1/2}$であるpathが存在.
\item 同様に$q_j$と$p_\varepsilon$を結ぶ$\widetilde{\omega_\varepsilon}$-lengthが高々$C\delta^{-1/2}$であるpathが存在. よって$\gamma_j$を作用させて, $\gamma_j(q_j)$と$\gamma_j(p_\varepsilon)$を結ぶ$\widetilde{\omega_\varepsilon}$-lengthが高々$C\delta^{-1/2}$であるpathが存在. 
\end{itemize}
したがって
\[
d_{\widetilde{\omega_\varepsilon}}(p_\varepsilon,\gamma_j(p_\varepsilon))
\leq 2C\delta^{-1/2}\eqqcolon D
\]
 \end{proof}


さて$X$上のRiemann計量を
\(
\widehat\omega_\varepsilon\coloneqq \varepsilon \omega_\varepsilon
\)
とおく. $\Ric(g)=-\frac{\sqrt{-1}}{2\pi}\partial \bar\partial\log \det(g)$なので, 特に定数倍ではRicci tensorは変わらず, 
\[
\Ric(\widehat\omega_\varepsilon)=\Ric(\omega_\varepsilon)\geq-\varepsilon \omega_\varepsilon=-\widehat\omega_\varepsilon
\geq -(2n-1)\widehat\omega_\varepsilon
\]
一般にRiemann計量$g$に関するpath$\gamma$の距離は$L(\gamma)=\int_{I}\sqrt{g(\dot{\gamma}, \dot{\gamma})}dt$なので, 特に$g$を定数$\varepsilon$倍すると長さは$\sqrt{\varepsilon}$倍される. 
よって
\(
d_{\pi^{*}\widehat\omega_\varepsilon}=\sqrt{\varepsilon}\,d_{\pi^{*}\omega_\varepsilon}\)であり特に
\begin{equation}
\label{eq-D-varepsilon}
d_{\pi^{*}\widehat\omega_\varepsilon}(p_\varepsilon,\gamma_j(p_\varepsilon))
=\sqrt{\varepsilon}\,d_{\pi^{*}\omega_\varepsilon}(p_\varepsilon,\gamma_j(p_\varepsilon))
\leq\sqrt{\varepsilon}D.
\end{equation}

今Theorem~\ref{thm-Generalized-Margulis-lem}での定数を$\delta(m)$とし, 
$\varepsilon>0$を十分小さく取って
\(
\sqrt{\varepsilon}D<\delta(m)
\)
とする. 
\eqref{eq-D-varepsilon}から
\[
\gamma_j \in B_{\widehat\omega_\varepsilon}(\pi(p_\varepsilon),\delta(m))
\]
である. ここでこの$\gamma_j $は始点$\pi(p_\varepsilon)$で終点$\pi(\gamma_j p_\varepsilon)$となるloopとみている. 
以上より,  $\gamma_1,\ldots,\gamma_k$は$\pi_1(X)$を生成するので,
\[
\pi_1(B_{\widehat\omega_\varepsilon}(\pi(p_\varepsilon),\delta(m)),\pi(p_\varepsilon))\longrightarrow\pi_1(X,\pi(p_\varepsilon))
\]
は全射である.よってGeneralized Margulis lemma (Theorem~\ref{thm-Generalized-Margulis-lem})から, $\pi_1(X,\pi(p_\varepsilon))$はalmost nilpotentである.
\end{proof}

残っていたClaim~\ref{claim-DPS-lem1.3}の証明をする.
\begin{proof}[Proof of Claim~\ref{claim-DPS-lem1.3}]
\(
\widetilde{\omega_\varepsilon}\coloneqq \pi^*\omega_\varepsilon
\),\(\widetilde{\omega_0}\coloneqq \pi^*\omega_0
\)
とおく.$\{\omega_\varepsilon\}=\{\omega_0\}$であるから
\[
\int_X\operatorname{tr}_{\omega_0}(\omega_\varepsilon)\,\omega_0^n
=n\int_X\omega_\varepsilon\wedge\omega_0^{n-1}
=n\int_X\omega_0^n.
\]
したがって$\lVert\omega_\varepsilon\rVert_{L^1(X,\omega_0)}$は$\varepsilon$に関して一様に有界である.

$U$がコンパクトなので, $U=K$で$K$が$\widetilde X$のあるcoordinate open set
$\Omega$内のcompact convex subsetである場合に帰着できる. この場合に示そう. 

$x_1,x_2\in K$に対して
\[
c_{x_1,x_2}(t)\coloneqq (1-t)x_1+tx_2,\qquad 0\leq t\leq1
\]
とおく.$K$はconvexなので$c_{x_1,x_2}([0,1])\subset K$である.
このsegmentの$\widetilde{\omega_\varepsilon}$-lengthを
\[
L_\varepsilon(x_1,x_2)
\coloneqq 
\int_0^1
\sqrt{
\widetilde{\omega_\varepsilon}
\bigl((1-t)x_1+tx_2\bigr)
(x_2-x_1,x_2-x_1)
}\,dt
\]
とする.Cauchy--Schwarzの不等式より
\[
L_\varepsilon(x_1,x_2)^2
\leq
\int_0^1
\widetilde{\omega_\varepsilon}
\bigl((1-t)x_1+tx_2\bigr)
(x_2-x_1,x_2-x_1)\,dt.
\]
固定された座標近傍上ではEuclid計量と
$\widetilde{\omega_0}$は一様に同値であり,
$\lvert x_2-x_1\rvert\leq\diam K$であるから,ある定数$C_K>0$について
\[
L_\varepsilon(x_1,x_2)^2
\leq
C_K\int_0^1
\left\lvert 
\widetilde{\omega_\varepsilon}
\bigl((1-t)x_1+tx_2\bigr)
\right\rvert_{\widetilde{\omega_0}}dt.
\]
したがってFubiniの定理より
\[
\int_{K\times K}L_\varepsilon(x_1,x_2)^2\,dx_1dx_2
\leq
C_K\int_0^1dt\int_{K\times K}
\left\lvert 
\widetilde{\omega_\varepsilon}
\bigl((1-t)x_1+tx_2\bigr)
\right\rvert_{\widetilde{\omega_0}}
dx_1dx_2
\leq
C\lVert\widetilde{\omega_\varepsilon}\rVert_{L^1(\Omega,\widetilde{\omega_0})}.
\]
最後の計算は変数変換による. 
$\Omega\Subset\widetilde X$は固定されており,
$\widetilde{\omega_\varepsilon}=\pi^*\omega_\varepsilon$であるから, $\lVert\omega_\varepsilon\rVert_{L^1(X,\omega_0)}$は$\varepsilon$に関して一様に有界であることを用いれば, この式の右辺は$\varepsilon$に依らず取ることができる.
そこで
\[
S\coloneqq 
\left\{
(x_1,x_2)\in K\times K\;\middle\vert\;
L_\varepsilon(x_1,x_2)>
\left(\frac{C_1}{\delta}\right)^{1/2}
\right\}
\]
とおく.すると
\[
C_1
\geq
\int_S L_\varepsilon(x_1,x_2)^2\,dx_1dx_2
>
\frac{C_1}{\delta}\Vol(S)
\Rightarrow \Vol(S)<\delta.
\]

各$x\in K$について
\[
S(x)\coloneqq \{y\in K\mid (x,y)\in S\}
\quad \text{and} \quad
Q\coloneqq 
\left\{
x\in K\;\middle\vert\;
\Vol(S(x))
\geq\frac12\Vol(K)
\right\}
\]
とおく. Fubiniの定理より
\[
\delta>\Vol(S)
=\int_K\Vol(S(x))\,dx
\geq
\frac12\Vol(K)\Vol(Q),
\Rightarrow 
\Vol(Q)
<
\frac{2\delta}{\Vol(K)}.
\]

ここで$x_1,x_2\in K\setminus Q$とする.このとき$Q$の定義から
\[
\Vol(S(x_1)), \Vol(S(x_2))
<
\frac12\Vol(K)
\Rightarrow
\Vol\bigl(S(x_1)\cup S(x_2)\bigr)
<
\Vol(K)
\Rightarrow
K\setminus\bigl(S(x_1)\cup S(x_2)\bigr)\neq\emptyset.
\]
そこで
\(
y\in K\setminus\bigl(S(x_1)\cup S(x_2)\bigr)
\)
を取る.すると
\(
(x_1,y)\notin S, (x_2,y)\notin S
\)
であるから
\[
L_{\widetilde{\omega_\varepsilon}}([x_1,y])
\leq
\left(\frac{C_1}{\delta}\right)^{1/2},
L_{\widetilde{\omega_\varepsilon}}([x_2,y])
\leq
\left(\frac{C_1}{\delta}\right)^{1/2}
\Rightarrow
d_{\widetilde{\omega_\varepsilon}}(x_1,x_2)
\leq
2\left(\frac{C_1}{\delta}\right)^{1/2}.
\]

また固定されたコンパクト集合 $K$上ではEuclid計量による体積と
$\widetilde{\omega_0}$による体積は一様に同値であるから,
ある定数$C_2>0$について
\[
\Vol_{\widetilde{\omega_0}}(Q)
\leq C_2\Vol(Q)
<
\frac{2C_2}{\Vol(K)}\delta.
\]
したがって最初の$\delta$を適当な定数倍に置き換えることにより
\(
\Vol_{\widetilde{\omega_0}}(Q)<\delta
\)
としてよい.よって
\(
K_{\varepsilon,\delta}\coloneqq K\setminus Q
\)
とおけば
\[
\Vol_{\widetilde{\omega_0}}
(K\setminus K_{\varepsilon,\delta})<\delta
\]
かつ任意の$x_1,x_2\in K_{\varepsilon,\delta}$を
$\widetilde{\omega_\varepsilon}$-lengthが$C\delta^{-1/2}$以下であるpathによって結ぶことができる.
\end{proof}


 \subsection{$-K_X$がnefな多様体の基本群の歴史}

Fanoが単連結の歴史はかなり古いらしい. 
藤野先生のサーベイ「ファノ多様体についてのいくつかの問題」\url{https://www.math.kyoto-u.ac.jp/~fujino/Fano-problems5.pdf}
によると
\begin{tcolorbox}[mybox]
厳密な証明がついたのは, Atiyahの論文(1976)
が最初のようである. $\cdots$ (中略)$\cdots$
さらに調べてみると, この問題への最
初の貢献は小林昭七先生の論文(1961)のようである.
\end{tcolorbox}

Fanoが単連結の証明はかなり面白い. 知ってる限り3つの証明がある. 
\begin{enumerate}
\item $\Ric>0$な計量を作ってMyersの定理を使う
\item $\pi_1(X)$に関するShafarevich mapが自明であることを示す. 
\item MRC fibtrationが自明であることを示して, $X$が有理連結であることを示す. (有理連結ならば単連結である)
\end{enumerate}

(2)の方向性は以下の通り
\begin{itemize}
\item KLT Fano (もっと弱くweak Fano)も単連結. \cite{Tak00}による. この証明はかなり多変数複素解析である.
\item KLT Fano (もっと弱くweak Fano)のregular locus $X_\reg$の基本群は有限\cite{Bra21}による. この証明(の半分くらいは)高山先生の証明をorbifoldに拡張することで証明する. かなり面白い証明である. ちなみにweak Fanoの$X_{\reg}$の代数的基本群の有限性は\cite{Xu14}で知られていた. 
\end{itemize}
(3)の方向性は以下の通り
\begin{itemize}
\item Fano多様体は有理連結は\cite{Cam92, KMM92}による. 
\item KLT weak Fanoは有理連結. \cite{Zhang04}による. この方向性, 他にも\cite{BP11, EG19, CCM21, EIM21}などいろいろあるが, 実は多変数複素解析を用いるとかなりわかりやすく理解できる. それが\cite{Deng21}である. \cite{Deng21}の主定理は代数幾何学の話題だが, 証明はOhsawa-Takegoshiを用いる解析的で面白い証明である. 
\item KLT Fanoの$X_\reg$の基本群の有限性に関して, \cite{JLR25}で有理連結っぽい感じで証明された. 1-free curveの存在がわかるようである. \cite{Bra21}のイントロで書かれていた内容が実現したようである. 
\end{itemize}

$-K_X$がnefの場合の基本群に関しては以下の通り 
\begin{itemize}
\item $X$がsmoothな場合は\cite{Paun98}が初め. \cite{DPS96}をかなり使う. 
\item $X$がKLTの場合は$X$も基本群がalmost Abelianであるかはわかっていない. 3次元の場合に$\pi_{1}(X)$のalmost Abelianが\cite{FGSW26}で解決. 
\item \cite{GGK19, Cam21}より$X$がKLTで$c_1(K_X)=0$の場合,  irreducible symplectic varietiesか偶数次元の強い意味でのCalabi-Yauならば, $\pi_1(X)$が有限がわかっている. 
これは$\chi(X, \mathcal{O}_X) \neq 0$であり, \cite{Cam95}を用いる. よって残りは奇数次元の強い意味でのCalabi-Yauである. おそらくこれの$\pi_1(X)$が有限がわかれば, $-K_X$がnefで$X$がKLTの場合の基本群$\pi_1(X)$がalmost Abelianは\cite{MW21}の構造定理からわかるはずである. 
\end{itemize}

$-K_X$がnefでKLTの場合の$\pi_1(X_\reg)$のalmost Abelianは全くわかっていない. 3次元でも未解決かもしれない. また$c_1(K_X)=0$の場合に帰着できるかも不明らしい. これは有理連結KLT多様体の$X_\reg$がどうなるかが難しいためである. なお$-K_X$がnefで$X$がKLTの場合は, $\pi_1(X_\reg)$の有限次元表現がalmost Abelianな像を持つことは\cite{GGK19}と\cite{MW21}を用いればわかる.(\cite{IMM24}参照)  

\newpage


\bibliographystyle{alpha}
\bibliography{ref_Yau_merged.bib}
\end{document}



















